% !TEX root = main.tex
\providecommand{\reg}[1]{\ensuremath{\mathsf{#1}}}

% ===================================================================
\section{Notation and Overview}
\label{app:map}
\label{app:formal}

\subsection{Notation and Conventions}

\Cref{tab:notation-all} collects the notation used across the main text and appendices by topic, followed by the formal definitions of the Markov counting task, two-block Transformer architecture, balanced small initialization, and projected population gradient flow used in \cref{app:architectures,app:damping-proof,app:joint,app:adaptive}.

\begin{table}[!htbp]
\caption{\textbf{Summary of core notation across the main text and appendices, grouped by theme.}}
\label{tab:notation-all}
\centering\footnotesize
\renewcommand{\arraystretch}{1.22}
\fitwidth{\setlength{\tabcolsep}{5.5pt}\begin{tabular}{@{}lll@{}}
\toprule
\textbf{Notation} & \textbf{Meaning} & \textbf{Definition / Reference}\\
\midrule
\multicolumn{3}{@{}l}{\emph{Markov Counting Task and Parity Contraction}}\\
$n,\; \Sigma,\; w^\star,\; w_1,\ldots,w_n$ & Horizon, alphabet, query letter, context letters & \Cref{sec:data}\\
$p,\; b_t,\; x_{t-1},\; \chi_t$ & Match rate $\Pr(w_t = w^\star)$, bit $\one[w_t = w^\star]$, state sign $(-1)^{s_{t-1}}$, flip sign $(-1)^{b_t}$ & \Cref{eq:tracking}\\
$s_t \in \{0,1\},\; s_n \in \{0,1\}$ & Running modulo-$2$ count and terminal state & \Cref{eq:tracking}\\
$e(s),\; y_t := e(s_t)$ & One-hot state map $e(s) := e_s \in \{e_0, e_1\}$ and label vector ($0 \le t \le n$) & \Cref{sec:notation}\\
$\mixrate = |1-2p| \in [0,1)$ & Single-step Markov parity contraction factor & \Cref{eq:lambda}\\
$q(s),\; d_t = y_t - q(s_{t-1})$ & One-step state prior $(1-p)e_s + p e_{1-s}$ and $2$D step rule innovation & \Cref{eq:local-variation}\\
$e_0, e_1,\; e_\pm,\; \bary,\; R$ & Basis of $\R^2$, Hadamard vectors $\frac{1}{\sqrt 2}(e_0 \pm e_1)$, uniform $\frac{1}{2}\one$, swap matrix & \Cref{sec:notation}\\
\midrule
\multicolumn{3}{@{}l}{\emph{Two-Block Causal Transformer and Supervision Objectives}}\\
$\dmodel,\; d_k,\; \dmlp$ & Residual stream width, query--key head dimension, ReLU MLP width & \Cref{ass:arch,prop:represent}\\
$W_Q^{(\ell,h)}, W_K^{(\ell,h)}, W_V^{(\ell,h)}$ & Query, key ($d_k \times \dmodel$), and value ($\dmodel \times \dmodel$) blocks & \Cref{eq:block-Equations}\\
$W_{\mathrm{up}}^{(\ell)}, b_{\mathrm{up}}^{(\ell)}, W_{\mathrm{down}}^{(\ell)},\; V$ & ReLU MLP weights/biases and final linear readout ($2 \times \dmodel$) & \Cref{ass:arch}\\
$\hat y_t(\theta),\; p_\theta(s_t),\; \Lout,\; \Lstep$ & Softmax distribution, state probability $\ip{\hat y_t(\theta)}{y_t}$, outcome and step losses & \Cref{eq:losses}\\
$\Jz,\; \outcredit(\theta),\; \rulecredit(\theta)$ & Logit Jacobian, outcome credit $\E[\langle z_\theta, y_n - \bary\rangle]$, step rule credit & \Cref{sec:credit-results}\\
\midrule
\multicolumn{3}{@{}l}{\emph{Structured Step Sub-Circuit and Full-Prefix Routing}}\\
$v,\; r_{\mathrm{align}},\; \theta_{\mathrm s}$ & Frozen letter-value scale, initial QK alignment, and Block-$2$ step parameters & \Cref{def:gate}\\
$(\cnt, g, a, \brel),\; c^\star,\; \Delta_0,\; \kappa(a)$ & Reduced coordinates, target prior logit $\log\frac{1-p}{p}$, Gram invariant, score mobility & \Cref{prop:reduced-transition}\\
$Z_t, N_t,\; \pi_{t,r}, \pi_t^\star,\; u_t, D_t, \mu_t$ & Letter/non-letter partition sums, letter weights, log-ratio, gate, rule signal & \Cref{prop:reduced-transition,lem:J-reduce}\\
$\mathcal W_t, \overline{\mathcal W}_t,\; \omega_\delta,\; c_{\mathrm F}$ & Letter and non-letter lags at step $t$, reciprocal lag weight $e^{-\brel_\delta}$, distractor-floor constant & \Cref{prop:reduced-transition,lem:U-lag-cancel,app:hypF}\\
\midrule
\multicolumn{3}{@{}l}{\emph{Neighborhoods, Norms, Timescales, and Asymptotic Conventions}}\\
$\|\cdot\|_{\op},\; \|\cdot\|_{\frob},\; \|\cdot\|_*$ & Operator (largest singular value), Frobenius, and nuclear (trace) norms & \Cref{rem:norms,prop:reduced-transition}\\
$\initscale,\; \ballradius,\; \initball,\; \mathcal D_{\initscale}$ & Initialization scale, certified radius, $\brel{=}0$ and unrestricted-$\brel$ balls & \Cref{ass:init,lem:bias-uniform}\\
$\tilde\tau := \initscale^2\tau,\; s := \initscale\tau$ & Dimensionless Euclidean prior time and Block-RMS time & \Cref{app:J-timescales,lem:rms-align}\\
$O, \Theta, \Omega, \lesssim, \asymp, \widetilde O$ & Constants depend only on $p$; $\widetilde O$ also hides logarithmic factors in $n, 1/\initscale, 1/r_{\mathrm{align}}$ & \Cref{sec:trajectory-results}\\
\bottomrule
\end{tabular}}
\end{table}

Writing $\ell_t(\theta) := \E[-\log p_\theta(s_t)]$ for the cross-entropy loss at step $t$ (where $p_\theta(s_t) = \ip{\cS(z_t(\theta))}{y_t}$ and $y_t = e(s_t)$), the step-supervised objectives over $t \ge 2$ and over all $t \ge 1$ are
\begin{equation}
 \Lstep(\theta) := \frac{1}{n-1}\sum_{t=2}^n \ell_t(\theta),
 \qquad
 \Lstep^{\mathrm{all}}(\theta) := \frac{1}{n}\sum_{t=1}^n \ell_t(\theta).
 \label{eq:step-objectives}
\end{equation}
At each step $t \in [n]$, the one-step transition prior from state $s \in \{0,1\}$ is $q(s) := (1-p)e_s + p e_{1-s}$, and the local step innovation is $d_t := y_t - q(s_{t-1})$. The outcome and step rule credits are defined by\label{tab:credit-summary}
\[
  \outcredit(\theta) := \E\bigl[\langle z_\theta, y_n - \bary\rangle\bigr],
  \qquad\qquad
  \rulecredit(\theta) := \frac{1}{n-1}\sum_{t=2}^n \E\bigl[\langle z_t(\theta), d_t\rangle\bigr].
\]

\begin{definition}[Modulo-$2$ Markov counting]
\label{def:task}
Fix an alphabet $\Sigma$, horizon $n \ge 2$, and match probability $p \in (0,1)$. A prompt $(s_0, w^\star, w_1, \ldots, w_n) \sim \mathcal D$ draws $s_0 \sim \mathrm{Unif}(\{0,1\})$, query $w^\star \in \Sigma$, and context letters $w_1, \ldots, w_n \in \Sigma$ conditionally independently given $(s_0, w^\star)$ with $\Pr(w_t = w^\star \mid s_0, w^\star) = p$. With $b_t := \one[w_t = w^\star]$, $x_{t-1} := (-1)^{s_{t-1}}$, $\chi_t := (-1)^{b_t}$, and $\mixrate := |\E[\chi_t]| = |1-2p| \in [0,1)$, states evolve as $s_t = s_{t-1} + b_t \pmod 2$ and $(-1)^{s_t} = x_{t-1}\chi_t$ with terminal label $y_n = e(s_n)$.
\end{definition}

\begin{definition}[Two-block causal Transformer]
\label{ass:arch}
From token-plus-position embeddings $h_i^{(0)} \in \R^{\dmodel}$, each block $\ell \in \{1,2\}$ applies three-head causal softmax attention (head width $d_k$, relative lag biases $\brel_{i-j}^{(\ell,h)}$ with $\brel_0 = 0$) over the full causal prefix $j \in \{0,\ldots,i\}$ without a target-lag mask, followed by a width-$\dmlp$ ReLU MLP without layer normalization and linear readout $\TF_\theta(x)_i = V h_i^{(2)} \in \R^2$ (\cref{app:architectures}).
\end{definition}

\begin{assumption}[Balanced small initialization]
\label{ass:init}
Non-bias weights and embeddings initialize at $\theta_0 = \initscale\tilde\theta$ ($0 < \initscale \le \initscale_0$) with $\brel(0) = 0$, balanced initial readout ($\cS(z_{\theta_0}) \equiv \bary$, e.g.\ $V(0) = 0$), $O(1)$ unscaled operator norms, non-degenerate state embeddings, and $|\tilde b_m| \ge 1$ so ReLU masks stay locally invariant near $\theta_0$. We write $\initball := \{\theta : \|\theta - \theta_0\|_2 \le \ballradius\initscale,\, \brel = 0\}$.
\end{assumption}

\begin{definition}[Projected population gradient flow]
\label{ass:gd}
Parameters evolve under projected population gradient flow
\[
 \ddtau{\theta} = -P\nabla_\theta\mathcal L(\theta)
\]
for $\mathcal L \in \{\Lout, \Lstep\}$ over continuous training time $\tau \ge 0$, where $P$ is the orthogonal projection onto a subset of trainable parameter blocks. \Cref{thm:main} bounds outcome credit uniformly across $\mathcal D_{\initscale}$ (including arbitrary $\brel$), so \cref{thm:outcome-end} allows $\brel$ to evolve under any piecewise-constant block projection schedule $P(\tau)$ via the general energy barrier of \cref{lem:energy-barrier}.
\end{definition}

\begin{remark}[Parameter norms and width-independent constants]
\label{rem:norms}
\label{rem:end-width}
We equip the parameter space with the Euclidean norm $\|\theta\|_2 := (\sum_{\mathrm B} \|\theta_{\mathrm B}\|_{\frob}^2)^{1/2}$ summed over all parameter blocks $\mathrm B$. In \cref{ass:init}, each unscaled weight and embedding block of $\tilde\theta$ has operator norm at most $M_0$, the count embeddings satisfy $\sigma_{\min}(\tilde E_{\mathrm{val}}) \ge \embconst > 0$, and the unscaled MLP biases satisfy $|\tilde b_m| \ge 1$ with $\|\tilde b\|_2 \le N_b$. We set the neighborhood radius and scale threshold to
\[
 \ballradius := \min\Bigl(\frac{1}{8}, \frac{\embconst}{4}\Bigr), \qquad \initscale_0 := c\min\bigl(1, N_b^{-1}\bigr),
\]
where the constants $c, C > 0$ depend only on $M_0$, $\embconst$, and the head and block counts.

On $\initball$, the margin $|\initscale\tilde b_m| \ge \initscale \gg O(\initscale^2)$ prevents ReLU pre-activations from changing sign across all inputs, ensuring the higher-order smoothness used in the attention-score expansion of \cref{prop:network} (\cref{sec:experiments} confirms empirically that the outcome plateau persists under standard PyTorch initialization without bias rescaling). Standard fan-in-scaled random initialization with $|\tilde b_m| \in [1,2]$ satisfies $N_b \le 2\sqrt{\dmlp}$ and $\embconst \ge 1/2$ with high probability at every width \citep{vershynin2018high}. At the linear executor widths $\dmodel, d_k, \dmlp = \Theta(n)$ of \cref{prop:represent}, all constants remain independent of $n$ whenever $\initscale \le c\min(\mixrate, n^{-1/2})$.
\end{remark}

\subsection{Comparison with Prior Work}

\Cref{tab:related} compares our setting, architecture, and dynamical guarantees with prior theoretical analyses of chain-of-thought, sequential state tracking, and parity learning.

\begin{table}[!htbp]
\caption{\textbf{Detailed comparison with prior theoretical analyses of chain-of-thought, sequential state tracking, and parity learning.}}
\label{tab:related}
\centering\footnotesize
\fitwidth{\setlength{\tabcolsep}{5.5pt}\begin{tabular}{@{}llllll@{}}
\toprule
\textbf{Dimension} & \citet{kim2024parity} & \citeauthor{wen2024sparse}; \citeauthor{huang2025cotgd} & \citeauthor{pmlr-v291-joshi25a}; \citeauthor{hanneke2026samplecomplexityautoregressivereasoning} & \citeauthor{abbe2025parities}; \citeauthor{pmlr-v336-giapitzakis26a} & \textbf{This paper}\\
\midrule
Target task & Hidden parity family & Sparse parity / linear CoT & Autoregressive CoT class & Sparse parity / automaton family & \textbf{Single fixed rule} ($\mathrm{SQ\text{-}dim}=1$), invertible FSM\\
Architecture & $1$-layer position-only & $1$--$2$-layer staged / masked & Abstract hypothesis class & $2$-layer MLP / SQ oracle & \textbf{$2$-block softmax $QK^\top$ + ReLU Transformer}\\
Outcome lower bound & Finite-precision GD over class & Sample / masked gap & ERM sample lower bound & Class SQ / noisy-GD bound & \textbf{Single-rule $\Omega(\mixrate^{-n})$ stay time (full model, CE)}\\
Step optimization & $1$-layer position-only, squared loss & Layerwise staged freezing & Statistical ERM (no GD) & --- & \textbf{Simultaneous causal-prefix $(W_Q, W_K, \brel, V)$ flow}\\
Saddle dynamics & Bypassed (no $QK^\top$) & Bypassed by staging & None (no dynamics) & --- & \textbf{Baseline-denoised $|a|^{3/2}$ escape, $O(\initscale^{-5})$}\\
Shared expressivity & Separate models & CoT-only model & Separate sample bounds & --- & \textbf{Both executors in one model; greedy rollout}\\
\bottomrule
\end{tabular}}
\end{table}

\subsection{Frequently Asked Questions}
\label{app:faq}

\begin{enumerate}[label=\textbf{Q\arabic*.},leftmargin=*,itemsep=3pt,topsep=3pt]
\item \textbf{Why do \cref{thm:joint,thm:block-rms} analyze steps $t \in \{2,\ldots,n\}$ rather than $t = 1$?}
For every $t \ge 2$, the query token at index $n+1+t$ is $s_{t-1}$ itself, and its matching letter $w_t$ sits at constant lookback lag $\delta^\star = n$. At $t = 1$, the query position $n+2$ holds the delimiter $\mathtt{SEP}$ while $s_0$ sits back at index $0$, requiring an extra position-$0$ hop (\cref{rem:first-state}). All end-to-end experiments supervise $t = 1,\ldots,n$ and evaluate rollout from $s_0$.

\item \textbf{Why does \cref{thm:outcome-end} analyze the full two-block network while \cref{thm:joint,thm:block-rms} use the Block-$2$ circuit (\cref{def:gate})?}
Under outcome supervision, we bound credit across all $27$ parameter blocks simultaneously (\cref{thm:outcome-end}). Under step supervision, unlike prior CoT analyses that rely on position-only attention, pre-masked prefixes, or layerwise curriculum freezing \citep{kim2024parity,wen2024sparse,huang2025cotgd}, \cref{def:gate} evolves $(W_Q, W_K, W_V, \brel, V)$ \emph{simultaneously} over the \emph{unmasked} causal prefix and reduces \emph{without approximation} to a four-variable ODE (\cref{prop:reduced-transition}), while our external-offset ablation (\cref{tab:clamp-main,fig:feature-learning}) isolates readout-to-Block-$1$ coupling in full two-block networks.

\item \textbf{What is Hypothesis~(F), and why does \cref{thm:block-rms} not need it?}
Under plain Euclidean flow, query--key scores stay near zero for $\Theta(\initscale^{-5})$ time while lag biases drift; Hypothesis~(F) prevents non-letter lag weights at late lags $\delta \ge n+2$ from collapsing during that long plateau (proved unconditionally for $\delta \le n+1$ and checked numerically to $n = 2{,}000$; \cref{lem:U-floor-vis,app:hypF}). Under stateless Block-RMS (\cref{thm:block-rms}), unit-RMS updates open the query--key gate before distractor weights decay, removing Hypothesis~(F) altogether.

\item \textbf{Could AdamW's adaptive scaling explain the gap instead of step supervision?}
No. From identical initializations and minibatch streams (\cref{tab:transfer-main}a), stateless Block-RMS solves $5/5$ step runs in $1.20$k median steps while matched outcome Block-RMS solves $0/5$ through $10$k updates ($1.001$ bits loss, $49.1\%$ accuracy): at batch size $B = 256$, sampling noise ($B^{-1/2} \approx 0.063$) swamps the contracted parity signal ($\mixrate^8 \approx 0.004$).

\item \textbf{Why do natural step runs solve $0/20$ seeds within $2.4$k updates at extreme skew ($p \le 0.10$) in \cref{tab:causal-shortcut}?}
Baseline calibration cancels wrong-lag distractors at every $p$, but it also scales the true-lag gradient by $2p(1-p)$. When $p \le 0.10$, that smaller signal stretches post-routing gate escape past the fixed $2.4$k intervention budget (\cref{sec:causal-shortcut}).

\item \textbf{Does omitting LayerNorm in \cref{ass:arch} drive the separation?}
No. Omitting normalization keeps the Taylor and ODE analyses exact, while \cref{tab:transfer-main}b verifies the same step-versus-outcome separation in Pre-LN GPT-2-2L and Pre-RMSNorm SmolLM2-135M and Qwen2.5-0.5B models.
\end{enumerate}

% ===================================================================
\section{Proof of \texorpdfstring{\Cref{prop:represent}}{Proposition 1}}
\label{app:architectures}
\label{app:construction}
\label{app:model}

\paragraph{Proof roadmap for \cref{prop:represent}.}
We construct explicit parameter realizations $\thstep$ and $\thout$ within the single normalization-free two-block Transformer of \cref{ass:arch} at linear widths $d_k = 2n+6$, $\dmodel = 2|\Sigma| + d_k + 11$, and $\dmlp = \max\{|\Sigma|, 6(n+1)\}$. The proof proceeds in three modular steps: \cref{app:subspace-construction} specifies the ten-subspace residual decomposition and explicit weight blocks shared by or specialized to $\thstep$ and $\thout$ (\cref{tab:executors}); \cref{lem:exact-hard} verifies exact step and outcome execution in the hard-attention limit ($B \to \infty$); and \cref{lem:softmax-leakage} bounds off-target softmax leakage across both blocks at finite inverse temperature $B = \max\{16, 3\log(8(n+3))\}$, certifying a uniform logit margin strictly above $19$ on an open parameter neighborhood.

\subsection{Explicit Weight Construction}
\label{app:subspace-construction}

Each block $\ell \in \{1,2\}$ contains three causal attention heads of query--key dimension $d_k = 2n + 6$ and a ReLU MLP of hidden width $\dmlp = \max\{|\Sigma|, 6(n+1)\}$. Given residual vectors $h_i^{(\ell-1)} \in \R^{\dmodel}$, block $\ell$ computes
\begin{align}
 a_i^{(\ell,h)} &= \sum_{j=0}^i \cS_{j\le i}\!\Biggl(\frac{1}{\sqrt{d_k}}\bigl(W^{(\ell,h)}_Q h_i^{(\ell-1)}\bigr)^\top W^{(\ell,h)}_K h_j^{(\ell-1)}\Biggr)W^{(\ell,h)}_V h_j^{(\ell-1)}, &
 \tilde h_i^{(\ell)} &= h_i^{(\ell-1)} + \sum_{h=1}^3 a_i^{(\ell,h)}, \nonumber\\
 m_i^{(\ell)} &= W_{\mathrm{down}}^{(\ell)}\relu\bigl(W_{\mathrm{up}}^{(\ell)}\tilde h_i^{(\ell)} + b_{\mathrm{up}}^{(\ell)}\bigr), &
 h_i^{(\ell)} &= \tilde h_i^{(\ell)} + m_i^{(\ell)},
 \label{eq:block-Equations}
\end{align}
and the linear readout produces logits $\TF_\theta(x)_i = V h_i^{(2)}$. We decompose $\R^{\dmodel}$ ($\dmodel = 2|\Sigma| + d_k + 11$) into ten orthogonal coordinate blocks, writing the residual vector at position $i$ directly in block coordinates as
\[
 h_i
 = \bigl(\mathrm{tok}_i,\; \mathrm{val}_i,\; \mathrm{pos}_i,\; \mathrm{qry}_i,\; \hat b_i,\; \mathrm{prev}_i,\; \mathrm{cur}_i,\; \mathrm{sum}_i,\; \mathrm{init}_i,\; \mathrm{out}_i\bigr)
 \in \R^{|\Sigma|} \times \R^2 \times \R^{d_k} \times \R^{|\Sigma|} \times \R \times \R^2 \times \R \times \R \times \R^2 \times \R^2.
\]
Here $\mathrm{tok}_i$ and $\mathrm{val}_i$ store input letter and state embeddings, $\mathrm{pos}_i$ stores the one-hot position vector $e_i$, $\mathrm{qry}_i$ stores the copied query letter $w^\star$, and $\hat b_i$ stores the match bit $b_t$. The remaining slots hold the routed previous state and current match bit ($\mathrm{prev}_i, \mathrm{cur}_i$) for $\thstep$, the match sum and initial state ($\mathrm{sum}_i, \mathrm{init}_i$) for $\thout$, and the output state ($\mathrm{out}_i$). Prompt tokens sit at position $0$ ($s_0$), position $1$ ($w^\star$), positions $1+t$ ($w_t$, $1\le t\le n$), and position $n+2$ ($\mathrm{SEP}$); the step continuation places $s_t$ at $n+2+t$, the delimiter at $2n+3$, and $s_n$ at $2n+4$, whereas the outcome continuation places the delimiter at $n+3$ and $s_n$ at $n+4$.

\begin{table}[!htbp]
\caption{\textbf{Submodule roles in the two-block executors $\thstep$ and $\thout$.} Both executors share Block~$1$ and inhabit the same parameter space.}
\label{tab:executors}
\centering\footnotesize
\fitwidth{\begin{tabular}{@{}lll@{}}
\toprule
\textbf{Submodule} & \textbf{Step executor $\thstep$} & \textbf{Outcome executor $\thout$}\\
\midrule
Block 1, Head 1 & Copy query letter $w^\star$ from pos.\ $1$ into $\mathrm{qry}_i$ & Copy query letter $w^\star$ from pos.\ $1$ into $\mathrm{qry}_i$\\
Block 1, MLP & Write match bit $b_t = \one[w_t = w^\star]$ into $\hat b_{1+t}$ & Write match bit $b_t = \one[w_t = w^\star]$ into $\hat b_{1+t}$\\
Block 2, Head 1 & Copy previous state $s_{t-1}$ into $\mathrm{prev}_{i_t}$ & Sum $\hat b_{1+t}$ over $w_1,\ldots,w_n$ into $\mathrm{sum}_{n+3}$\\
Block 2, Head 2 & Copy current match bit $b_t$ into $\mathrm{cur}_{i_t}$ & Copy initial state $s_0$ from pos.\ $0$ into $\mathrm{init}_{n+3}$\\
Block 2, Head 3 & Copy final state $s_n$ into $\mathrm{out}_{2n+3}$ & Zero\\
Block 2, MLP & Modulo-$2$ sum $s_{t-1} + b_t \bmod 2 \to \mathrm{out}_{i_t}$ & Parity spline of $s_0 + \sum_{t=1}^n b_t \to \mathrm{out}_{n+3}$\\
Readout $V$ & Read $20\,\mathrm{out}_i$ at state and answer positions & Read $20\,\mathrm{out}_i$ at the answer position\\
\bottomrule
\end{tabular}}
\end{table}

To route uniformly over a prefix index set $\mathcal I(i) \subseteq \{0, \ldots, i\}$, fix $B := \max\{16, 3\log(8(n+3))\}$ and $C_Q := B\sqrt{d_k}$, and define $W_K, W_Q \in \R^{d_k \times \dmodel}$ on the $\mathrm{pos}$ block by $W_K h_j = \mathrm{pos}_j = e_j$ and $(W_Q h_i)_j = C_Q\,\one[j \in \mathcal I(i)]$, so
\begin{equation}
 \frac{1}{\sqrt{d_k}}(W_Q h_i)^\top W_K h_j = B\,\one[j \in \mathcal I(i)],
 \label{eq:route}
\end{equation}
paired with a coordinate-copy value matrix $W_V$ that copies the source slot into the destination slot and zeros all other coordinates. At input $h_i^{(0)}$, state tokens $s \in \{0,1\}$ set $\mathrm{val}_i = e(s)$, letter tokens $a \in \Sigma$ set $\mathrm{tok}_i = e_a$, delimiters set both to $0$, and position $i$ sets $\mathrm{pos}_i = e_i$.

\paragraph{Block~1 (shared query-matching circuit).}
Block~1 is identical in $\thstep$ and $\thout$ (with $\brel = 0$). Head~1 applies \cref{eq:route} with $\mathcal I(i) = \{1\}$ for $i \ge 1$ ($\mathcal I(0) = \{0\}$) to copy $\mathrm{tok}_1$ into $\mathrm{qry}_i$, while Heads~2 and~3 are zero. The Block~1 MLP uses $|\Sigma|$ active ReLU units indexed by $a \in \Sigma$ (padded with zeros to width $\dmlp$) to test $w_t = w^\star$ and write into the $\hat b$ slot:
\begin{equation}
 \hat b_{1+t} = \sum_{a\in\Sigma} 2\,\relu\Bigl((\mathrm{tok}_{1+t})_a + (\mathrm{qry}_{1+t})_a - \tfrac{3}{2}\Bigr).
 \label{eq:block1-mlp}
\end{equation}

\paragraph{Block~2 of the step executor $\thstep$.}
At each step-prediction position $i_t = n+1+t$ ($1 \le t \le n$):
\begin{itemize}[leftmargin=*]
\item \textbf{Head~1 (previous state):} Copies $\mathrm{val}$ into $\mathrm{prev}_{i_t}$ from $\mathcal I(i_1) = \{0\}$ at $t = 1$ or $\mathcal I(i_t) = \{i_t\}$ at $t \ge 2$ (holding $s_{t-1}$), and routes to $\{1\}$ elsewhere ($\mathcal I(0) = \{0\}$). Since position~$1$ holds $w^\star$ with $\mathrm{val}_1 = 0$, this fallback contributes only $O(e^{-B})$ softmax leakage at $i = 2n+3$.
\item \textbf{Head~2 (current match bit):} Copies $\hat b_{1+t}$ into $\mathrm{cur}_{i_t}$ via $\mathcal I(n+1+t) = \{1+t\}$ ($1 \le t \le n$, and $\mathcal I(i) = \{0\}$ otherwise).
\item \textbf{Head~3 (final-answer copy):} Copies $\mathrm{val}_{2n+2}$ into $\mathrm{out}_{2n+3}$ via $\mathcal I(2n+3) = \{2n+2\}$, and routes to $\{1\}$ for $1 \le i \ne 2n+3$ ($\mathcal I(0) = \{0\}$).
\end{itemize}
Four ReLU units indexed by $(s', b') \in \{0,1\}^2$ form pre-activations $q_{s',b'} := (\mathrm{prev}_i)_{s'} + (2b' - 1)\,\mathrm{cur}_i + (1 - b') - \tfrac{3}{2}$ and write modulo-$2$ addition into $\mathrm{out}_i$:
\begin{equation}
 \mathrm{out}_i = \sum_{s',b'\in\{0,1\}} 2\,e(s'+b'\bmod 2)\,\relu(q_{s',b'}),
 \label{eq:lookup}
\end{equation}
while staying inactive at $i = 2n+3$.

\paragraph{Block~2 of the outcome executor $\thout$ and shared readout.}
At the answer-prediction position $i = n+3$, Head~1 routes uniformly over $\mathcal I(n+3) = \{2,\ldots,n+1\}$ ($\mathcal I(i)=\{0\}$ for $i \ne n+3$) with value scale $n$ to write $\sum_{t=1}^n \hat b_{1+t}$ into $\mathrm{sum}_{n+3}$, Head~2 routes $\mathcal I(i) = \{0\}$ to copy $\mathrm{val}_0$ into $\mathrm{init}_{n+3}$, and Head~3 is zero. Using $6(n+1)$ ReLU units indexed by $(j, a, \varsigma) \in \{0,1\}\times\{0,\ldots,n\}\times\{+1,0,-1\}$ with $(c_{+1}, c_0, c_{-1}) = (1, -2, 1)$, the Block~2 MLP writes
\[
 \mathrm{out}_{n+3} = \sum_{j\in\{0,1\}}\sum_{a=0}^n e(j+a\bmod 2)\sum_{\varsigma\in\{+1,0,-1\}} c_\varsigma\,\relu\Bigl(\mathrm{sum}_{n+3} - a + \varsigma + (n+2)\bigl((\mathrm{init}_{n+3})_j - 1\bigr)\Bigr).
\]
Both executors share the linear readout $V h_i^{(2)} = 20\,\mathrm{out}_i \in \R^2$.

\subsection{Exactness and Margin Bounds}
\label{app:exact}
\label{tab:leakage-bounds}

We verify next that the explicit weights $\thstep$ and $\thout$ in \cref{app:subspace-construction} generate the exact step-by-step and direct-answer continuations on every prompt. \Cref{lem:exact-hard} checks exact execution with $\Omega(1)$ ReLU margins in the hard-attention limit $B \to \infty$, and \cref{lem:softmax-leakage} bounds the off-target softmax tail $\delta = (d_k - 1)e^{-B}$ across both blocks at finite inverse temperature $B$.

\begin{lemma}[Exact execution in the hard-attention limit]
\label{lem:exact-hard}
In the limit $B \to \infty$, where \cref{eq:route} assigns uniform attention weight $1/|\mathcal I(i)|$ to $\mathcal I(i)$ and zero outside $\mathcal I(i)$:
\begin{enumerate}[label=(\roman*),leftmargin=*]
  \item Shared Block~$1$ writes $\hat b_{1+t} = b_t = \one[w_t = w^\star]$ for every $t \in [n]$ with all ReLU pre-activations in $\{-3/2, -1/2, +1/2\}$.
  \item Under $\thstep$, Block~$2$ writes $\mathrm{out}_{n+1+t} = y_t$ ($1 \le t \le n$) and $\mathrm{out}_{2n+3} = y_n$ with all ReLU pre-activations bounded away from $0$ by at least $1/2$.
  \item Under $\thout$, Block~$2$ writes $\mathrm{out}_{n+3} = y_n$ via integer hat splines with all off-branch ReLU pre-activations $\le -1$.
\end{enumerate}
\end{lemma}
\begin{proof}[Proof of \cref{lem:exact-hard}]
For~(i), Head~1 places weight $1$ on position~$1$ and copies $\mathrm{tok}_1 = e_{w^\star}$ into $\mathrm{qry}_i$ for all $i \ge 1$. At context position $1+t$, $(\mathrm{tok}_{1+t})_a + (\mathrm{qry}_{1+t})_a$ equals $2$ when $a = w_t = w^\star$ and lies in $\{0,1\}$ otherwise, so the pre-activations in \cref{eq:block1-mlp} lie in $\{-3/2, -1/2, +1/2\}$ and write $\hat b_{1+t} = b_t$.

For~(ii), under $\thstep$ at $i = n+1+t$, Head~1 copies $y_{t-1} = e(s_{t-1})$ into $\mathrm{prev}_i$ and Head~2 copies $b_t$ into $\mathrm{cur}_i$. Among $(s', b') \in \{0,1\}^2$, $q_{s',b'}$ equals $+1/2$ at $(s_{t-1}, b_t)$ and is at most $-1/2$ at the other three pairs, so \cref{eq:lookup} writes $\mathrm{out}_i = y_t$. At $i = 2n+3$, Head~3 copies $\mathrm{val}_{2n+2} = y_n$ into $\mathrm{out}_{2n+3}$ while $\mathrm{prev}_{2n+3} = 0$ keeps all $q_{s',b'} \le -1/2$.

For~(iii), for any integers $m, a \in \{0, \ldots, n\}$, the three-ReLU hat spline satisfies
\[
  \relu(m - a + 1) - 2\relu(m - a) + \relu(m - a - 1) = \one[m = a].
\]
Under $\thout$ at $i = n+3$, Head~1 assigns weight $1/n$ to each $j \in \{2, \ldots, n+1\}$ and scales by $n$ to write $\mathrm{sum}_{n+3} = \sum_{t=1}^n b_t =: m \in \{0,\ldots,n\}$, while Head~2 copies $y_0 = e(s_0)$ into $\mathrm{init}_{n+3}$.

On the wrong initial-state branch $j = 1 - s_0$, $(\mathrm{init}_{n+3})_j - 1 = -1$ shifts every pre-activation to
\[
  m - a + \varsigma - (n+2) \le n - 0 + 1 - (n+2) = -1 < 0.
\]
On the true branch $j = s_0$, $(\mathrm{init}_{n+3})_{s_0} - 1 = 0$ evaluates the hat indicator at $a = m$:
\begin{equation}
 \mathrm{out}_{n+3} = \sum_{j\in\{0,1\}}\sum_{a=0}^n e(j+a\bmod 2)\sum_{\varsigma\in\{+1,0,-1\}} c_\varsigma\,\relu\bigl(\mathrm{sum}_{n+3} - a + \varsigma + (n+2)((\mathrm{init}_{n+3})_j - 1)\bigr) = y_n.\qedhere
 \label{eq:outcome-hat}
\end{equation}
\end{proof}

Next, setting $B = \max\{16, 3\log(8(n+3))\}$ keeps every ReLU unit on the same side of zero as in \cref{lem:exact-hard} and preserves a logit margin above $19$.

\begin{lemma}[Layerwise softmax leakage bound at finite $B$]
\label{lem:softmax-leakage}
At finite inverse temperature $B = \max\{16, 3\log(8(n+3))\}$, every routed head in \cref{eq:route} leaks total off-target softmax weight at most $\delta := (d_k - 1)e^{-B} < \frac{1}{256(n+3)^2}$, and at every supervised position $i$ of $\thstep$ and $\thout$ (with target state $s^\star(i) \in \{0,1\}$) the output slot satisfies $\|\mathrm{out}_i - e(s^\star(i))\|_\infty < \frac{1}{128}$, giving true-versus-false logit margin $> 19$.
\end{lemma}
\begin{proof}[Proof of \cref{lem:softmax-leakage}]
For any target set $\mathcal I(i) \subseteq \{0, \ldots, i\}$, off-target softmax weights sum to at most
\[
 \sum_{j \notin \mathcal I(i)} \alpha_{ij}
 \le \frac{i + 1 - |\mathcal I(i)|}{|\mathcal I(i)|e^{B}}
 \le (d_k - 1)e^{-B} = \delta < \frac{1}{256(n+3)^2}.
\]
We propagate $\delta$ across both blocks relative to the $B \to \infty$ states of \cref{lem:exact-hard}:
\begin{enumerate}[label=(\arabic*),leftmargin=*]
\item \emph{Block~$1$ query and match bit.} By convexity of the softmax average, $\|\mathrm{qry}_i - e_{w^\star}\|_\infty \le \delta < 2^{-10}$ for all $i \ge 1$. Because ideal Block~$1$ pre-activations lie in $\{-3/2, -1/2, +1/2\}$, no ReLU unit flips sign, so $|\hat b_{1+t} - b_t| \le 2\delta < 2^{-9}$.

\item \emph{Block~$2$ inputs.} Under $\thstep$ at $i = n+1+t$, Head~$1$ and Head~$2$ satisfy
\[
  \bigl\|\mathrm{prev}_i - y_{t-1}\bigr\|_\infty \le \delta
  \qquad\text{and}\qquad
  \bigl|\mathrm{cur}_i - b_t\bigr| \le 2\delta + (1+2\delta)\delta \le 4\delta < 2^{-7}.
\]
Under $\thout$ at $i = n+3$, Head~$2$ and Head~$1$ satisfy
\[
  \bigl\|\mathrm{init}_{n+3} - y_0\bigr\|_\infty \le \delta
  \qquad\text{and}\qquad
  \bigl|\mathrm{sum}_{n+3} - m\bigr| \le 4n\delta < \frac{1}{64(n+3)}.
\]

\item \emph{Block~$2$ outputs and readout margin.} Under $\thstep$ in \cref{eq:lookup}, each ideal pre-activation has margin $1/2$ from $0$ and shifts by at most $5\delta < 1/64$; adding Head~$3$ leakage $\le \delta$ gives $\|\mathrm{out}_i - e(s^\star(i))\|_\infty \le 11\delta < \frac{1}{128}$.

Under $\thout$ in \cref{eq:outcome-hat}, set $\tilde x_j := \mathrm{sum}_{n+3} + (n+2)((\mathrm{init}_{n+3})_j - 1)$. Wrong-branch units ($j = 1 - s_0$) stay inactive since their pre-activations are at most $-1 + (5n+2)\delta < 0$, while on the true branch $j = s_0$ we have $|\tilde x_{s_0} - m| \le (5n+2)\delta < \frac{1}{128}$. For $\tilde x_{s_0} \in (m - \frac{1}{2}, m + \frac{1}{2})$, triplets $a < m$ sum slopes to $1 - 2 + 1 = 0$, triplets $a > m$ are inactive, and only $(a=m, \varsigma=+1)$ varies as $\tilde x_{s_0} - m + 1$, giving $\|\mathrm{out}_{n+3} - y_n\|_\infty = |\tilde x_{s_0} - m| < \frac{1}{128}$.
\end{enumerate}
Reading out $20\,\mathrm{out}_i$ gives true-versus-false logit margin $\ge 20(1 - 2/128) > 19$ at every supervised position.
\end{proof}

Together with the weight constructions in \cref{app:subspace-construction}, \cref{lem:exact-hard,lem:softmax-leakage} prove \cref{prop:represent}.

\par\bigskip\noindent\textbf{Proposition~\ref{prop:represent} (Shared expressivity; formal statement).}\enspace{\itshape
For every $|\Sigma|$ and $n \ge 2$, a single Transformer as in \cref{ass:arch} with linear widths ($d_k = 2n+6$, $\dmodel = 2|\Sigma| + d_k + 11$, $\dmlp = \max\{|\Sigma|, 6(n+1)\}$) admits parameter settings $\thstep$ and $\thout$ under which greedy decoding generates the exact step-by-step and outcome-only continuations on every prompt with uniform $\Omega(1)$ logit margin.}\par\smallskip

\begin{proof}[Proof of \cref{prop:represent}]
Combine the shared-architecture weight realizations of \cref{app:subspace-construction} with \cref{lem:exact-hard,lem:softmax-leakage}: at every supervised step and terminal answer position, the true token exceeds the competing token by logit margin $> 19 > 0$, so greedy argmax decoding generates the exact continuation on every prompt. By continuity of $\theta \mapsto \TF_\theta(x)$ on the compact finite prompt set $\Sigma^{n+1}\times\{0,1\}$, the uniform $\Omega(1)$ margin persists on an open neighborhood of $\thstep$ and $\thout$.
\end{proof}

% ===================================================================
\section{Proof of \texorpdfstring{\Cref{thm:outcome-end}}{Theorem 1}}
\label{app:damping-proof}
\label{app:mainproofs}

\paragraph{Proof roadmap for \cref{thm:outcome-end}.}
Throughout this section \cref{ass:arch,ass:init} are in force. The proof separates a general optimization principle from the Transformer-specific correlation decay:
\begin{enumerate}[leftmargin=*,nosep,topsep=3pt]
  \item \textbf{Gradient split and Markov contraction (\cref{app:contraction}).} \Cref{lem:gradient-split} decomposes $\nabla_\theta\Lout$ into centering drift minus $\nabla_\theta\outcredit$, and \cref{lem:target-contraction} proves that any feature reading $|S|$ context letters contracts by $\mixrate^{n-|S|}$ against $y_n - \bary$.
  \item \textbf{Score interpolation and letter support (\cref{app:network}).} Scaling only the bilinear content scores by $\zeta \in [0,1]$ \cref{eq:score-split}, \cref{lem:frozen-relu-zeta} freezes all ReLU masks on $\mathcal U_{\initscale}$, and \cref{prop:network,prop:score-degree} show that order-$k$ interactions carry weight $O((\attconst\initscale^4)^k)$ while reading at most $2k + 1$ letters.
  \item \textbf{$\brel$-uniform geometric summation (\cref{app:thm1}).} \Cref{thm:main} sums $\sum_k \mixrate^{n-2k-1}(\attconst\initscale^4)^k$ uniformly across $\mathcal D_{\initscale}$ (arbitrary $\brel$), yielding $\Delta := \sup_{\mathcal D_{\initscale}}|\outcredit(\theta)| \le C_2\initscale^2\mixrate^n$.
  \item \textbf{General energy barrier and exit time (\cref{app:end}).} \Cref{lem:energy-barrier} establishes $T_{\mathrm{exit}} \ge R^2/\Delta$ (and $J\stepsize > R^2/(2\Delta)$ for projected GD) for any projected gradient dynamics with $F(\theta) \ge F(\theta_0) - \Delta$ on $\bar B(\theta_0, R)$; plugging in $R = \ballradius\initscale$ and $\Delta = C_2\initscale^2\mixrate^n$ immediately yields \textbf{\cref{thm:outcome-end}} and \textbf{\cref{cor:reversible}}.
\end{enumerate}

\paragraph{Parameter sets and standing constants.}
Write $\vartheta$ for all parameters other than the relative lag biases $\brel$; thus $\vartheta$ contains every weight, embedding, readout, and MLP bias. Besides $\initball$ of \cref{ass:init} we use
\[
  \initball' := \{\theta : \|\theta - \theta_0\|_2 \le 2\ballradius\initscale,\; \brel = 0\}, \qquad
  \mathcal D_{\initscale} := \{\theta : \|\vartheta - \vartheta_0\|_2 \le 2\ballradius\initscale\}, \qquad
  \mathcal U_{\initscale} := \{\theta : \|\vartheta - \vartheta_0\|_2 < 4\ballradius\initscale\},
\]
where $\brel$ is unrestricted in the last two, so $\initball \subset \initball' \subset \mathcal D_{\initscale} \subset \mathcal U_{\initscale}$. We prove the bounds of \cref{app:network,app:thm1} on $\mathcal U_{\initscale}$ or $\mathcal D_{\initscale}$ for arbitrary $\brel$. Write $H := 3$ for the number of heads per block and set $M := M_0 + \frac{1}{2}$, $\mu := M + 1 = M_0 + \frac{3}{2}$:
\begin{enumerate}[label=(\Alph*),leftmargin=*,nosep,topsep=2pt]
\item[(N1)] \emph{Block norms.} Since $\ballradius \le \frac{1}{8}$ (\cref{rem:norms}), every $\theta \in \mathcal U_{\initscale}$ satisfies $\|\vartheta - \vartheta_0\|_2 < \frac{1}{2}\initscale$, so every weight, embedding, and readout block $W$ satisfies $\|W\|_{\op} \le \|\initscale\widetilde W\|_{\op} + \|W - \initscale\widetilde W\|_{\frob} \le M\initscale$.
\item[(N2)] \emph{MLP biases.} Every bias coordinate satisfies $|b_m| \ge \initscale|\tilde b_m| - \frac{1}{2}\initscale \ge \frac{1}{2}\initscale$, and $\|b\|_2 \le \initscale N_b + \frac{1}{2}\initscale$.
\item[(E)] \emph{Scale condition.} Fix $c$ in $\initscale_0 = c\min(1, N_b^{-1})$ of \cref{rem:norms} as $c := (16 H \mu^2)^{-1}$. Then $\initscale \le \initscale_0$ implies
\begin{equation}
  \text{(E1)}\;\; \initscale N_b \le \frac{1}{2}
  \qquad\text{and}\qquad
  \text{(E2)}\;\; 16 H \mu^2\initscale \le 1,
  \label{eq:scale-E}
\end{equation}
giving $\|b\|_2 \le 1$, $H\mu\initscale \le \frac{1}{16\mu} \le \frac{1}{24}$, and $\initscale \le \frac{1}{108}$.
\item[(In)] \emph{Inputs.} At position $j$, $x_j := h_j^{(0)} = E_{\mathrm{tok}} e_{\mathrm{tok}(j)} + E_{\mathrm{pos}} e_j$ satisfies $\|x_j\|_2 \le 2M\initscale$ and $\|\partial_v x_j\|_2 \le 2\|v\|_2$.
\item[(A)] \emph{Positional weights.} For block $\ell$, head $h$, and query position $j$, $A_{jj'}^{(\ell,h)} := \cS(\brel_{j-\cdot}^{(\ell,h)})_{j'}$ ($j' \le j$) are nonnegative, sum to $1$ over $j'$, and do not depend on the prompt; at $\brel = 0$, $A_{jj'}^{(\ell,h)} = 1/m_j$ with $m_j := j + 1$.
\end{enumerate}
The answer position is $i = n + 3$, and the context letters $w_1,\ldots,w_n$ sit at positions $2,\ldots,n+1$ (\cref{app:subspace-construction}).

\subsection{Conditional Credit Contraction}
\label{app:assumptions}
\label{app:target-spectrum}
\label{app:contraction}
\label{rem:two-block-bias}

\Cref{lem:gradient-split} separates the label-dependent correlation term in $\nabla_\theta\Lout$, and \cref{lem:target-contraction} bounds Markov parity decay across unread context letters.

\begin{lemma}[Cross-entropy gradient decomposition]
\label{lem:gradient-split}
For any differentiable model with logits $z_\theta \in \R^2$, predicted distribution $\hat y_\theta := \cS(z_\theta)$, uniform vector $\bary := \frac{1}{2}\one$, and outcome credit $\outcredit(\theta) := \E[\langle z_\theta, y_n - \bary\rangle]$,
\begin{equation}
 \nabla_\theta\Lout(\theta) = \E\bigl[\Jz^\top(\hat y_\theta - \bary)\bigr] - \nabla_\theta\outcredit(\theta), \qquad
 \nabla_\theta\outcredit(\theta) = \E\bigl[\Jz^\top(y_n - \bary)\bigr].
 \label{eq:grad-split-display}
\end{equation}
\end{lemma}
\begin{proof}[Proof of \cref{lem:gradient-split}]
Since $-\log\ip{\cS(z)}{y_n} = \log\sum_a e^{z_a} - \ip{z}{y_n}$ and $\nabla_z\log\sum_a e^{z_a} = \cS(z)$, differentiating with respect to $z$ gives $\nabla_z(-\log\ip{\cS(z)}{y_n}) = \cS(z) - y_n$. Because the prompt distribution has finite support, differentiation and expectation commute; applying the chain rule with $\hat y_\theta - y_n = (\hat y_\theta - \bary) - (y_n - \bary)$ yields
\begin{align*}
  \nabla_\theta\Lout(\theta)
  &= \E\bigl[\Jz^\top(\hat y_\theta - y_n)\bigr] \\
  &= \E\bigl[\Jz^\top(\hat y_\theta - \bary)\bigr] - \E\bigl[\Jz^\top(y_n - \bary)\bigr],
\end{align*}
which is \cref{eq:grad-split-display}.
\end{proof}

Each transition $s_t = s_{t-1} + b_t \pmod 2$ multiplies $(-1)^{s_{t-1}}$ by $\chi_t = (-1)^{b_t}$. Because $\E[\chi_t] = 1 - 2p$ has modulus $\mixrate = |1-2p|$, the single-step conditional variance equals
\begin{equation}
 \E\Bigl[\bigl|(-1)^{s_t} - (1 - 2p)\,(-1)^{s_{t-1}}\bigr|^2 \Bigm| s_{t-1}\Bigr] = \operatorname{Var}(\chi_t) = 4p(1-p) = 1 - \mixrate^2.
 \label{eq:move-variance}
\end{equation}
Since $d_t = y_t - q(s_{t-1}) = (b_t - p)(e(1-s_{t-1}) - e(s_{t-1})) = \pm\frac{1}{\sqrt 2}(\chi_t - \E[\chi_t])e_-$, \cref{eq:move-variance} yields the horizon-free local innovation covariance
\begin{equation}
  \operatorname{Cov}(d_t) = \frac{1 - \mixrate^2}{2}\,e_- e_-^\top.
  \label{eq:local-variation}
\end{equation}
Over $n$ steps, each unread letter contracts the conditional expectation of $(-1)^{s_n}$ by $\mixrate$.

\begin{lemma}[Conditional contraction across unseen letters]
\label{lem:target-contraction}
In the counting chain of \cref{sec:data} with $s_0$ and $w^\star$ fixed, let $G := y_n - \bary \in \R^2$, and let $\mathcal F_S := \sigma(w_t : t \in S)$ for $S \subseteq [n]$. Then
\begin{equation}
 \bigl\|\E[G \mid \mathcal F_S]\bigr\|_2^2 = \frac{1}{2}\,\mixrate^{2(n-|S|)}, \qquad
 \bigl\|\E[H^\top G]\bigr\|_2 \le \mixrate^{n-|S|}\,\|H\|_{L^2}
 \label{eq:horizon}
\end{equation}
for every square-integrable $\mathcal F_S$-measurable matrix $H$ with two rows, where $\|H\|_{L^2} := (\E\|H\|_{\frob}^2)^{1/2}$, and the exponent $n - |S|$ is sharp.
\end{lemma}
\begin{proof}[Proof of \cref{lem:target-contraction}]
Because $G = y_n - \bary = \frac{1}{2}(-1)^{s_n}(e_0 - e_1)$, $(-1)^{s_n} = (-1)^{s_0}\prod_{t=1}^n \chi_t$, and the unread signs $\chi_t$ ($t \notin S$) are independent of $\mathcal F_S$ with mean $1 - 2p$, conditioning on $\mathcal F_S$ yields
\[
  \E[G \mid \mathcal F_S]
  = \frac{1}{2}(-1)^{s_0}(1-2p)^{n-|S|}\Biggl(\prod_{t\in S}\chi_t\Biggr)(e_0 - e_1).
\]
Since $|\chi_t| = 1$ and $\|e_0 - e_1\|_2^2 = 2$, taking the squared Euclidean norm pointwise gives $\|\E[G \mid \mathcal F_S]\|_2^2 = \frac{1}{2}\mixrate^{2(n-|S|)}$, proving the first identity in \cref{eq:horizon}.

Next, fix any unit vector $u$. Since $H u$ is $\mathcal F_S$-measurable and $\|H u\|_2 \le \|H\|_{\frob}$, the tower property and Cauchy--Schwarz give
\[
  u^\top\E[H^\top G]
  = \E\bigl[\langle H u, \E[G \mid \mathcal F_S]\rangle\bigr]
  \le \bigl(\E\|H u\|_2^2\bigr)^{1/2}\bigl(\E\|\E[G \mid \mathcal F_S]\|_2^2\bigr)^{1/2}
  \le \frac{1}{\sqrt 2}\,\mixrate^{n-|S|}\|H\|_{L^2},
\]
proving the second bound in \cref{eq:horizon}. Finally, taking $H := (-1)^{s_0}(\prod_{t\in S}\chi_t)e_-$ gives $\|H\|_{L^2} = 1$ and $|\E[H^\top G]| = \frac{1}{\sqrt 2}\mixrate^{n-|S|}$, establishing sharpness of $n - |S|$.
\end{proof}
\label{rem:tight}

\subsection{Attention Score Expansion}
\label{app:model-spectrum}
\label{app:network}

Expanding the network directly in powers of $\theta$ around $0$ mixes parameter scale with token degree because the $O(\initscale)$ MLP biases and residual skip paths create terms of arbitrary depth that read no new letters. Instead, we interpolate only the bilinear content scores $\psi_{il}^{(\ell)} \mapsto \zeta\psi_{il}^{(\ell)}$ for $\zeta \in [0,1]$, keeping the lag biases $\brel$. Let $z(\zeta)$ denote the output logits at interpolation parameter $\zeta$, so $z(1) = z_\theta$ and $z(0)$ replaces every attention matrix by the input-independent positional averaging $A^{(\ell,h)}$ of item~(A). Taylor's theorem around $\zeta = 0$ to order $k_0 \ge 0$ gives
\begin{equation}
 z_\theta = \sum_{k=0}^{k_0} z_k + z_{\mathrm{rem}}, \qquad
 z_k := \frac{1}{k!}\frac{\partial^k z}{\partial\zeta^k}(0), \qquad
 z_{\mathrm{rem}} := \int_0^1 \frac{(1 - \zeta)^{k_0}}{k_0!}\frac{\partial^{k_0+1} z}{\partial\zeta^{k_0+1}}(\zeta)\,\rmd\zeta,
 \label{eq:score-split}
\end{equation}
with parameter Jacobians $J_k := \partial_\theta z_k$ and $J_{\mathrm{rem}} := \partial_\theta z_{\mathrm{rem}}$. Write $x_i := h_i^{(0)}$ for the token-plus-position embedding at site $i$, and $\partial_v F := \frac{\partial}{\partial s}F(\theta + sv)\big|_{s=0}$ for the directional derivative along $v$.

\begin{lemma}[Frozen-ReLU affine reduction and softmax derivative bound]
\label{lem:frozen-relu-zeta}
Uniformly over $\theta \in \initball'$, positions $i$, inputs, and $\zeta \in [0,1]$:
\begin{enumerate}[label=(\alph*),leftmargin=*]
  \item Every MLP activation mask $\Lambda_\ell := \operatorname{diag}(\one[b_{\mathrm{up}}^{(\ell)} > 0])$ is constant, reducing block $\ell \in \{1,2\}$ to the affine map $h_i^{(\ell)} = M_\ell(h_i^{(\ell-1)} + a_i^{(\ell)}) + o_\ell$ with $M_\ell := I + W_{\mathrm{down}}^{(\ell)}\Lambda_\ell W_{\mathrm{up}}^{(\ell)}$, $o_\ell := W_{\mathrm{down}}^{(\ell)}\Lambda_\ell b_{\mathrm{up}}^{(\ell)}$, $\|M_\ell - I\|_{\op} \le C\initscale^2$, $\|o_\ell\|_2 \le C\initscale$, and $\|\partial_v o_\ell\|_2 \le C\|v\|_2$.
  \item For every base score vector $s \in \R^{i+1}$, integer $k \ge 1$, and directions $u_1, \ldots, u_k \in \R^{i+1}$, the $k$-th directional derivative of $\alpha(s) = \cS(s)$ satisfies
  \begin{equation}
    \bigl\|D^k\alpha(s)[u_1, \ldots, u_k]\bigr\|_1 \le 2^k k!\prod_{a=1}^k \|u_a\|_\infty.
    \label{eq:softmax-avg}
  \end{equation}
\end{enumerate}
\end{lemma}
\begin{proof}[Proof of \cref{lem:frozen-relu-zeta}]
We prove~(a) on the larger open set $\mathcal U_{\initscale}$ for arbitrary $\brel$, with $C := 2\mu^2$.

\proofstep{Step 1: Hidden-state norms under arbitrary masks.}
By~(N1), $\|x_j\|_2 \le 2M\initscale$, and $H M\initscale \le \frac{1}{16}$ from \cref{eq:scale-E}. For any row-stochastic attention weights, we obtain $\|\tilde h_i^{(1)}\|_2 \le \frac{17}{16}\cdot 2M\initscale \le 3M\initscale$. Since $\|\relu(u)\|_2 \le \|u\|_2$, $\|b\|_2 \le 1$, and $3M^2\initscale^2 \le 1$, the first MLP adds at most $M\initscale(3M^2\initscale^2+1) \le 2M\initscale$, yielding
\[
  \bigl\|\tilde h_i^{(1)}\bigr\|_2 \le 3M\initscale, \qquad
  \bigl\|h_i^{(1)}\bigr\|_2 \le 5M\initscale, \qquad
  \bigl\|\tilde h_i^{(2)}\bigr\|_2 \le \frac{17}{16}\cdot 5M\initscale \le 6M\initscale.
\]

\proofstep{Step 2: Frozen ReLU masks and affine block bounds.}
For every unit $m$ and block $\ell \in \{1,2\}$, combining~(N1), Step~1, $6M^2\initscale \le \frac{6}{16H} < \frac{1}{2}$ from \cref{eq:scale-E}, and~(N2) gives
\[
  \bigl|\bigl(W_{\mathrm{up}}^{(\ell)}\tilde h_i^{(\ell)}\bigr)_m\bigr|
  \le 6M^2\initscale^2 < \frac{1}{2}\initscale \le |b_m|.
\]
Hence $\relu(W_{\mathrm{up}}^{(\ell)}\tilde h_i^{(\ell)} + b_{\mathrm{up}}^{(\ell)}) = \Lambda_\ell(W_{\mathrm{up}}^{(\ell)}\tilde h_i^{(\ell)} + b_{\mathrm{up}}^{(\ell)})$ throughout $\mathcal U_{\initscale}$. Substituting into \cref{eq:block-Equations} yields
\begin{align*}
  \|M_\ell - I\|_{\op} &\le M^2\initscale^2, \qquad\qquad \|o_\ell\|_2 \le M\initscale,\\
  \|\partial_v o_\ell\|_2 &= \bigl\|(\partial_v W_{\mathrm{down}})\Lambda_\ell b + W_{\mathrm{down}}\Lambda_\ell\partial_v b\bigr\|_2 \le \|v\|_2\bigl(\|b\|_2 + M\initscale\bigr) \le 2\|v\|_2,
\end{align*}
completing part~(a).

\proofstep{Step 3: Softmax higher-order derivatives.}
Using $\partial_u\alpha_j = \alpha_j u_j - \alpha_j\sum_{j'}\alpha_{j'}u_{j'}$, induction on $m$ expands $D^m\alpha(s)[u_1,\ldots,u_m]_j$ into at most $2^m m!$ signed products of the form $\pm\alpha_j\alpha_{j_1}\cdots\alpha_{j_q}\prod_{a=1}^m u_{a,\iota_a}$. Summing absolute values over $j$ and the free indices via $\sum_j\alpha_j = 1$ gives \cref{eq:softmax-avg}.
\end{proof}

\paragraph{Term expansions and majorant series.}
Fix $\theta \in \mathcal U_{\initscale}$. By \cref{lem:frozen-relu-zeta}(a), for real $\zeta \in [0,1]$ and real $s$ near $0$, the network at $\theta + s v$ with content scores scaled by $\zeta$ reduces to the frozen analytic network
\begin{align}
  \alpha_{j\cdot}^{(\ell,h)} &= \cS\bigl(\brel_{j-\cdot}^{(\ell,h)} + s v_{\brel\,j-\cdot}^{(\ell,h)} + \zeta\psi_{j\cdot}^{(\ell,h)}\bigr), &
  a_j^{(\ell)} &= \sum_{h=1}^H \sum_{j'\le j} \alpha_{jj'}^{(\ell,h)} W_V^{(\ell,h)} h_{j'}^{(\ell-1)}, \nonumber\\
  h_j^{(\ell)} &= M_\ell(h_j^{(\ell-1)} + a_j^{(\ell)}) + o_\ell, &
  z_j &= V h_j^{(2)},
  \label{eq:frozen-net}
\end{align}
where every parameter block, $M_\ell$, and $o_\ell$ is evaluated at $\theta + sv$, and $v_\brel$ is the $\brel$-part of $v$.

Expanding \cref{eq:frozen-net} as a double power series $Q = \sum_{k,m\ge 0}\zeta^k s^m Q_{k,m}$, each coefficient $Q_{k,m}$ is an absolutely convergent sum of terms built from entries of $A_{j\cdot}^{(\ell,h)}$ and $v_\brel$, parameter blocks of $\theta$ or $v$, and inputs $x_j$ or $\partial_v x_j$. We say that a term \emph{reads} position $l$ if $x_l$ or $\partial_v x_l$ is one of its factors, and write $Q^{\ni l}$ for the sum of terms of $Q$ reading $l$. For a scalar series $\widehat Q(\zeta, s) = \sum q_{k,m}\zeta^k s^m$ with $q_{k,m} \ge 0$, we write $Q \preccurlyeq \widehat Q$ iff
\[
  \|Q_{k,m}\|_{\mathrm{abs}} := \sum_{\text{terms } T \text{ of } Q_{k,m}} \sup_{\|v\|_2=1,\, w} \|T\| \le q_{k,m} \qquad \forall k, m \ge 0,
\]
so that $z_k = z_{k,0}$ and $\partial_v z_k = z_{k,1}$.

\begin{proposition}[Blockwise bounds for the score expansion]
\label{prop:network}
In the Transformer of \cref{ass:arch} under \cref{ass:init}, uniformly over $\theta \in \mathcal D_{\initscale}$ (arbitrary $\brel$), positions $i$, and inputs, there exist constants $C, \attconst > 0$ independent of $n, k, \brel$, and widths such that:
\begin{enumerate}[label=(\roman*),leftmargin=*]
\item with input-independent bias offsets $O^{(1)} = o_1$ and $O^{(2)} = o_1 + o_2$, we have $h_i^{(\ell)} = x_i + O^{(\ell)} + O(\initscale^2)$, $|\psi_{il}^{(\ell)}| \le C\initscale^4$, $\|z_i\|_2 \le C\initscale^2$, and $\|\partial_v z_i\|_2 \le C\initscale\|v\|_2$;
\item for all $k \ge 1$, $\zeta \in [0,1]$, and $\|v\|_2 = 1$,
\[
  \frac{1}{k!}\Bigl\|\frac{\partial^k z_i}{\partial\zeta^k}\Bigr\|_2 \le C\initscale^3(\attconst\initscale^4)^k
  \qquad\text{and}\qquad
  \frac{1}{k!}\Bigl\|\partial_v\frac{\partial^k z_i}{\partial\zeta^k}\Bigr\|_2 \le C\initscale^2(\attconst\initscale^4)^k
\]
(and likewise for $\|\cdot\|_{\mathrm{abs}}$ at $\zeta = 0$ and for $h^{(2)}_i$ with one fewer factor of $\initscale$);
\item at $\zeta = 0$, the network is affine in $(x_j)$, attention-transported logit terms have tangent $O(\initscale^2)$, direct and bias terms have tangent $O(\initscale)$, and every order-$k$ monomial contains at most $2k + 1$ inputs $x_j$; and
\item when $\brel = 0$ ($\theta \in \initball'$) under full-prefix causal attention, monomials of the order-$k$ coefficient at $\zeta = 0$ at position $i$ containing an earlier input $x_l$ ($l < i$) satisfy the bounds of~(ii) multiplied by
\begin{equation}
 d_{l,i} := \frac{1}{m_i}\Biggl(1 + \sum_{l < j \le i} \frac{1}{m_j}\Biggr), \qquad m_i := i + 1.
 \label{eq:dilution}
\end{equation}
\end{enumerate}
\end{proposition}
\begin{proof}[Proof of \cref{prop:network}]
We prove all parts for every $\theta \in \mathcal U_{\initscale}$ and every $\brel$ (only \cref{eq:dilution} uses $\brel = 0$), setting
\begin{equation}
  R := \frac{\rho}{\initscale^4}, \qquad \rho := \frac{1}{256\mu^4}, \qquad \attconst := \frac{2}{\rho} = 512\mu^4, \qquad C_z := 120 H\mu^3,
  \label{eq:majorant-consts}
\end{equation}
so $\mu\initscale \le \frac{1}{24}$ gives $R \ge 24^4/256 > 2$.

\proofstep{Step 1: Majorant calculus rules.}
Let $\bar\sigma(u) := e^u/(2 - e^u)$ and $\bar\sigma_2(u) := e^{2u}/(2 - e^u)^2$. Four elementary majorant rules hold:
\begin{itemize}[leftmargin=*]
\item \textbf{(M1) Algebra and letter tracking:} If $Q \preccurlyeq \widehat Q$ and $Q' \preccurlyeq \widehat Q'$, then $Q + Q' \preccurlyeq \widehat Q + \widehat Q'$, $QQ' \preccurlyeq \widehat Q\widehat Q'$, and $(QQ')^{\ni l} \preccurlyeq \widehat Q^{\ni l}\widehat Q' + \widehat Q\widehat Q'^{\ni l}$.
\item \textbf{(M2) Base blocks:} $W + s v_W \preccurlyeq M\initscale + s$, $x_j + s\partial_v x_j \preccurlyeq 2M\initscale + 2s$, $M_\ell \preccurlyeq 1 + (M\initscale + s)^2$, and $o_\ell \preccurlyeq (M\initscale + s)(1 + s)$.
\item \textbf{(M3) Softmax expansion:} For zero-constant-term $w_{j'} \preccurlyeq \hat w$ and $Q_{j'} \preccurlyeq \widehat Q$ with $A := A_{j\cdot}^{(\ell,h)}$, expanding $\cS(\brel_{j-\cdot}^{(\ell,h)} + w)_{j'} = \sum_{r\ge 0}(-1)^r A_{j'}e^{w_{j'}}(\sum_{l'}A_{l'}(e^{w_{l'}}-1))^r$ yields
\begin{align*}
  \sum_{j'}\cS(\brel_{j-\cdot}^{(\ell,h)}+w)_{j'}Q_{j'} &\preccurlyeq \bar\sigma(\hat w)\widehat Q,\\
  \Bigl(\sum_{j'}\cS(\brel_{j-\cdot}^{(\ell,h)}+w)_{j'}Q_{j'}\Bigr)^{\ni l}
  &\preccurlyeq \bar\sigma(\hat w)\sum_{j'}A_{j'}\bigl(\widehat Q_{j'}^{\ni l} + \hat w_{j'}^{\ni l}\widehat Q\bigr) + \bar\sigma_2(\hat w)\widehat Q\sum_{j'}A_{j'}\hat w_{j'}^{\ni l}.
\end{align*}
\item \textbf{(M4) Cauchy coefficient bound:} If $Q \preccurlyeq \widehat Q$ and $\widehat Q(R', S') < \infty$, then $\|Q_{k,m}\|_{\mathrm{abs}} \le \widehat Q(R', S')R'^{-k}S'^{-m}$.
\end{itemize}

\proofstep{Step 2: Recursive majorants of the frozen network.}
Set $\hat h^{(0)} := 2M\initscale + 2s$ and define recursively for $\ell \in \{1,2\}$:
\begin{align*}
  \hat\psi^{(\ell)} &:= (M\initscale + s)^2\bigl(\hat h^{(\ell-1)}\bigr)^2, \qquad
  \hat w^{(\ell)} := s + \zeta\hat\psi^{(\ell)}, \qquad
  \hat a^{(\ell)} := H\bar\sigma(\hat w^{(\ell)})(M\initscale + s)\hat h^{(\ell-1)},\\
  \hat h^{(\ell)} &:= \bigl(1 + (M\initscale + s)^2\bigr)\bigl(\hat h^{(\ell-1)} + \hat a^{(\ell)}\bigr) + (M\initscale + s)(1 + s),\\
  \hat z &:= (M\initscale + s)\hat h^{(2)}, \qquad
  \hat z_{\mathrm{att}} := (M\initscale + s)\bigl(1 + (M\initscale + s)^2\bigr)\Bigl[\bigl(1 + (M\initscale + s)^2\bigr)\hat a^{(1)} + \hat a^{(2)}\Bigr].
\end{align*}
By~(M1)--(M3), $\psi_{jj'}^{(\ell,h)} \preccurlyeq \hat\psi^{(\ell)}$, $a_j^{(\ell)} \preccurlyeq \hat a^{(\ell)}$, $h_j^{(\ell)} \preccurlyeq \hat h^{(\ell)}$, $z_j \preccurlyeq \hat z$, and $\|z_{k,m}\|_{\mathrm{abs}} \le [\zeta^k s^m]\hat z_{\mathrm{att}}$ for all $k \ge 1$. Evaluating at $(\zeta, s) = (R, \initscale)$ (where $M\initscale + s = \mu\initscale$) via \cref{eq:scale-E} gives $\hat w^{(1)}, \hat w^{(2)} \le \frac{1}{2}$ ($\bar\sigma(\frac{1}{2}) < 5$) and
\[
  \hat h^{(1)} \le 8\mu\initscale, \qquad
  \hat h^{(2)} \le 26\mu\initscale, \qquad
  \hat\psi^{(2)} \le 64\mu^4\initscale^4, \qquad
  \hat z \le 26\mu^2\initscale^2, \qquad
  \hat z_{\mathrm{att}} \le C_z\initscale^3.
\]

\proofstep{Step 3: Proof of parts (i)--(iii).}
Evaluating at $(\zeta, s) = (1, 0)$ with $R > 2$ and applying~(M4) at $(R, \initscale)$ gives $|\psi_{il}^{(\ell)}| \le 64\mu^4\initscale^4$, $\|z_i\|_2 \le 26\mu^2\initscale^2$, $\|\partial_v z_i\|_2 \le \hat z(R,\initscale)/\initscale \le 26\mu^2\initscale$, and $h_i^{(\ell)} = x_i + O^{(\ell)} + O(\initscale^2)$. Moreover, using $\binom{j}{k}\zeta^{j-k}(R/2)^k \le R^j$ for $\zeta \in [0,1]$, we obtain for $m \in \{0,1\}$:
\[
  \Bigl\|\frac{1}{k!}\frac{\partial^k}{\partial\zeta^k}\partial_v^m z_i(\zeta)\Bigr\|_2
  \le \Bigl(\frac{2}{R}\Bigr)^k \frac{\hat z_{\mathrm{att}}(R,\initscale)}{\initscale^m}
  \le C_z\initscale^{3-m}(\attconst\initscale^4)^k,
\]
proving~(i)--(ii). At $\zeta = 0$, $z_{0,0}$ is affine in $(x_j)$, with attention-transported tangent $\|z_{0,1}\|_{\mathrm{abs}} \le C_z\initscale^2$ and direct/bias tangent $\le 26\mu^2\initscale$; inducting over \cref{eq:frozen-net} shows that $[\zeta^k]h_j^{(\ell)}$ reads at most $2k+1$ positions and $[\zeta^k](s v_\brel + \zeta\psi)$ reads at most $2k$ positions, so $z_{k,m}$ reads at most $2k + 1$ letters, proving~(iii).

\proofstep{Step 4: Prefix dilution bound in part (iv).}
Fix $l < i$ and let $g_j^{(\ell)}$ denote the value at $(R,\initscale)$ of the majorant of $(h_j^{(\ell)})^{\ni l}$, starting from $g_j^{(0)} = 2\mu\initscale\,\one[j = l]$. Applying~(M3) with $\bar\sigma(\frac{1}{2}) \le 5$ and $\bar\sigma_2(\frac{1}{2}) < 23$ yields
\[
  g_j^{(\ell)} \le 3 g_j^{(\ell-1)} + 24\mu\initscale\sum_{h=1}^H\sum_{j'\le j} A_{jj'}^{(\ell,h)} g_{j'}^{(\ell-1)}.
\]
Unrolling $\ell = 1, 2$ gives $g_i^{(2)} \le 144\mu^2\initscale^2\omega_l(i)$ with $\omega_l(i) := \sum_h A_{il}^{(1,h)} + \sum_h A_{il}^{(2,h)} + \sum_{h,h'}(A^{(2,h)}A^{(1,h')})_{il}$, and hence $\|(z_{k,m})^{\ni l}\|_{\mathrm{abs}} \le 144\mu^3\initscale^{3-m}\omega_l(i)(\initscale^4/\rho)^k$. At $\brel = 0$, $A_{jl}^{(\ell,h)} = 1/m_j$ gives $\omega_l(i) \le 3H^2 d_{l,i} = 27 d_{l,i}$, proving~(iv) with $C := 27 C_z$.
\end{proof}

Grouping the power-series terms of \cref{prop:network} by the exact subset of context letters $S \subseteq [n]$ they read translates the majorant bounds directly into degree-graded Jacobian bounds.

\begin{proposition}[Score count and letter support]
\label{prop:score-degree}
Fix $s_0$ and $w^\star$. Uniformly over $\theta \in \mathcal D_{\initscale}$ (arbitrary $\brel$), $\Jz = \sum_{k=0}^{k_0} J_k + J_{\mathrm{rem}}$, and there exist constants $C_\star, \attconst > 0$ depending only on $M_0$ and $\embconst$ (independent of $n, k, k_0, \brel$, and widths) such that:
\begin{enumerate}[label=(\roman*),leftmargin=*]
\item $J_0 = J_{0,\emptyset} + \sum_{t=1}^n J_{0,t}$, where $J_{0,\emptyset}$ reads no context letter, $J_{0,t}$ reads only $w_t$, $\sup \|J_{0,\emptyset}\|_{\frob} \le C_\star\initscale$, and $\sum_{t=1}^n \sup \|J_{0,t}\|_{\frob} \le C_\star\initscale^2$;
\item for each $k \ge 1$, $J_k = \sum_r T_{k,r}$ where each monomial $T_{k,r}$ reads at most $2k + 1$ letters and $\sum_r \sup \|T_{k,r}\|_{\frob} \le C_\star\initscale^2(\attconst\initscale^4)^k$; and
\item $\sup \|J_{\mathrm{rem}}\|_{\frob} \le C_\star\initscale^2(\attconst\initscale^4)^{k_0+1}$.
\end{enumerate}
The same bounds hold for $z_k$ and $z_{\mathrm{rem}}$ with an extra factor of $\initscale$.
\end{proposition}
\begin{proof}[Proof of \cref{prop:score-degree}]
Group the terms of $z_{k,0}$ by the exact letter subset $S \subseteq [n]$ they read ($|S| \le 2k+1$ by Step~3 of \cref{prop:network}), setting $T_{k,S} := \partial_\theta z_{k,S}$. Since $T_{k,S}$ has two rows, we have $\|T_{k,S}\|_{\frob} \le \sqrt 2\,\|T_{k,S}\|_{\op} = \sqrt 2\sup_{\|v\|_2=1}\|\partial_v z_{k,S}\|_2$, and therefore
\[
  \sum_{S\subseteq [n]} \sup\|T_{k,S}\|_{\frob} \le \sqrt 2\,\|z_{k,1}\|_{\mathrm{abs}}.
\]
At $k = 0$, we decompose $z_0 = z_{0,\emptyset} + \sum_{t=1}^n z_{0,\{t\}}$ with $\sup\|J_{0,\emptyset}\|_{\frob} \le 26\sqrt 2\,\mu^2\initscale$. Since $i = n+3 > 1+t$ and each $A_{i\cdot}^{(\ell,h)}$ is row-stochastic ($\sum_{t=1}^n \omega_{1+t}(i) \le H + H + H^2 = 15$), Step~4 of \cref{prop:network} gives
\[
  \sum_{t=1}^n \sup\|J_{0,t}\|_{\frob} \le 15\cdot 144\sqrt 2\,\mu^3\initscale^2
\]
uniformly over $\theta \in \mathcal U_{\initscale}$ and $\brel$. Parts~(ii)--(iii) and the corresponding $z$-bounds follow directly from Step~3 of \cref{prop:network} and \cref{eq:score-split} with $C_\star := 15\cdot 144\sqrt 2\,C_z$.
\end{proof}

\subsection{Uniform Credit Bound}
\label{app:thm1}

Combining the $\mixrate^{n-|S|}$ contraction bound (\cref{lem:target-contraction}, with $|S| \le 2k + 1$) and the $(\attconst\initscale^4)^k$ majorant bound on order-$k$ score interactions (\cref{prop:score-degree}) bounds total outcome credit over $\mathcal D_{\initscale}$. When $\attconst\initscale^4 \le \frac{1}{2}\mixrate^2$ ($\initscale \lesssim \mixrate$), successive orders decay by at least $1/2$ and sum to $O(\initscale^2\mixrate^n)$ uniformly in $\brel$.

\begin{lemma}[Degree-wise Markov contraction and $\brel$-uniform outcome credit]
\label{thm:main}
\label{lem:score-sum}
\label{lem:bias-uniform}
In the Transformer of \cref{ass:arch} under \cref{ass:init}, consider the counting chain \cref{eq:tracking,eq:lambda} conditioned on $(s_0, w^\star)$ (and hence unconditionally under any law on $(s_0, w^\star)$). There exist constants $c_2, C_2 > 0$ independent of $n, \initscale, \brel$, and widths such that:
\begin{enumerate}[label=(\roman*),leftmargin=*,itemsep=1pt,topsep=2pt]
\item \emph{Degree-wise contraction.} For every $S \subseteq [n]$ and square-integrable feature matrix $H$ measurable with respect to $\{w_t : t \in S\}$, $\|\E[H^\top(y_n - \bary)]\|_2 \le \mixrate^{n-|S|}\|H\|_{L^2}$, and the exponent $n - |S|$ is sharp.
\item \emph{$\brel$-uniform full-network credit.} For all $n \ge 2$, $\mixrate \in (0,1)$, $0 < \initscale \le \min\{\initscale_0, c_2\mixrate\}$ (in particular $\initscale \le c'\min\{\mixrate, \dmlp^{-1/2}\}$), and $\theta \in \mathcal D_{\initscale}$ (with $\brel$ unrestricted), bilinear score interactions $J_{\ge 1} := \Jz - J_0$ satisfy $\|\E[J_{\ge 1}^\top(y_n - \bary)]\|_2 \le C_2\initscale^6\mixrate^{n-3}$, and
\begin{equation}
  \|\nabla_\theta\outcredit(\theta)\|_2 \le C_2\initscale\mixrate^n, \qquad \lvert\outcredit(\theta)\rvert \le C_2\initscale^2\mixrate^n.
  \label{eq:thm1-full}
\end{equation}
\end{enumerate}
\end{lemma}
\begin{proof}[Proof of \cref{thm:main}]
Part~(i) is \cref{lem:target-contraction} conditionally and follows unconditionally by Jensen's inequality. For~(ii), set $c_2 := \min\{1, (2\attconst)^{-1/4}\}$ (so $\initscale \le c_2\mixrate \implies \attconst\initscale^4 \le \frac{1}{2}\mixrate^2$), $k_0 := \lceil(n-1)/2\rceil$, and $G := y_n - \bary$ ($\|G\|_2 = 1/\sqrt 2$). Conditioning on $(s_0, w^\star)$ and applying~(i) term-by-term to \cref{prop:score-degree}(i)--(iii) across $\mathcal D_{\initscale}$ gives:
\begin{align*}
  \bigl\|\E\bigl[(J_0)^\top G \mid s_0, w^\star\bigr]\bigr\|_2
  &\le C_\star\initscale\mixrate^n + C_\star\initscale^2\mixrate^{n-1} \le 2C_\star\initscale\mixrate^n \qquad (\text{since } \initscale \le \mixrate),\\
  \bigl\|\E\bigl[(J_k)^\top G \mid s_0, w^\star\bigr]\bigr\|_2
  &\le C_\star\attconst\initscale^6\mixrate^{n-3}\Bigl(\frac{\attconst\initscale^4}{\mixrate^2}\Bigr)^{k-1}
  \le C_\star\attconst\initscale^6\mixrate^{n-3}\,2^{-(k-1)} \qquad (1 \le k \le k_0),\\
  \bigl\|\E[J_{\mathrm{rem}}^\top G \mid s_0, w^\star]\bigr\|_2
  &\le C_\star\attconst\initscale^6\mixrate^{n-3}\,2^{-k_0} \qquad (\text{since } 2k_0 \ge n-1).
\end{align*}
Summing over $1 \le k \le k_0$ and the remainder gives $\|\E[J_{\ge 1}^\top G]\|_2 \le 3C_\star\attconst\initscale^6\mixrate^{n-3}$. Adding the order-$0$ bound (and repeating verbatim for $z_\theta$, which carries an extra factor of $\initscale$) yields \cref{eq:thm1-full} with $C_2 := (2 + 3\attconst)C_\star$, while $|\tilde b_m| \ge 1$ gives $N_b \ge \sqrt{\dmlp}$ and $\initscale_0 \le (16H\mu^2)^{-1}\dmlp^{-1/2}$.
\end{proof}

\subsection{Exit-Time Lower Bound}
\label{app:end}

To pass from $|\outcredit(\theta)| \le \Delta = C_2\initscale^2\mixrate^n$ on $\mathcal D_{\initscale}$ to an $\Omega(\mixrate^{-n})$ stay-time bound, \cref{lem:energy-barrier} gives a general exit-time inequality for projected gradient flow and projected gradient descent, \cref{lem:end-smooth} checks smoothness of $\Lout$, and \cref{thm:outcome-end,cor:reversible} follow.

\begin{lemma}[Energy barrier for projected gradient dynamics]
\label{lem:energy-barrier}
\label{rem:flow}
Let $F : \mathcal U \to \R$ be $C^1$ on an open set $\mathcal U \supset \bar B(\theta_0, R)$ in $\R^d$ ($R, \Delta > 0$) satisfying
\[
  F(\theta) \ge F(\theta_0) - \Delta \qquad \forall \theta \in \bar B(\theta_0, R).
\]
\begin{enumerate}[label=(\roman*),leftmargin=*,itemsep=1pt,topsep=2pt]
\item \emph{(Projected gradient flow.)} For any piecewise-constant orthogonal projection $P(\tau)$ ($P^\top = P = P^2$) and piecewise-$C^1$ trajectory of $\dot\theta(\tau) = -P(\tau)\nabla F(\theta(\tau))$ with $\theta(0) = \theta_0$, the first exit time $T_{\mathrm{exit}} := \inf\{\tau \ge 0 : \|\theta(\tau) - \theta_0\|_2 = R\}$ ($\inf\emptyset := \infty$) satisfies
\begin{equation}
  \|\theta(\tau) - \theta_0\|_2^2 \le \tau\bigl(F(\theta_0) - F(\theta(\tau))\bigr) \le \tau\Delta \quad (0 \le \tau < T_{\mathrm{exit}}), \qquad
  T_{\mathrm{exit}} \ge \frac{R^2}{\Delta}.
  \label{eq:energy-barrier}
\end{equation}
\item \emph{(Projected gradient descent.)} Suppose moreover that $\bar B(\theta_0, \frac{3}{2}R) \subset \mathcal U$ with $F(\theta) \ge F(\theta_0) - \Delta$, $\|\nabla F(\theta)\|_2 \le G$, and $|\partial_s^2 F(\theta + sv)|_{s=0}| \le L$ for all $\theta \in \bar B(\theta_0, \frac{3}{2}R)$ and $\|v\|_2 = 1$. Then for every step size $0 < \stepsize \le \min\{\frac{1}{L}, \frac{R}{2G}\}$, projected GD $\theta_{m+1} = \theta_m - \stepsize P_m\nabla F(\theta_m)$ satisfies $\|\theta_m - \theta_0\|_2^2 \le 2(m\stepsize)\Delta$ up to its first exit iterate $J := \inf\{m \ge 0 : \|\theta_m - \theta_0\|_2 > R\}$, hence $J\stepsize > \frac{R^2}{2\Delta}$.
\end{enumerate}
\end{lemma}
\begin{proof}[Proof of \cref{lem:energy-barrier}]
For~(i), away from switching instants, orthogonality gives $\ddtau{F(\theta(\tau))} = -\langle\nabla F(\theta), P(\tau)\nabla F(\theta)\rangle = -\|P(\tau)\nabla F(\theta)\|_2^2 = -\|\dot\theta(\tau)\|_2^2$. For $0 \le \tau \le T_{\mathrm{exit}}$ ($\tau < \infty$), Cauchy--Schwarz gives
\[
  \|\theta(\tau) - \theta_0\|_2^2
  = \Biggl\|\int_0^\tau \dot\theta(u)\,\rmd u\Biggr\|_2^2
  \le \tau\int_0^\tau \|\dot\theta(u)\|_2^2\,\rmd u
  = \tau\bigl(F(\theta_0) - F(\theta(\tau))\bigr)
  \le \tau\Delta;
\]
evaluating at $\tau = T_{\mathrm{exit}}$ when $T_{\mathrm{exit}} < \infty$ yields $R^2 \le T_{\mathrm{exit}}\Delta$.

For~(ii), while $\|\theta_m - \theta_0\|_2 \le R$, the step bound $\|\theta_{m+1} - \theta_m\|_2 \le \stepsize G \le \frac{1}{2}R$ keeps the line segment $[\theta_m, \theta_{m+1}]$ inside $\bar B(\theta_0, \frac{3}{2}R)$. Directional $L$-smoothness with $\stepsize L \le 1$ gives
\[
  F(\theta_{m+1})
  \le F(\theta_m) - \stepsize\Bigl(1 - \frac{\stepsize L}{2}\Bigr)\|P_m\nabla F(\theta_m)\|_2^2
  \le F(\theta_m) - \frac{1}{2\stepsize}\|\theta_{m+1} - \theta_m\|_2^2.
\]
Summing from $m = 0$ to $J - 1 < \infty$ with $\theta_J \in \bar B(\theta_0, \frac{3}{2}R)$ yields
\[
  R^2 < \|\theta_J - \theta_0\|_2^2
  \le J\sum_{m=0}^{J-1}\|\theta_{m+1} - \theta_m\|_2^2
  \le 2(J\stepsize)\bigl(F(\theta_0) - F(\theta_J)\bigr)
  \le 2(J\stepsize)\Delta,
\]
completing the proof.
\end{proof}

To apply the discrete-time clause \cref{lem:energy-barrier}(ii) to $\Lout$, we next record uniform first- and second-order directional derivative bounds across $\mathcal D_{\initscale}$.

\begin{lemma}[Smoothness of the outcome objective on $\mathcal D_{\initscale}$]
\label{lem:end-smooth}
Under \cref{ass:arch,ass:init} with $\initscale \le \initscale_0$, there exists $C_1 > 0$ independent of $n$ and $\brel$ such that for all $\theta \in \mathcal D_{\initscale}$ and all unit directions $v$ (including $\brel$),
\[
 \|\nabla_\theta\Lout(\theta)\|_2 \le C_1\initscale, \qquad
 \Biggl|\frac{\partial^2}{\partial s^2}\Lout(\theta + s v)\Big|_{s=0}\Biggr| \le 2C_1.
\]
\end{lemma}
\begin{proof}[Proof of \cref{lem:end-smooth}]
Set $C_1 := 52\mu^2$. For $\theta \in \mathcal D_{\initscale}$ and $\|v\|_2 = 1$, rule~(M4) and Step~2 of \cref{prop:network} give (including lag-bias directions via $s v_\brel \preccurlyeq s$)
\[
  \|\partial_v z_\theta\|_2 \le \sum_{k\ge 0} \|z_{k,1}\|_{\mathrm{abs}} \le \frac{\hat z(R,\initscale)}{\initscale} \le 26\mu^2\initscale,
  \qquad
  \|\partial_v^2 z_\theta\|_2 \le 2\sum_{k\ge 0} \|z_{k,2}\|_{\mathrm{abs}} \le \frac{2\hat z(R,\initscale)}{\initscale^2} \le 52\mu^2.
\]
Since $\|\hat y_\theta - y_n\|_2 \le \sqrt 2$ and the softmax Jacobian satisfies $\|\cJ(\hat y_\theta)\|_{\op} = 2\hat y_0\hat y_1 \le \frac{1}{2}$, we obtain by \cref{eq:scale-E}:
\[
  \lvert\langle\nabla_\theta\Lout, v\rangle\rvert
  = \bigl|\E\bigl[\langle\hat y_\theta - y_n, \partial_v z_\theta\rangle\bigr]\bigr|
  \le 26\sqrt 2\,\mu^2\initscale \le C_1\initscale,
  \qquad
  \Bigl|\partial_s^2\Lout(\theta+sv)\big|_{s=0}\Bigr|
  \le \frac{1}{2}(26\mu^2\initscale)^2 + 52\sqrt 2\,\mu^2 \le 2C_1.\qedhere
\]
\end{proof}

Subtracting $\KL(\bary \,\|\, \hat y_\theta) = \log\sum_a e^{(z_\theta)_a} - \log 2 - \langle\bary, z_\theta\rangle$ from $-\log\ip{\cS(z_\theta)}{y_n}$ and taking expectations gives, by \cref{thm:main}(ii), the exact KL-credit lower bound across $\mathcal D_{\initscale}$:
\begin{equation}
 \Lout(\theta)
 = \log 2 + \E\bigl[\KL(\bary \,\|\, \hat y_\theta)\bigr] - \outcredit(\theta)
 \ge \log 2 - C_2\initscale^2\mixrate^n.
 \label{eq:lo-identity}
\end{equation}

\par\bigskip\noindent\textbf{Theorem~\ref{thm:outcome-end} (Exponential plateau under outcome supervision; formal statement).}\enspace{\itshape
Under \cref{ass:init}, let $P(\tau)$ be a piecewise-constant orthogonal projection onto any subset of parameter blocks (including $\brel$), and evolve $\dot\theta(\tau) = -P(\tau)\nabla_\theta\Lout(\theta(\tau))$ (or projected GD $\theta_{m+1} = \theta_m - \stepsize P_m\nabla_\theta\Lout(\theta_m)$ for small $\stepsize > 0$). There exist constants $c, C > 0$ independent of $n$, $\initscale$, widths, and $P(\tau)$ such that for all $n \ge 2$, $\mixrate \in (0,1)$, $\initscale \le c\min\{\mixrate, \dmlp^{-1/2}\}$, and all $0 \le \tau \le c\,\mixrate^{-n}$ (or $m\stepsize \le c\,\mixrate^{-n}$): (i)~$\|\theta(\tau) - \theta_0\|_2 \le \ballradius\initscale$ (including any movement in $\brel$); (ii)~$\Lout(\theta(\tau)) \ge \log 2 - C\initscale^2\mixrate^n$; and (iii)~$\E[p_{\theta(\tau)}(s_n)] \le \frac{1}{2} + C\initscale\mixrate^{n/2}$.}\par\smallskip

\begin{proof}[Proof of \cref{thm:outcome-end}]
Set $R := \ballradius\initscale$, $\Delta := C_2\initscale^2\mixrate^n$, $c := \min\{c_2, (16H\mu^2)^{-1}, \frac{\ballradius^2}{2C_2}\}$, and $C := \max\{C_2, \sqrt{C_2/2}\}$. Since $\bar B(\theta_0, \frac{3}{2}R) \subset \mathcal D_{\initscale} \subset \mathcal U_{\initscale}$ and $\hat y_{\theta_0} \equiv \bary$ gives $\Lout(\theta_0) = \log 2$, the objective $\Lout$ is $C^\infty$ on $\mathcal U_{\initscale}$ (\cref{lem:frozen-relu-zeta}(a)) with bounded derivatives $G = C_1\initscale$ and $L = 2C_1$ (\cref{lem:end-smooth}), so concatenating constant-$P$ pieces yields a unique piecewise-$C^1$ trajectory up to $T_{\mathrm{exit}}$.

By \cref{eq:lo-identity}, $\Lout(\theta) \ge \Lout(\theta_0) - \Delta$ throughout $\mathcal D_{\initscale}$. Applying \cref{lem:energy-barrier} with $0 < \stepsize \le \min\{1, \ballradius\}/(2C_1)$ gives
\[
  T_{\mathrm{exit}} \ge \frac{R^2}{\Delta} = \frac{\ballradius^2}{C_2}\,\mixrate^{-n} \ge 2c\,\mixrate^{-n}
  \qquad\text{and}\qquad
  J\stepsize > \frac{R^2}{2\Delta} \ge c\,\mixrate^{-n},
\]
which proves~(i) and~(ii).

Finally, for $0 \le \tau \le c\mixrate^{-n}$, combining $\Lout(\theta(\tau)) \le \log 2$ with \cref{eq:lo-identity} gives $\E[\KL(\bary \,\|\, \hat y_{\theta(\tau)})] \le \Delta$. Applying Pinsker's inequality $|p_\theta(s_n) - \frac{1}{2}| \le \sqrt{\frac{1}{2}\KL(\bary \,\|\, \hat y_\theta)}$ and Jensen's inequality yields
\[
  \E\bigl[p_{\theta(\tau)}(s_n)\bigr] \le \frac{1}{2} + \sqrt{\frac{\Delta}{2}} = \frac{1}{2} + \sqrt{\frac{C_2}{2}}\,\initscale\mixrate^{n/2},
\]
proving~(iii).
\end{proof}

\begin{remark}[Cancellation of $\initscale^2$ and expected probability versus hard argmax]
\label{rem:outcome-scale}
In \cref{lem:energy-barrier}, lowering $\initscale$ scales both the maximum loss drop $\Delta = C_2\initscale^2\mixrate^n$ and the squared radius $R^2 = (\ballradius\initscale)^2$ by $\initscale^2$, so the ratio $R^2/\Delta = (\ballradius^2/C_2)\mixrate^{-n}$ does not depend on $\initscale$. Under unmasked step supervision (\cref{app:J-timescales}), by contrast, leaving $\initball$ requires growing $|a|$ from $\Theta(r_{\mathrm{align}}\initscale^4/\sqrt{d_k})$ against full-prefix dilution $1/n$ via $g \asymp \initscale d_k^{1/4}\sqrt{|a|}$, which takes $O(n\sqrt{d_k}\log(8/r_{\mathrm{align}})\,\initscale^{-5})$ time. Part~(iii) bounds expected softmax probability $\E[p_{\theta(\tau)}(s_n)]$; hard argmax accuracy $\Pr(\arg\max_a (z_\theta)_a = s_n)$ is discontinuous at $z_\theta = 0$, where $O(\initscale^2)$ logits aligned with $s_0$ attain the state-only baseline $\frac{1+\mixrate^n}{2}$ (\cref{app:exp-regime}).
\end{remark}

Finally, because the majorant bounds of \cref{app:network} depend on the readout dimension $K$ only through $\|V\|_{\op} \le M\initscale$, replacing scalar parity multiplication by isometric permutation action on the zero-sum subspace $\one_K^\perp$ extends \cref{thm:outcome-end} directly to any invertible $K$-state finite-state machine with contracting average transition operator.

\begin{corollary}[Invertible FSMs with contracting average operator]
\label{cor:reversible}
For any $K$-state set $\mathcal Q = \{0,\ldots,K-1\}$ ($K \ge 2$) and bijective token transitions $\{\tau_w\}_{w\in\Sigma}$ represented by permutation matrices $P_w \in \{0,1\}^{K\times K}$ (writing $e(s) \in \{0,1\}^K$ for the basis state vector and $y_n := e(s_n)$), let $s_t = \tau_{w_t}(s_{t-1})$ for $w_1,\ldots,w_n \sim \mu$ drawn i.i.d.\ conditioned on $s_0$ (hence under any law on $s_0$), with readout $V \in \R^{K\times\dmodel}$, $\bary := \frac{1}{K}\one_K$, average transition operator $\bar P := \E_{w\sim\mu}[P_w]$, and $\mixrate := \sigma_2(\bar P) = \|\bar P|_{\one_K^\perp}\|_{\op} \in (0,1)$. Then $\|\E[H^\top(y_n - \bary)]\|_2 \le \mixrate^{n-|S|}(\E\|H\|_{\op}^2)^{1/2}$ for every $\mathcal F_S$-measurable $H$, and on $\mathcal D_{\initscale}$ ($\initscale \le c\min\{\mixrate, \dmlp^{-1/2}\}$), $\|\nabla_\theta\outcredit(\theta)\|_2 \le C\initscale\mixrate^n$ with $c, C > 0$ independent of $K, n$, widths, and $P(\tau)$. From any balanced readout ($\hat y_{\theta_0} \equiv \bary$), projected gradient flow and projected GD on $\Lout$ satisfy $\|\theta(\tau)-\theta_0\|_2 \le \ballradius\initscale$, $\Lout(\theta(\tau)) \ge \log K - C\initscale^2\mixrate^n$, and $\E[\ip{\cS(z_{\theta(\tau)})}{y_n}] \le \frac{1}{K} + C\initscale\mixrate^{n/2}$ for $0 \le \tau \le c\mixrate^{-n}$.
\end{corollary}
\begin{proof}[Proof of \cref{cor:reversible}]
Every permutation matrix $P_w$ fixes $\one_K$ and maps $\one_K^\perp := \{u \in \R^K : \langle u, \one_K\rangle = 0\}$ isometrically onto itself, so $\bar P$ is block-diagonal on $\operatorname{span}(\one_K)\oplus\one_K^\perp$ with $\|\bar P|_{\one_K^\perp}\|_{\op} = \mixrate$. Writing $G := y_n - \bary = P_{w_n}\cdots P_{w_1}(y_0 - \bary)$ ($y_0 := e(s_0)$, $\|y_0 - \bary\|_2 = \sqrt{1 - 1/K}$) yields
\[
  \bigl\|\E[G \mid s_0, \mathcal F_S]\bigr\|_2 \le \mixrate^{n-|S|}
  \qquad\text{and}\qquad
  \bigl\|\E[H^\top G]\bigr\|_2 \le \mixrate^{n-|S|}\bigl(\E\|H\|_{\op}^2\bigr)^{1/2}.
\]
Because \cref{lem:frozen-relu-zeta,prop:network,prop:score-degree} use $V$ only through $\|V\|_{\op} \le M\initscale$ and $\|H\|_{\op} \le \|H\|_{\frob}$, the bound $\sup_{\mathcal D_{\initscale}}|\outcredit(\theta)| \le C_2\initscale^2\mixrate^n =: \Delta$ holds with $K$-independent constants. Applying \cref{lem:energy-barrier} to $\Lout(\theta) = \log K + \E[\KL(\bary \,\|\, \hat y_\theta)] - \outcredit(\theta) \ge \Lout(\theta_0) - \Delta$ ($C_1$ doubled via $\|\cJ(\hat y)\|_{\op} \le 1$) together with Pinsker's inequality $|\ip{\hat y_\theta}{y_n} - \frac{1}{K}| \le \sqrt{\frac{1}{2}\KL(\bary \,\|\, \hat y_\theta)}$ completes the proof.
\end{proof}

% ===================================================================
\section{Proof of \texorpdfstring{\Cref{thm:joint}}{Theorem 2}}
\label{app:joint}
\label{app:routing}

\paragraph{Proof roadmap and dependency chain for \cref{thm:joint,cor:matched}.}
Under the full causal mask $j \in \{0, \ldots, i_t\}$, population gradient flow on $\Lstep$ concentrates full-prefix softmax attention onto the target lag $\delta^\star = i_t - j_t = n$ and grows the bilinear state-to-letter gate from $\Theta(r_{\mathrm{align}}\initscale^4/\sqrt{d_k})$ to $\Omega(1)$ for every $p \in (0,1)$.\label{tab:J-phases} The proof has five parts:
\begin{enumerate}[leftmargin=*,nosep,topsep=3pt]
  \item \textbf{Exact symmetry reduction (\cref{app:J-reduce-sub}):} \Cref{lem:J-features,lem:J-symmetry,lem:J-reduce} show that unconstrained gradient flow on the nine Block-$2$ blocks $\theta_{\mathrm s}$ (\cref{def:gate}) preserves the joint parity-and-head swap manifold \cref{eq:U-manifold} and reduces exactly to \stepode{} (\cref{prop:reduced-transition}).
  \item \textbf{Timescale hierarchy (\cref{app:J-timescales}):} Identifies the three separated velocities governing prior warm-up ($\dot\cnt = \Theta(\initscale^2)$), additive lag routing ($\partial\brel/\partial|a| \asymp \sqrt{d_k}/\initscale^2 \gg 1$), and $|a|^{3/2}$ bilinear saddle escape.
  \item \textbf{Correlation cancellation and reciprocal weights (\cref{app:J-cancel-sub}):} \Cref{lem:U-corr} proves that calibrating $\sigma(-\cnt) \approx p$ cancels all $n-1$ distractor correlations to $O(p(1-p)/n)$, and \cref{lem:U-lag-cancel} derives the reciprocal-weight ODEs ($\omega_\delta := e^{-\brel_\delta}$) showing strict decay of $\omega_n + \omega_1$ and exact cancellation of non-letter drift in $\omega_\delta - \omega_1$.
  \item \textbf{Trajectory analysis across Phases 0--3 (\cref{app:J-geom-sub}):} \Cref{lem:U-mobility} establishes the conserved-Gram mobility and coupled $(g,a)$ integral envelope; \cref{lem:U-floor-vis} proves the unconditional $\frac{1}{7n}$ visible-distractor floor via a tilted Lyapunov invariant; \cref{lem:U-preroute} tracks prior warm-up ($\tau_0$) and full-prefix lag routing ($\tau_1$); and \cref{lem:U-postroute,lem:joint-explicit} establish permanent post-routing centering and worst-case margin mastery ($T_{\mathrm{step}}$).
  \item \textbf{Silent plateau routing, the clamped prior, and synthesis (\cref{app:J-silent-obstruction}):} \Cref{lem:silent-routing} bounds the $\widetilde O((1 + n^{3/2}d_k^{-1/2})\initscale^4)$ loss change across $[T_0, \tau_1]$, \cref{lem:clamped-shortcut} proves the clamped-prior reversal at $\cnt \equiv 0$, and combining these lemmas with \cref{cor:joint-outcome,thm:block-rms} completes \textbf{\cref{thm:joint}} and \textbf{\cref{cor:matched}}.
\end{enumerate}

\paragraph{Recurring quantities and prefix lag partition ($\mathcal W_t$ versus $\overline{\mathcal W}_t$).}
Throughout this section, $C > 0$ denotes a universal constant that may change from line to line; $O$, $\Theta$, $\Omega$, $\lesssim$ and $\asymp$ hide constants that depend only on $p$, and $\widetilde O$ additionally hides factors logarithmic in $n$, $1/\initscale$ and $1/r_{\mathrm{align}}$; and $c_{\mathrm F} := \frac{2}{7}$ denotes the distractor-floor constant of Hypothesis~(F) (\cref{app:hypF}).
\begin{center}
\small
\renewcommand{\arraystretch}{1.18}
\begin{tabular}{@{}ll@{\hspace{2.2em}}ll@{}}
\toprule
\textbf{Notation} & \textbf{Meaning} & \textbf{Notation} & \textbf{Meaning}\\
\midrule
$Z_t,\ N_t$ & partition sums on $\mathcal W_t$ / $\overline{\mathcal W}_t$ & $D_t$ & two-head gate $\sigma(u_t+a)-\sigma(u_t-a)$\\
$\pi_{t,r},\ \pi_t^\star := \pi_{t,t}$ & within-letter weights and lag-$n$ share & $\gamma_t := g D_t$ & effective gate strength at step $t$\\
$u_t$ & log-ratio $\log(Z_t/N_t)$ & $\mu_t$ & rule signal $\E[\sigma(-M_t)\chi_t\rho_t]$\\
$\rho_t$ & routed letter sign $\sum_r \pi_{t,r}\chi_r$ & $\xi_t$ & centering factor $\sigma(u_t+a)+\sigma(u_t-a)-1$\\
$M_t$ & true-class margin at step $t$ & $\omega_\delta := e^{-\brel_\delta}$ & reciprocal lag weight\\
\bottomrule
\end{tabular}
\end{center}
At step $t \in \{2,\ldots,n\}$, state $s_{t-1}$ sits at position $i_t = n + 1 + t$ and predicts $s_t = s_{t-1} + b_t \pmod 2$, while each context letter $w_r$ ($r \in [n]$) sits at position $j = 1 + r$, i.e., lookback lag $\delta = i_t - (1+r) = n + t - r$ behind $i_t$ (placing $w_t$ at constant lag $\delta^\star = n$). Partitioning $\{0, \ldots, i_t\}$ into the $n$ letter lags $\mathcal W_t := \{n+t-r : r \in [n]\}$ and the $t + 2$ non-letter lags $\overline{\mathcal W}_t := \{0, \ldots, i_t\} \setminus \mathcal W_t$ (illustrated below for $n = 4$) yields four structural properties for every $n \ge 2$:
\begin{center}
\small
\renewcommand{\arraystretch}{1.15}
\begin{tabular}{@{}llll@{}}
\toprule
Step $t$ ($n = 4$) & Query site $i_t$ & Letter lags $\mathcal W_t$ ($n$ lags) & Non-letter lags $\overline{\mathcal W}_t$ ($t+2$ lags)\\
\midrule
$t = 2$ & $i_2 = 7$ ($s_1$) & $\{2, 3, \mathbf{4}, 5\}$ & $\{0, \mathbf{1}, 6, 7\}$\\
$t = 3$ & $i_3 = 8$ ($s_2$) & $\{3, \mathbf{4}, 5, 6\}$ & $\{0, \mathbf{1}, 2, 7, 8\}$\\
$t = 4$ & $i_4 = 9$ ($s_3$) & $\{\mathbf{4}, 5, 6, 7\}$ & $\{0, \mathbf{1}, 2, 3, 8, 9\}$\\
\bottomrule
\end{tabular}
\end{center}
\begin{enumerate}[label=(\roman*),leftmargin=*,nosep,topsep=3pt]
\item Lag $\delta = n$ belongs to $\mathcal W_t$ at every step $t \in \{2,\ldots,n\}$ (reading $\chi_t$) and never belongs to $\overline{\mathcal W}_t$.
\item Lags $\delta \in \{0, 1\}$ belong to $\overline{\mathcal W}_t$ at every step $t \in \{2,\ldots,n\}$ and never belong to $\mathcal W_t$; since $\brel_0 \equiv 0$, $N_t := \sum_{\delta\in\overline{\mathcal W}_t} e^{\brel_\delta} \ge 1$.
\item At $t = 2$, $\overline{\mathcal W}_2 = \{0, 1, n+2, n+3\}$ contains exactly four lags for every horizon $n \ge 2$.
\item Every distractor $\delta \in \{2,\ldots,\min(n+3, 2n-1)\}\setminus\{n\}$ is visible at every step ($\delta \le i_2 \le i_t$) and enters $\mathcal W_t$ for at least one step (while for $n \in \{2,3\}$, any $\delta \in \{n+2, n+3\} \cap \{2n, 2n+1\}$ belongs to $\overline{\mathcal W}_t$ at every step, sharing the membership of $\delta = 1$). For $n \ge 5$, every $\delta \in \{n+4,\ldots,2n-1\}$ is invisible for $t < \delta-n-1$, in $\overline{\mathcal W}_t$ only at $t \in \{\delta-n-1, \delta-n\}$, and in $\mathcal W_t$ for all $t \in \{\delta-n+1,\ldots,n\}$ (including $t = n$), while $\delta \in \{2n, 2n+1\}$ never enters $\mathcal W_t$.
\end{enumerate}

%-----------------------------------------------------------------------------
\subsection{Symmetry Reduction}
\label{app:J-reduce-sub}

\Cref{def:gate} specifies the nine-block Block-$2$ sub-circuit for joint full-prefix lag routing and query--key amplification. By parity-flip and head-swap invariance (\cref{lem:J-symmetry}), unconstrained nine-block gradient flow stays on a four-coordinate manifold and reduces to \stepode{} (\cref{prop:reduced-transition,lem:J-reduce}).

\begin{definition}[Structured Block-$2$ sub-circuit for step learning]
\label{def:gate}
Within \cref{ass:arch}, given \emph{frozen precomputed match signs} $\chi_r = (-1)^{b_r}$ at $j = 1+r$ ($r \in [n]$), evolve the nine trainable blocks $\theta_{\mathrm s} := (U, w_1, w_2, Q_1, Q_2, k_1, k_2, \brel^{(1)}, \brel^{(2)})$ over continuous training time $\tau \ge 0$ (denoted with overdots $\dot f := \ddtau{f}$) by unit-rate population gradient flow $\dot\theta_{\mathrm s} = -\nabla_{\theta_{\mathrm s}}\Lstep(\theta_{\mathrm s})$ ($t \ge 2$) over the full causal prefix $j \in \{0, \ldots, i_t\}$ without a target-lag mask:
\begin{enumerate}[label=(\roman*),leftmargin=*,nosep,topsep=2pt]
\item \emph{Frozen features and inactive blocks:} using five orthonormal vectors $(u_0, b_0, \ell_0, g_1, g_2)$ orthogonal to the $2$D state slice (with $x_{t-1} := (-1)^{s_{t-1}}$, $R = \bigl(\begin{smallmatrix}0&1\\1&0\end{smallmatrix}\bigr)$, and $e_\pm = \frac{1}{\sqrt 2}(e_0 \pm e_1)$), supply
\[
  h^{(0)}_{i_t} = \frac{\initscale}{\sqrt 2}\bigl(u_0 + x_{t-1}e_-\bigr), \qquad
  h^{(0)}_{1+r} = \initscale\bigl(b_0 + \chi_r\ell_0\bigr) \quad (r \in [n]),
\]
and $0$ on $(b_0, \ell_0)$ at all non-letter sites; freeze $W_V^{(2,h)} = v\,g_h \ell_0^\top$ ($v \in [\initscale/2, 2\initscale]$), and zero Block~$1$, Head~$3$, and MLP weights.
\item \emph{Jointly trained blocks:} leave Block-$2$ Heads~$1,2$ queries $Q_h \in \R^{d_k\times 2}$ on the state slice, keys $k_h \in \R^{d_k}$ on $b_0$, lag biases $\brel^{(h)} \in \R^{2n+2}$ ($\brel^{(h)}_0 = 0$), and readout $V = (U, w_1, w_2)$ ($U \in \R^{2\times 2}$, $w_h \in \R^2$) unconstrained.
\item \emph{Paired initialization:} set $U(0) = 0$, $w_h(0) = 0$, $\brel^{(h)}(0) = 0$, $Q_h(0) = (-1)^{h-1}q(0)e_-^\top$, and $k_h(0) = k(0)$ with $\|q(0)\|_2, \|k(0)\|_2, v \in [\initscale/2, 2\initscale]$ and $\lvert\langle q(0), k(0)\rangle\rvert \ge \frac{1}{2}r_{\mathrm{align}}\,\|q(0)\|_2\|k(0)\|_2$ ($r_{\mathrm{align}} \in (0,1]$).
\end{enumerate}
\end{definition}

On the paired two-head subspace \cref{eq:U-manifold}, full-prefix softmax attention factors into a scalar letter-versus-nonletter gate and a shared within-letter distribution.

\begin{lemma}[Softmax factorization and step margin]
\label{lem:J-features}
\label{lem:U-features}
On the subspace
\begin{equation}
 U = \bar u(I - R), \quad w_1 = -w_2 = \bar w(e_0 - e_1), \quad Q_1 = -Q_2 = q e_-^\top, \quad k_1 = k_2 = k, \quad \brel^{(1)} = \brel^{(2)} = \brel,
 \label{eq:U-manifold}
\end{equation}
define $\cnt := 2\initscale\bar u$, $g := 2\initscale v\bar w$, $a := \frac{\initscale^2\langle q, k\rangle}{\sqrt{2d_k}}$, $Z_t := \sum_{\delta\in\mathcal W_t} e^{\brel_\delta}$, $N_t := \sum_{\delta\in\overline{\mathcal W}_t} e^{\brel_\delta}$, $\pi_{t,r} := e^{\brel_{n+t-r}}/Z_t$, and $u_t := \log(Z_t/N_t)$. Then at each step $t \in \{2,\ldots,n\}$, Head $h \in \{1,2\}$ assigns full-prefix softmax weight
\begin{equation}
 \alpha^{(h)}_{t,\,1+r} = \sigma\bigl(u_t + (-1)^{h-1}x_{t-1}a\bigr)\,\pi_{t,r} \qquad (r \in [n]),
 \label{eq:U-factor}
\end{equation}
and with $D_t := \sigma(u_t + a) - \sigma(u_t - a)$ and $\rho_t := \sum_{r=1}^n \pi_{t,r}\chi_r$, the true-class margin is
\begin{equation}
 M_t := \ip{z_t}{y_t - (\one - y_t)} = \chi_t\bigl(\cnt + g\,D_t\,\rho_t\bigr).
 \label{eq:J-logits}
\end{equation}
\end{lemma}
\begin{proof}[Proof of \cref{lem:J-features}]
Since $h^{(1)}_j = h^{(0)}_j$, Head~$h$ has query $\frac{\initscale}{\sqrt 2}(-1)^{h-1}x_{t-1}q$ at $i_t$, key $\initscale k$ at letters $j = 1+r$ ($r \in [n]$), and key $0$ on $\overline{\mathcal W}_t$. Thus the unnormalized score at lag $\delta = i_t - j$ equals $\brel_\delta + (-1)^{h-1}x_{t-1}a$ on $\mathcal W_t$ and $\brel_\delta$ on $\overline{\mathcal W}_t$, giving partition sum $e^{(-1)^{h-1}x_{t-1}a}Z_t + N_t$ and proving \cref{eq:U-factor}.

Next, $W_V^{(2,h)}h^{(1)}_{1+r} = \initscale v\chi_r g_h$ on $\mathcal W_t$ and $0$ on $\overline{\mathcal W}_t$. Using $\sigma(u_t + x_{t-1}a) - \sigma(u_t - x_{t-1}a) = x_{t-1}D_t$ for $x_{t-1} \in \{\pm 1\}$ and $w_1 = -w_2 = \bar w(e_0 - e_1)$, the two heads contribute $\frac{1}{2}g x_{t-1}D_t\rho_t(e_0 - e_1)$. Adding the prior term $U(\frac{\initscale}{\sqrt 2}x_{t-1}e_-) = \frac{1}{2}\cnt x_{t-1}(e_0 - e_1)$ gives
\[
  (z_t)_0 - (z_t)_1 = x_{t-1}\bigl(\cnt + g\,D_t\,\rho_t\bigr),
\]
and multiplying by $(-1)^{s_t} = x_{t-1}\chi_t$ yields \cref{eq:J-logits}.
\end{proof}

Before proving invariance of \cref{eq:U-manifold}, we record the algebraic bounds and lag-bias derivatives of the two-head gate $D_t$.

\begin{lemma}[Two-head gate and prefix softmax calculus]
\label{lem:gate-calculus}
For any $(u_t, a, \brel)$ with $a > 0$, $D_t = \sigma(u_t + a) - \sigma(u_t - a)$, $\xi_t = \sigma(u_t + a) + \sigma(u_t - a) - 1$, and $\rho_t = \sum_{r=1}^n \pi_{t,r}\chi_r$ ($\chi_r \in \{\pm 1\}$):
\begin{enumerate}[label=(\roman*),leftmargin=*,itemsep=1pt,topsep=2pt]
\item \emph{(Algebraic and tail bounds.)} $\operatorname{sign}\xi_t = \operatorname{sign}u_t = \operatorname{sign}(\frac{1}{N_t} - \frac{1}{Z_t})$, $1 - D_t, |\xi_t| \le 2e^{-(a - |u_t|)}$,
\begin{equation}
  D_t = \frac{2(\sinh a)\,Z_t N_t}{Z_t^2 + N_t^2 + 2(\cosh a)Z_t N_t}, \qquad
  \tanh a \le \frac{D_t}{\sigma'(u_t+a)+\sigma'(u_t-a)} \le \sinh a.
  \label{eq:U-rational}
\end{equation}
\item \emph{(Lag-bias derivatives and uniform $\frac{1}{2}$-Lipschitzness.)} $\pd{D_t}{u_t} = -D_t\xi_t$, and for every $\delta \ge 1$,
\[
  \pd{(D_t\rho_t)}{\brel_\delta} = \begin{cases}
    D_t\pi_{t,r}\bigl(\chi_r - (1 + \xi_t)\rho_t\bigr), & \delta \in \mathcal W_t \;\; (r = n + t - \delta),\\[2pt]
    -D_t\xi_t\,\dfrac{e^{\brel_\delta}}{N_t}\,\rho_t, & \delta \in \overline{\mathcal W}_t,
  \end{cases}
  \qquad
  \Bigl\lvert\pd{(D_t\rho_t)}{\brel_\delta}\Bigr\rvert \le \frac{1}{2}.
\]
\end{enumerate}
\end{lemma}
\begin{proof}[Proof of \cref{lem:gate-calculus}]
\proofstep{Part (i): Rational form and ratio bounds.}
Substituting $\sigma(u_t \pm a) = \frac{e^{u_t\pm a}}{1 + e^{u_t\pm a}}$ with $e^{u_t} = Z_t/N_t$ yields the rational formula for $D_t$, the sign of $\xi_t$, and the tail bound $1 - D_t, |\xi_t| \le 1 - \sigma(a - |u_t|) + \sigma(-a + |u_t|) < 2e^{-(a - |u_t|)}$. Moreover, $x \mapsto \frac{x + \cosh a}{x\cosh a + 1}$ is decreasing on $x = \cosh u_t \ge 1$, giving
\[
  \tanh a \le \frac{\sinh a(\cosh u_t + \cosh a)}{\cosh u_t\cosh a + 1} = \frac{D_t}{\sigma'(u_t+a)+\sigma'(u_t-a)} \le \sinh a.
\]

\proofstep{Part (ii): Lag-bias derivatives and $\frac{1}{2}$-Lipschitz bound.}
Combining $\sigma'(u_t + a) - \sigma'(u_t - a) = -D_t\xi_t$, $\pd{u_t}{\brel_\delta} = \pi_{t,r}$ ($\delta \in \mathcal W_t$) or $-e^{\brel_\delta}/N_t$ ($\delta \in \overline{\mathcal W}_t$), and $\pd{\rho_t}{\brel_\delta} = \pi_{t,r}(\chi_r - \rho_t)$ ($\delta \in \mathcal W_t$) gives the piecewise derivative formulas.

To verify the uniform $\frac{1}{2}$ bound, write $D := D_t$ and $S := 1 + \xi_t \in [0, 2]$, so $0 \le D \le \min\{S, 2 - S\}$ and $D|S - 1| \le \frac{1}{4}$. On $\overline{\mathcal W}_t$, this immediately gives $|\pd{(D_t\rho_t)}{\brel_\delta}| \le \frac{1}{4}$. On $\mathcal W_t$, write $\rho_t = \pi\chi_r + (1 - \pi)z$ with $\pi := \pi_{t,r} \in [0,1]$ and $|z| \le 1$; by convexity in $\zeta := \chi_r z \in [-1, 1]$, $D\pi|\chi_r - S\rho_t|$ is maximized either at $\zeta = +1$ (where it is at most $D|1 - S| \le \frac{1}{4}$) or at $\zeta = -1$, where
\[
  D\pi\lvert 1 + S - 2S\pi\rvert
  \le \max\Bigl\{\frac{\min\{S, 2-S\}}{S}\frac{(1+S)^2}{8},\; (2 - S)(S - 1)\Bigr\}
  \le \frac{1}{2},
\]
using $\frac{2-S}{S}\frac{(1+S)^2}{8} = \frac{1}{2} - \frac{(S-1)(S^2+S+2)}{8S} \le \frac{1}{2}$ for $S \in [1, 2]$.
\end{proof}

Next, we verify that a joint state-flip and head-swap involution preserves $\Lstep$ and traps unconstrained gradient flow on \cref{eq:U-manifold} for all $\tau \ge 0$.

\begin{lemma}[Joint parity-and-head swap symmetry]
\label{lem:J-symmetry}
\label{lem:U-symmetry}
\label{rem:J-line}
Define the orthogonal linear involution
\[
 \mathcal T\bigl(U, w_1, w_2, Q_1, Q_2, k_1, k_2, \brel^{(1)}, \brel^{(2)}\bigr) := \bigl(RUR,\; Rw_2,\; Rw_1,\; Q_2 R,\; Q_1 R,\; k_2,\; k_1,\; \brel^{(2)},\; \brel^{(1)}\bigr).
\]
Under $s_0 \sim \mathrm{Unif}(\{0,1\})$, $\Lstep(\mathcal T\theta_{\mathrm s}) = \Lstep(\theta_{\mathrm s})$ for all $\theta_{\mathrm s}$, and gradient flow from \cref{def:gate}(iii) preserves \cref{eq:U-manifold} for all $\tau \ge 0$.
\end{lemma}
\begin{proof}[Proof of \cref{lem:J-symmetry}]
Holding $(w_1,\ldots,w_n)$ fixed and flipping $s_0 \mapsto 1 - s_0$ negates $x_{t-1} \mapsto -x_{t-1}$ and maps $y_t \mapsto R y_t$ while leaving every $\chi_r$ fixed. Since $R e_- = -e_-$, replacing $(Q_1, Q_2)$ by $(Q_2 R, Q_1 R)$ and negating $x_{t-1}$ swaps the queries of Head~$1$ and Head~$2$; swapping $(k_1, \brel^{(1)}) \leftrightarrow (k_2, \brel^{(2)})$ swaps their scalar outputs, and $(RUR, Rw_2, Rw_1)$ maps $z_t \mapsto R z_t$. Because $-\log\ip{\cS(R z_t)}{R y_t} = -\log\ip{\cS(z_t)}{y_t}$ and $s_0 \sim \mathrm{Unif}(\{0,1\})$, we obtain $\Lstep(\mathcal T\theta_{\mathrm s}) = \Lstep(\theta_{\mathrm s})$.

Orthogonality of $\mathcal T$ ensures that $\nabla\Lstep$ preserves the fixed-point subspace $\{U = RUR,\; w_1 = Rw_2,\; Q_1 = Q_2 R,\; k_1 = k_2,\; \brel^{(1)} = \brel^{(2)}\}$. Together with the nullspace identities
\[
  (\nabla_{Q_h}\Lstep)e_+ = 0, \qquad
  (\nabla_U\Lstep)e_+ = 0, \qquad
  \one^\top\nabla_U\Lstep = 0, \qquad
  \one^\top\nabla_{w_h}\Lstep = 0,
\]
the initialization \cref{def:gate}(iii) and ODE uniqueness preserve \cref{eq:U-manifold} for all $\tau \ge 0$.
\end{proof}

Projecting the nine-block gradient flow onto the invariant manifold \cref{eq:U-manifold} yields the exact reduced ODE system (\cref{prop:reduced-transition}), whose conserved Gram-difference nuclear norm $\Delta_0$ and explicit coordinate partial derivatives are recorded in \cref{lem:J-reduce}.

\begin{proposition}[Symmetry-reduced step gradient flow]
\label{prop:reduced-transition}
Under \cref{def:gate}, let $\Delta_0 := \|q(0)q(0)^\top - k(0)k(0)^\top\|_*$ (nuclear norm, i.e., sum of singular values) and $\rho_t := \sum_{r=1}^n \pi_{t,r}\chi_r$, and define
\[
  D_t := \sigma(u_t + a) - \sigma(u_t - a), \qquad
  M_t := \chi_t(\cnt + g\,D_t\,\rho_t),
\]
so $\Lstep(\cnt,g,a,\brel) = \frac{1}{n-1}\sum_{t=2}^n \E[\log(1 + e^{-M_t})]$. Under $\ddtau{\theta_{\mathrm s}} = -\nabla_{\theta_{\mathrm s}}\Lstep$, $(\cnt, g, a, \brel)$ obey the exact system
\begin{equation*}
  \ddtau{\cnt} = -\initscale^2\pd{\Lstep}{\cnt}, \qquad
  \ddtau{g} = -\initscale^2 v^2\pd{\Lstep}{g}, \qquad
  \ddtau{a} = -\initscale^2\kappa(a)\pd{\Lstep}{a}, \qquad
  \ddtau{\brel_\delta} = -\frac{1}{2}\pd{\Lstep}{\brel_\delta} \;\; (\delta \ge 1),
  \tag{\textsf{Step-ODE}}
  \label{eq:reduced-transition}
\end{equation*}
with mobility $\kappa(a) := \frac{1}{4d_k}\sqrt{\initscale^4\Delta_0^2 + 8d_k a^2}$, $\cnt(0) = g(0) = 0$, $\brel(0) = 0$, and $a(0) = \initscale^2\langle q(0), k(0)\rangle/\sqrt{2d_k}$.
\end{proposition}

\begin{lemma}[Conserved Gram matrix and explicit ODE]
\label{lem:J-reduce}
\label{lem:U-reduce}
On \cref{eq:U-manifold}, $q(\tau)q(\tau)^\top - k(\tau)k(\tau)^\top$ is constant in $\tau$ with nuclear norm $\Delta_0 := \|q(0)q(0)^\top - k(0)k(0)^\top\|_*$, and $(\cnt, g, a, \brel)$ satisfy \stepode{}:
\begin{equation}
 \dot\cnt = -\initscale^2\pd{\Lstep}{\cnt}, \qquad
 \dot g = -\initscale^2 v^2\pd{\Lstep}{g}, \qquad
 \dot a = -\initscale^2\kappa(a)\pd{\Lstep}{a}, \qquad
 \dot\brel_\delta = -\frac{1}{2}\pd{\Lstep}{\brel_\delta} \;\; (\delta \ge 1),
 \label{eq:J-flow}
\end{equation}
where $\kappa(a) = \frac{1}{4d_k}\sqrt{\initscale^4\Delta_0^2 + 8d_k a^2}$. Setting $\gamma_t := g D_t$, $\mu_t := \E[\sigma(-M_t)\chi_t\rho_t]$, and $\xi_t := \sigma(u_t+a)+\sigma(u_t-a)-1 \in (-1,1)$, the partial derivatives of $\Lstep$ are given by
\begin{subequations}
\label{eq:U-explicit-grads}
\begin{align}
 -\pd{\Lstep}{\cnt} &= \frac{1}{n-1}\sum_{t=2}^n \E\bigl[\sigma(-M_t)\chi_t\bigr], \qquad\qquad
 -\pd{\Lstep}{g} = \frac{1}{n-1}\sum_{t=2}^n D_t\,\mu_t,\\
 -\pd{\Lstep}{a} &= \frac{1}{n-1}\sum_{t=2}^n g\bigl(\sigma'(u_t+a)+\sigma'(u_t-a)\bigr)\mu_t,\\
 -\pd{\Lstep}{\brel_\delta} &= \frac{1}{n-1}\sum_{\substack{t\,:\,\delta\in\mathcal W_t\\(r = n+t-\delta)}} \gamma_t\pi_{t,r}\Bigl(\E\bigl[\sigma(-M_t)\chi_t\chi_r\bigr] - (1+\xi_t)\mu_t\Bigr) + \frac{1}{n-1}\sum_{t\,:\,\delta\in\overline{\mathcal W}_t} \gamma_t\,\frac{e^{\brel_\delta}}{N_t}\,\xi_t\,\mu_t.
 \label{eq:U-beta-explicit}
\end{align}
\end{subequations}
\end{lemma}
\begin{proof}[Proof of \cref{lem:J-reduce} and \cref{prop:reduced-transition}]
By \cref{lem:J-features,lem:J-symmetry}, $\theta_{\mathrm s}(\tau)$ remains on \cref{eq:U-manifold}, where $\|U\|_{\frob}^2 = 4\bar u^2$, $\|w_1\|_2^2 + \|w_2\|_2^2 = 4\bar w^2$, $\|Q_1\|_{\frob}^2 + \|Q_2\|_{\frob}^2 = 2\|q\|_2^2$, $\|k_1\|_2^2 + \|k_2\|_2^2 = 2\|k\|_2^2$, and $\sum_{h=1}^2 \|\brel^{(h)}\|_2^2 = 2\|\brel\|_2^2$. Unit-rate gradient flow in $(\bar u, \bar w, q, k, \brel)$ therefore reads
\[
 \dot{\bar u} = -\frac{1}{4}\pd{\Lstep}{\bar u}, \qquad
 \dot{\bar w} = -\frac{1}{4}\pd{\Lstep}{\bar w}, \qquad
 \dot q = -\frac{\initscale^2}{2\sqrt{2d_k}}\pd{\Lstep}{a}\,k, \qquad
 \dot k = -\frac{\initscale^2}{2\sqrt{2d_k}}\pd{\Lstep}{a}\,q, \qquad
 \dot\brel_\delta = -\frac{1}{2}\pd{\Lstep}{\brel_\delta}.
\]
Rescaling $\cnt = 2\initscale\bar u$ and $g = 2\initscale v\bar w$ gives $\dot\cnt$ and $\dot g$. From $\dot q$ and $\dot k$ we obtain $\ddtau{(qq^\top - kk^\top)} = 0$, conserving $\Delta_0 = \|qq^\top - kk^\top\|_*$, and $(\|q\|_2^2 + \|k\|_2^2)^2 = \Delta_0^2 + 4\langle q, k\rangle^2 = \Delta_0^2 + \frac{8d_k}{\initscale^4}a^2$ yields
\[
 \dot a = \frac{\initscale^2}{\sqrt{2d_k}}\bigl(\langle\dot q, k\rangle + \langle q, \dot k\rangle\bigr) = -\frac{\initscale^4}{4d_k}\bigl(\|q\|_2^2 + \|k\|_2^2\bigr)\pd{\Lstep}{a} = -\initscale^2\kappa(a)\pd{\Lstep}{a}.
\]
Differentiating $\log(1 + e^{-M_t})$ via $-\frac{\partial}{\partial M}\log(1 + e^{-M}) = \sigma(-M)$ gives the $(\cnt, g, a)$ partial derivatives directly, and \cref{lem:gate-calculus}(ii) yields \cref{eq:U-beta-explicit}.
\end{proof}

%-----------------------------------------------------------------------------
\subsection{Timescale Separation}
\label{app:J-timescales}
\label{tab:J-speeds}

Introducing the dimensionless prior time $\tilde\tau := \initscale^2\tau$ and unit-scale value factor $v_0 := v/\initscale \in [\frac{1}{2}, 2]$ absorbs the ubiquitous $\initscale^2$ prefactor in \stepode{}, yielding
\[
  \frac{\rmd\cnt}{\rmd\tilde\tau} = -\pd{\Lstep}{\cnt}, \qquad
  \frac{\rmd g}{\rmd\tilde\tau} = -\initscale^2 v_0^2\pd{\Lstep}{g}, \qquad
  \frac{\rmd a}{\rmd\tilde\tau} = -\kappa(a)\pd{\Lstep}{a}, \qquad
  \frac{\rmd\brel_\delta}{\rmd\tilde\tau} = -\frac{1}{2\initscale^2}\pd{\Lstep}{\brel_\delta} \;\; (\delta \ge 1).
\]
This rescaling makes the three-stage hierarchy transparent:
\begin{enumerate}[leftmargin=*,itemsep=2pt,topsep=3pt]
\item \emph{Prior calibration ($\tilde\tau = O(\log n)$):} $\frac{\rmd\cnt}{\rmd\tilde\tau} = \Theta(1)$ drives $\cnt \to c^\star$ in $O(\log n)$ prior time ($O(\initscale^{-2}\log n)$ physical time) before $g$, $a$, or $\brel$ move appreciably.
\item \emph{Silent lag routing:} Once $\cnt \approx c^\star$ ($\brel \approx 0$, $|a| \asymp r_{\mathrm{align}}\initscale^4/\sqrt{d_k}$), full-prefix dilution sets $\pi_t^\star \approx 1/n$ and $D_t \asymp a$. Because $\frac{\rmd\brel}{\rmd\tilde\tau}$ carries a $\frac{1}{2\initscale^2}$ multiplier whereas $\frac{\rmd a}{\rmd\tilde\tau}$ is throttled by $\kappa(a)$, their relative velocity satisfies
\[
  \frac{\rmd\brel}{\rmd|a|} \asymp \frac{\sqrt{d_k}}{\initscale^2} \gg 1,
\]
concentrating $\pi_t^\star$ onto $\delta^\star = n$ while $|a| \ll 1$ and $\Lstep \approx H(p)$.
\item \emph{Saddle escape:} Since $\kappa(a) \ge \max\{\frac{\initscale^4}{8\sqrt 2\,d_k}, \frac{a}{\sqrt{2d_k}}\}$, $|a|$ crosses $a_\star := \initscale^4/\sqrt{d_k}$ and then obeys $\frac{\rmd|a|}{\rmd\tilde\tau} \gtrsim \frac{\initscale}{n d_k^{1/4}}|a|^{3/2}$ with $g \asymp \initscale d_k^{1/4}\sqrt{|a|}$; integrating $\int_{a_\star}^\infty a^{-3/2}\,\rmd a = 2d_k^{1/4}/\initscale^2$ gives total physical time $O(n\sqrt{d_k}\log(8/r_{\mathrm{align}})\,\initscale^{-5})$.
\end{enumerate}

%-----------------------------------------------------------------------------
\subsection{Residual Correlation Bounds}
\label{app:J-cancel-sub}

By the reflection symmetry $(g, a) \mapsto (-g, -a)$ of \stepode{}, assume without loss of generality that $a(0) > 0$, so
\begin{equation}
 a_0 := a(0) = \frac{\initscale^2\langle q(0), k(0)\rangle}{\sqrt{2d_k}} \in \Biggl[\frac{r_{\mathrm{align}}\initscale^4}{8\sqrt{2d_k}},\; \frac{2\sqrt 2\,\initscale^4}{\sqrt{d_k}}\Biggr].
 \label{eq:U-a0}
\end{equation}
Set $c^\star := \log\frac{1-p}{p}$ ($c^\star = 0$ at $p = 1/2$). Since $\E[\sigma(-\chi_t\cnt)\chi_t] = (1-p)\sigma(-\cnt) - p\sigma(\cnt) = \sigma(-\cnt) - p$, the direct prior term is calibrated to the one-step transition prior $q(s_{t-1})$ precisely when $\sigma(-\cnt) = p$.

\begin{lemma}[Correlation cancellation at prior quasi-equilibrium]
\label{lem:U-corr}
Suppose $a \in (0, 1/4]$, $g \ge 0$, $|\sigma(-\cnt) - p| \le \frac{p(1-p)}{1024n}$, and $\max_{2\le t\le n}\gamma_t \le \frac{p(1-p)}{1024n}$. Then for every step $t \in \{2,\ldots,n\}$:
\begin{enumerate}[label=(\roman*),leftmargin=*]
\item $\E[\sigma(-M_t)\chi_t^2] = \E[\sigma(-M_t)] \in \bigl[\frac{3}{2}p(1-p),\; \frac{5}{2}p(1-p)\bigr]$.
\item For every distractor $r \in [n]\setminus\{t\}$, $|\E[\sigma(-M_t)\chi_t\chi_r]| \le |\sigma(-\cnt) - p| + \frac{1}{2}\gamma_t < \frac{p(1-p)}{512n}$.
\item $|\sum_{r\ne t}\pi_{t,r}\E[\sigma(-M_t)\chi_t\chi_r]| \le (1 - \pi_t^\star)\frac{p(1-p)}{512n}$, and if $\pi_t^\star \ge \frac{1}{8n}$, then $\mu_t \in \bigl[p(1-p)\pi_t^\star,\; 3p(1-p)\pi_t^\star\bigr]$.
\end{enumerate}
\end{lemma}
\begin{proof}[Proof of \cref{lem:U-corr}]
Write $\epsilon := \sigma(-\cnt) - p$ and $G := \frac{p(1-p)}{1024n}$ ($|\epsilon| \le G$, $\max_t\gamma_t \le G$).

\proofstep{Part (i): True-letter expectation.}
Since $\E[\sigma(-\chi_t\cnt)] = 2p(1-p) + (1-2p)\epsilon$ and $|\sigma(-M_t) - \sigma(-\chi_t\cnt)| \le \frac{1}{4}\gamma_t$ (by $\frac{1}{4}$-Lipschitzness of $\sigma$ and $|\rho_t| \le 1$), we obtain
\[
  \bigl|\E[\sigma(-M_t)] - 2p(1-p)\bigr| \le |\epsilon| + \frac{1}{4}\gamma_t \le \frac{5}{4}G < \frac{1}{2}p(1-p),
\]
proving~(i).

\proofstep{Parts (ii)--(iii): Distractor cancellation and $\mu_t$ sandwich.}
For $r \ne t$, split $M_t = M_t^{(-r)} + \gamma_t\pi_{t,r}\chi_t\chi_r$, where $M_t^{(-r)}$ omits index $r$ and is independent of $\chi_r$. Using $|\sigma(-M_t) - \sigma(-M_t^{(-r)})| \le \frac{1}{4}\gamma_t\pi_{t,r} \le \frac{1}{4}\gamma_t$, $|\E[\sigma(-M_t^{(-r)})\chi_t] - \epsilon| \le \frac{1}{4}\gamma_t$, and $|\E[\chi_r]| = |1-2p| \le 1$ yields
\[
  \bigl|\E[\sigma(-M_t)\chi_t\chi_r] - (1-2p)\epsilon\bigr| \le \frac{1}{2}\gamma_t
  \quad\implies\quad
  \bigl|\E[\sigma(-M_t)\chi_t\chi_r]\bigr| \le |\epsilon| + \frac{1}{2}\gamma_t \le \frac{3}{2}G < \frac{p(1-p)}{512n},
\]
proving~(ii). Summing over $r \ne t$ with $\sum_{r\ne t}\pi_{t,r} = 1 - \pi_t^\star$ and adding $\pi_t^\star\E[\sigma(-M_t)]$ proves~(iii).
\end{proof}

\begin{remark}[What clamping the transition prior changes]
\label{rem:shortcut-paradox}
The distractor estimate in the proof of \cref{lem:U-corr}(ii) uses none of the lemma's hypotheses: for every $(\cnt, g, a, \brel)$ and $r \ne t$, $|\E[\sigma(-M_t)\chi_t\chi_r] - (1-2p)(\sigma(-\cnt) - p)| \le \frac{1}{2}|\gamma_t|$. If the prior is clamped ($\cnt \equiv 0$, $\sigma(-\cnt) = \frac{1}{2}$) or miscalibrated by $|\sigma(-\cnt) - p| = \delta_0 > 0$, every distractor retains pull $\mixrate\delta_0$ ($\frac{1}{2}\mixrate^2$ at $\cnt \equiv 0$) up to $\frac{1}{2}|\gamma_t|$, shifting $\mu_t$ across non-letter lags via $\xi_t\mu_t$ (\cref{lem:clamped-shortcut}). This changes the \emph{direction} of early lag drift; integrating \stepode{} with $\cnt \equiv 0$, the clamped isolated circuit still routes and solves because the $1/2$ residual carries a larger true-lag signal when upstream match signs $\chi_r$ are fixed (\cref{sec:what-calibration-buys}).
\end{remark}

We analyze full-prefix lag routing in the reciprocal lag weights $\omega_\delta := e^{-\brel_\delta}$ rather than ratios of growing exponentials. \Cref{lem:U-lag-cancel} splits $-\dot\omega_\delta$ into within-letter contrasts $P_\delta$ and non-letter gate remainders $R_\delta$, and records two exact cancellations against the non-letter state lag $\delta = 1$.

\begin{lemma}[Reciprocal-weight decomposition and cancellation identities]
\label{lem:U-lag-cancel}
Set $\omega_\delta := e^{-\brel_\delta}$ for $\delta \in \{0,\ldots,2n+1\}$ ($\omega_0 \equiv 1$) and $E_{t,r} := \E[\sigma(-M_t)\chi_t\chi_r]$. Along \cref{eq:U-manifold}, multiplying $-\frac{1}{2}\pd{\Lstep}{\brel_\delta}$ by $\omega_\delta$ decomposes $-\dot\omega_\delta = P_\delta + R_\delta$ ($\delta \ge 1$) into letter contrasts and non-letter gate remainders:
\begin{equation}
 -\dot\omega_\delta = \underbrace{\frac{1}{2(n-1)}\sum_{\substack{t\,:\,\delta\in\mathcal W_t\\(r = n+t-\delta)}}\frac{\gamma_t}{Z_t}\bigl(E_{t,r} - \mu_t\bigr)}_{=: P_\delta}
 \;+\;
 \underbrace{\frac{1}{2(n-1)}\sum_{t\,:\,\delta\in\overline{\mathcal W}_t}\frac{\gamma_t\,\xi_t\,\mu_t}{N_t} - \frac{1}{2(n-1)}\sum_{t\,:\,\delta\in\mathcal W_t}\frac{\gamma_t\,\xi_t\,\mu_t}{Z_t}}_{=: R_\delta}.
 \label{eq:U-wd-ode}
\end{equation}
In particular, whenever $a > 0$, $g \ge 0$, and $\mu_t \ge 0$ for all $t \in \{2,\ldots,n\}$ (without any calibration assumption),
\begin{equation}
 -\ddtau{(\omega_n + \omega_1)} = P_n + \frac{1}{2(n-1)}\sum_{t=2}^n \gamma_t\,\xi_t\,\mu_t\Bigl(\frac{1}{N_t} - \frac{1}{Z_t}\Bigr) \ge P_n \ge 0.
 \label{eq:U-post-wnw1}
\end{equation}
Moreover, under the hypotheses of \cref{lem:U-corr} with $\min_{2\le t\le n}\pi_t^\star \ge \frac{1}{8n}$ and $a \in (0, 1/4]$:
\begin{enumerate}[label=(\roman*),leftmargin=*]
\item $\frac{3}{4\,N_t} \le \frac{1 + \xi_t}{Z_t} + \frac{\xi_t}{N_t} \le \frac{2}{N_t}$ and $\frac{D_k\,\pi_k^\star}{N_k} \le (\cosh a)\,\frac{D_2\,\pi_2^\star}{N_2}$ for all $k \in \{2,\ldots,n\}$.
\item Adding $\delta = n$ and $\delta = 1$ gives the quantitative decay bound
\begin{equation}
 -\ddtau{(\omega_n + \omega_1)} \ge \frac{p(1-p)}{2(n-1)}\sum_{t=2}^n \gamma_t\max\!\Biggl\{\frac{1 - \pi_t^\star}{Z_t},\; \frac{\pi_t^\star\,\one[Z_t \ge 2N_t]}{64\,N_t}\Biggr\}.
 \label{eq:U-wn-w1}
\end{equation}
\item For every $\delta \in \{2,\ldots,\min(n+3, 2n-1)\}\setminus\{n\}$,
\begin{equation}
 \ddtau{(\omega_\delta - \omega_1)} \ge \frac{1}{2(n-1)}\sum_{t\,:\,\delta\in\mathcal W_t} \gamma_t\Biggl(\frac{3\,\mu_t}{4\,N_t} - \frac{p(1-p)}{512n\,Z_t}\Biggr).
 \label{eq:U-wdelta-w1}
\end{equation}
\end{enumerate}
\end{lemma}
\begin{proof}[Proof of \cref{lem:U-lag-cancel}]
\proofstep{Step 1: Proof of \cref{eq:U-wd-ode,eq:U-post-wnw1} and part (i).}
Multiplying \cref{eq:U-beta-explicit} by $\frac{1}{2}\omega_\delta$ with $-\dot\omega_\delta = \omega_\delta\dot\brel_\delta$, $\omega_\delta\pi_{t,r} = 1/Z_t$, and $\omega_\delta e^{\brel_\delta}/N_t = 1/N_t$ yields \cref{eq:U-wd-ode}, and \cref{eq:U-post-wnw1} follows from $E_{t,t} - \mu_t = \E[\sigma(-M_t)(1 - \chi_t\rho_t)] \ge 0$ and \cref{lem:gate-calculus}(i).

Setting $x := Z_t/N_t$ on $a \in (0, 1/4]$ (where $1 < \cosh a < \frac{5}{4}$ and $x^2 + 1 \ge 2x$), we obtain
\[
  N_t\Bigl(\frac{1+\xi_t}{Z_t} + \frac{\xi_t}{N_t}\Bigr)
  = \frac{x^2 + 2x + 2\cosh a - 1}{x^2 + 2(\cosh a)x + 1} \in \Bigl[\frac{3}{4},\; 2\Bigr].
\]
Meanwhile, \cref{eq:U-rational}, $Z_2 + N_2 \le Z_k + N_k$, and $\pi_k^\star = (\omega_n Z_k)^{-1}$ yield
\[
  \frac{D_k\pi_k^\star}{N_k}
  \le \frac{2(\sinh a)/\omega_n}{(Z_k + N_k)^2}
  \le (\cosh a)\,\frac{D_2\pi_2^\star}{N_2},
\]
proving part~(i).

\proofstep{Step 2: Quantitative decay in part (ii).}
Adding \cref{eq:U-wd-ode} at $\delta = n \in \mathcal W_t$ and $\delta = 1 \in \overline{\mathcal W}_t$ with $\mu_t = \pi_t^\star E_t + R_t$ ($E_t := \E[\sigma(-M_t)]$, $R_t := \sum_{r\ne t}\pi_{t,r}\E[\sigma(-M_t)\chi_t\chi_r]$) gives the step-$t$ summand $\frac{\gamma_t}{2(n-1)}\mathcal B_t$, where
\[
  \mathcal B_t = E_t\Biggl[\frac{1-\pi_t^\star}{Z_t} + \pi_t^\star\xi_t\Bigl(\frac{1}{N_t} - \frac{1}{Z_t}\Bigr)\Biggr]
  + R_t\Bigl(\frac{\xi_t}{N_t} - \frac{1+\xi_t}{Z_t}\Bigr),
\]
with $\xi_t(\frac{1}{N_t} - \frac{1}{Z_t}) \ge 0$. When $Z_t < 2N_t$, the remainder $|R_t(\cdots)| \le \frac{(1-\pi_t^\star)p(1-p)}{128n Z_t}$ is absorbed into $\frac{3p(1-p)(1-\pi_t^\star)}{2Z_t}$. When $Z_t \ge 2N_t$, $\xi_t \ge \sigma(\log 2 + \frac{1}{4}) + \sigma(\log 2 - \frac{1}{4}) - 1 > \frac{1}{4}$ and $\pi_t^\star \ge \frac{1}{8n}$ give
\[
  \mathcal B_t
  \ge \frac{3p(1-p)(1-\pi_t^\star)}{2Z_t} + \frac{3p(1-p)\pi_t^\star}{16N_t} - \frac{p(1-p)}{512n\,N_t}
  \ge p(1-p)\max\Biggl\{\frac{1-\pi_t^\star}{Z_t},\; \frac{\pi_t^\star}{64N_t}\Biggr\},
\]
proving \cref{eq:U-wn-w1}.

\proofstep{Step 3: Exact non-letter cancellation in part (iii).}
Every visible distractor $\delta \in \{2,\ldots,\min(n+3, 2n-1)\}\setminus\{n\}$ satisfies $\delta \in \mathcal W_t \cup \overline{\mathcal W}_t$ at every step, sharing the summand $\frac{\gamma_t\xi_t\mu_t}{2(n-1)N_t}$ with $\delta = 1$ whenever $\delta \in \overline{\mathcal W}_t$. Subtracting \cref{eq:U-wd-ode} at $\delta$ from $-\dot\omega_1$ cancels all non-letter steps and applies part~(i) and \cref{lem:U-corr}(ii) on $\delta \in \mathcal W_t$ to yield \cref{eq:U-wdelta-w1}.
\end{proof}

%-----------------------------------------------------------------------------
\subsection{Phase-by-Phase Dynamics}
\label{app:J-geom-sub}
\label{app:J-constants}

We track \stepode{} across Phases~0--3 in order. \Cref{lem:U-mobility} uses the conserved Gram invariant $\Delta_0$ to bound $\kappa(a)$ from both sides and integrate $g\,\partial g/\partial a$. On any routing-window interval satisfying \cref{def:routing-window}, \cref{lem:U-floor-sign,lem:U-floor-vis,lem:U-floor-max,lem:U-floor-pi} control the signs and lower bounds of distractor reciprocal weights. \Cref{lem:U-preroute,lem:U-postroute,lem:joint-explicit} then bound prior warm-up ($\tau_0$), full-prefix lag routing ($\tau_1$), and worst-case margin convergence ($T_{\mathrm{step}}$).

\begin{lemma}[Conserved-Gram mobility and coupled $(g, a)$ envelope]
\label{lem:U-mobility}
Under \cref{def:gate} with $a(0) = a_0 > 0$ and $a_\star := \initscale^4/\sqrt{d_k}$, for all $a \ge a_0$ the mobility in \stepode{} satisfies
\[
 \max\Bigl\{\frac{\initscale^4}{8\sqrt 2\,d_k},\; \frac{a}{\sqrt{2d_k}}\Bigr\} \le \kappa(a) \le \frac{2\initscale^4}{d_k} + \frac{a}{\sqrt{2d_k}}.
\]
Moreover, on any interval $[\tau_0, T)$ where $g(\tau_0) \ge 0$, $a(\tau_0) \ge a_0$, and $\mu_t(\tau) > 0$ for all $t \in \{2,\ldots,n\}$, both $g$ and $a$ strictly increase, $g(\tau) \le g(\tau_0) + 4\initscale d_k^{1/4}\sqrt{\max\{a(\tau), a_\star\}}$ while $a(\tau) \le \frac{1}{4}$, and
\begin{equation}
 g(\tau) \ge \begin{cases}
 \dfrac{\sqrt{d_k}}{4\sqrt 2\,\initscale}\sqrt{a(\tau)^2 - a(\tau_0)^2}, & a(\tau) \in \bigl[a(\tau_0),\; \max\{a(\tau_0), 2a_\star\}\bigr],\\[6pt]
 \dfrac{1}{6}\,\initscale d_k^{1/4}\sqrt{a(\tau) - \one[a(\tau_0) \ge 2a_\star]\,a(\tau_0)}, & a(\tau) \in \bigl[\max\{a(\tau_0), 2a_\star\},\; 4\bigr].
 \end{cases}
 \label{eq:U-ga-envelope}
\end{equation}
\end{lemma}
\begin{proof}[Proof of \cref{lem:U-mobility}]
Since $\|q_0\|_2, \|k_0\|_2 \in [\initscale/2, 2\initscale]$, we have $\Delta_0 \le \|q_0\|_2^2 + \|k_0\|_2^2 \le 8\initscale^2$ and
\[
  \Delta_0^2 \ge \bigl\|q_0 q_0^\top - k_0 k_0^\top\bigr\|_{\frob}^2 \ge \frac{1}{8}\initscale^4 - \frac{4d_k}{\initscale^4}a_0^2,
  \qquad\text{so}\qquad
  \initscale^4\Delta_0^2 + 8d_k a^2 \ge \frac{1}{8}\initscale^8 + 4d_k a^2 \quad (\forall a \ge a_0),
\]
giving the two-sided bound on $\kappa(a)$.

Dividing $g\dot g$ by $\dot a$ in \cref{eq:J-flow} cancels the positive step average $\frac{1}{n-1}\sum_t(\sigma'(u_t+a)+\sigma'(u_t-a))\mu_t > 0$ and applies \cref{eq:U-rational} to yield
\[
  \frac{v^2\tanh a}{\kappa(a)} \le g\,\frac{\partial g}{\partial a} \le \frac{v^2\sinh a}{\kappa(a)}.
\]
Integrating $g\,\rmd g$ from $\tau_0$ yields the upper bound on $[a(\tau_0), 1/4]$ (using $\sinh a \le \frac{5}{4}a$ and $\kappa(a) \ge a/\sqrt{2d_k}$) together with both branches of \cref{eq:U-ga-envelope} (using $\kappa(a) < 4\initscale^4/d_k$, $\tanh a \ge a/2$ on $[a(\tau_0), 2a_\star]$ and $\kappa(a) \le 2a/\sqrt{d_k}$, $\tanh a > a/5$ on $[\max\{a(\tau_0), 2a_\star\}, 4]$).
\end{proof}

\begin{definition}[Routing-window invariants]
\label{def:routing-window}
An interval $J \subseteq [\tau_0, \tau^\star)$ satisfies the \emph{routing-window invariants} if throughout $J$:
\begin{equation}
  \text{(R)}\;\; \min_{2\le t\le n}\pi_t^\star \ge \frac{1}{8n}
  \qquad\text{and}\qquad
  \text{(U)}\;\; \min_{2\le t\le n} u_t \ge -\log 15 .
  \label{eq:U-floor-RU}
\end{equation}
On any such interval $J$, the hypotheses of \cref{lem:U-corr} hold by \cref{lem:U-preroute}(i) and the definition of $\tau^\star$, so \cref{lem:U-corr}(i)--(iii) and \cref{lem:U-lag-cancel} apply. We write the step-$t$ summands of \cref{eq:U-wd-ode} as
\[
  A_t := \frac{\gamma_t\,\xi_t\,\mu_t}{2(n-1)N_t},
  \qquad
  L_{t,\delta} := \frac{\gamma_t\bigl(\E[\sigma(-M_t)\chi_t\chi_r] - (1+\xi_t)\mu_t\bigr)}{2(n-1)Z_t}
  \quad (r = n+t-\delta),
\]
so that $-\dot\omega_\delta = \sum_{t:\delta\in\mathcal W_t} L_{t,\delta} + \sum_{t:\delta\in\overline{\mathcal W}_t} A_t$, $\operatorname{sign}A_t = \operatorname{sign}u_t$, $W := \omega_n + \omega_1$, and $-\dot W = \sum_{t=2}^n S_t$.
\end{definition}

Under the routing-window invariants \cref{eq:U-floor-RU}, prior calibration makes every letter-step contrast repulsive on distractors ($L_{t,\delta} \le 0$) while throttling positive non-letter drift.

\begin{lemma}[Sign of distractor drift]
\label{lem:U-floor-sign}
On any interval $J \subseteq [\tau_0, \tau^\star)$ satisfying \cref{def:routing-window}, for every distractor $\delta \in \mathcal W_t$ with $r = n+t-\delta \ne t$ we have $L_{t,\delta} \le 0$; that is, letter steps never decrease $\omega_\delta$. Moreover, for every lag $\delta \in \overline{\mathcal W}_t$,
\begin{equation}
  A_t^+ \le \frac{\gamma_t\,\mu_t}{2(n-1)}\,\min\{\omega_1, \omega_\delta\}.
  \label{eq:U-floor-throttle}
\end{equation}
\end{lemma}
\begin{proof}[Proof of \cref{lem:U-floor-sign}]
Since $\sigma$ is strictly increasing, condition~(U) and $a \le \frac{1}{4}$ give
\[
  1 + \xi_t = \sigma(u_t+a) + \sigma(u_t-a)
  \ge 2\sigma\bigl(u_t - \tfrac{1}{4}\bigr)
  \ge 2\sigma\bigl(-\log 15 - \tfrac{1}{4}\bigr) > \frac{1}{16}.
\]
Combining $\mu_t \ge p(1-p)\pi_t^\star \ge \frac{p(1-p)}{8n}$ (\cref{lem:U-corr}(iii) and~(R)) with $\lvert\E[\sigma(-M_t)\chi_t\chi_r]\rvert < \frac{p(1-p)}{512n}$ (\cref{lem:U-corr}(ii)) yields
\[
  \E\bigl[\sigma(-M_t)\chi_t\chi_r\bigr] - (1 + \xi_t)\mu_t
  < \frac{p(1-p)}{512n} - \frac{p(1-p)}{128n} < 0,
\]
proving $L_{t,\delta} \le 0$. For \cref{eq:U-floor-throttle}, $A_t^+ > 0$ forces $\xi_t \in (0,1)$, and $N_t \ge e^{\brel_0} + e^{\brel_1} + e^{\brel_\delta} \ge \max\{\omega_1^{-1}, \omega_\delta^{-1}\}$ because lags $0$, $1$, and $\delta$ all lie in $\overline{\mathcal W}_t$.
\end{proof}

\Cref{lem:U-lag-cancel}(iii) and a tilted functional against $\omega_1$ give an unconditional $\Omega(1/n)$ lower bound for every visible distractor lag in $\mathcal W_2$, even when $\omega_1(s) \to 0$.

\begin{lemma}[Visible-distractor Lyapunov floor]
\label{lem:U-floor-vis}
On any interval $J \subseteq [\tau_0, \tau^\star)$ satisfying \cref{def:routing-window}, set $\theta_n := \frac{5}{32(n-1)} \in (0, \frac{5}{32}]$. Then:
\begin{enumerate}[label=(\roman*),leftmargin=*]
\item For every visible distractor $\delta \in \{2,\ldots,\min(n+3, 2n-1)\}\setminus\{n\}$, the map $s \mapsto \omega_\delta(s) - \omega_1(s)$ is nondecreasing on $J$, so $\omega_\delta(s) \ge \omega_1(s) - \frac{1}{49n^2}$.
\item For every $\delta \in \{2,\ldots,\min(n+1, 2n-1)\}\setminus\{n\}$, the tilted functional $\Phi_\delta(s) := \omega_\delta(s) - (1-\theta_n)\omega_1(s)$ is nondecreasing on $J$; consequently $\omega_\delta(s) > \frac{1}{7n} = \frac{c_{\mathrm F}}{2n}$ for all $s \in J$, even as $\omega_1(s) \to 0$.
\end{enumerate}
\end{lemma}
\begin{proof}[Proof of \cref{lem:U-floor-vis}]
For~(i), every visible distractor $\delta \le n+3 = i_2 \le i_t$ lies in $\mathcal W_t \cup \overline{\mathcal W}_t$ while $1 \in \overline{\mathcal W}_t$ at every step. Using $e^{-u_t} \le 15$ from~(U), $\pi_t^\star \ge \frac{1}{8n}$ from~(R), and $\mu_t \ge p(1-p)\pi_t^\star$ from \cref{lem:U-corr}(iii), we have
\[
  \frac{p(1-p)}{512n\,Z_t}
  \le \frac{15}{512n\pi_t^\star}\,\frac{\mu_t}{N_t}
  < \frac{\mu_t}{4\,N_t},
\]
so \cref{lem:U-lag-cancel}(iii) gives
\[
  \ddtau{(\omega_\delta - \omega_1)} = \sum_{t:\delta\in\mathcal W_t}(A_t - L_{t,\delta})
  \ge \sum_{t:\delta\in\mathcal W_t}\frac{\gamma_t\,\mu_t}{4(n-1)N_t} > 0,
\]
while $\lvert\brel_\delta(\tau_0)\rvert, \lvert\brel_1(\tau_0)\rvert \le \frac{1}{100n^2}$ (\cref{lem:U-preroute}(i)) gives $\omega_\delta(\tau_0) - \omega_1(\tau_0) \ge -2\sinh\frac{1}{100n^2} \ge -\frac{1}{49n^2}$.

For~(ii), every $\delta \in \{2,\ldots,\min(n+1, 2n-1)\}\setminus\{n\}$ lies in $\mathcal W_2$, so keeping the $t = 2$ summand gives $\ddtau{(\omega_\delta - \omega_1)} \ge \frac{p(1-p)g}{4(n-1)}\frac{D_2\pi_2^\star}{N_2}$. Meanwhile, for every $k \in \{2,\ldots,n\}$, $\xi_k^+ < 1$, $\mu_k \le 3p(1-p)\pi_k^\star$, $\cosh(1/4) < \frac{16}{15}$, and prefix nesting in \cref{lem:U-lag-cancel}(i) yield
\[
  A_k^+ = \frac{g D_k\xi_k^+\mu_k}{2(n-1)N_k} \le \frac{8p(1-p)g}{5(n-1)}\,\frac{D_2\pi_2^\star}{N_2} \le \frac{32}{5}\,\ddtau{(\omega_\delta - \omega_1)}.
\]
Summing over $k = 2,\ldots,n$ bounds $-\dot\omega_1 = \sum_{k=2}^n A_k \le \sum_{k=2}^n A_k^+ \le \theta_n^{-1}\ddtau{(\omega_\delta - \omega_1)}$, so $\dot\Phi_\delta = \ddtau{(\omega_\delta - \omega_1)} - \theta_n(-\dot\omega_1) \ge 0$. Since $\omega_1(s) > 0$, $\omega_\delta(\tau_0) \ge 1 - \frac{1}{100n^2}$, and $\omega_1(\tau_0) \le 1 + \frac{1}{99n^2}$, we get $\omega_\delta(s) > \Phi_\delta(\tau_0) \ge \theta_n - \frac{2}{99n^2} > \frac{1}{7n}$ for all $n \ge 2$.
\end{proof}

In the opposite direction, comparing positive distractor drift at each step against $-\dot W = -\ddtau{(\omega_n + \omega_1)}$ gives a uniform upper bound on all distractor reciprocal weights.

\begin{lemma}[Bounded distractor reciprocal weights]
\label{lem:U-floor-max}
Let $\Omega_{\max}(s) := \max_{\delta\in\{1,\ldots,2n+1\}\setminus\{n\}}\omega_\delta(s)$. On every subinterval $J' \subseteq J$ on which, for each step $t$, either $\pi_t^\star \le \frac12$ or $u_t \ge \log 2$,
\[
  \Omega_{\max}(s) \le C_{\max} := 390 \qquad (s \in J'),
\]
and on subintervals where $\max_t\pi_t^\star \le \frac12$ the sharper bound $\Omega_{\max}(s) \le 14$ holds.
\end{lemma}
\begin{proof}
Fix $\delta \ne n$. At a letter step, $-L_{t,\delta} \le \frac{\gamma_t}{2(n-1)Z_t}\bigl(2\mu_t + \frac{p(1-p)}{512n}\bigr)$; at a non-letter step, $-A_t > 0$ holds only if $u_t < 0$ ($N_t > Z_t$), so $-A_t \le \frac{\gamma_t\mu_t}{2(n-1)Z_t}$. Using $\mu_t \le 3p(1-p)\pi_t^\star$ (\cref{lem:U-corr}(iii)), every step contributes to $\dot\omega_\delta$ at most
\[
  \frac{\gamma_t p(1-p)}{2(n-1)Z_t}\Bigl(6\pi_t^\star + \frac{1}{512n}\Bigr)
  \le \begin{cases}
    \bigl(6 + \tfrac{1}{256n}\bigr)\,S_t, & \text{if } \pi_t^\star \le \frac{1}{2},\\[4pt]
    64\bigl(6 + \tfrac{1}{512n\pi_t^\star}\bigr)\dfrac{N_t}{Z_t}\,S_t \le 193\,S_t, & \text{if } \pi_t^\star > \frac{1}{2} \text{ and } u_t \ge \log 2,
  \end{cases}
\]
by \cref{eq:U-wn-w1}. Summing over $t$ and integrating from $\tau_0$ with $\omega_\delta(\tau_0) \le 1.0026$ and $W(\tau_0) \le 2.0051$ yields $\Omega_{\max} < 390$ (and $\Omega_{\max} < 14$ whenever $\max_t\pi_t^\star \le \frac{1}{2}$).
\end{proof}

A lower floor on distractor reciprocal weights (unlike an upper bound on $\Omega_{\max}$ alone) preserves the target-lag routing condition~(R) whenever $W(s) \le W(\tau_0)$.

\begin{lemma}[Floors imply the routing condition (R)]
\label{lem:U-floor-pi}
If at time $s \in [\tau_0, \tau^\star)$ every distractor satisfies $\omega_\delta(s) \ge f$ and $W(s) \le W(\tau_0)$, then $\pi_t^\star(s) \ge \bigl(1 + 2.006(n-1)/f\bigr)^{-1}$ for all $t$; in particular $f \ge \frac{2.006(n-1)}{8n-1}$ (e.g.\ $f \ge 0.26$) implies (R). An upper bound on $\Omega_{\max}$ alone gives only the reverse inequality $\pi_t^\star \le (1 + (n-1)\omega_n/\Omega_{\max})^{-1}$.
\end{lemma}
\begin{proof}
This follows immediately from $\pi_t^\star = \bigl(1 + \sum_{r\ne t}\omega_n/\omega_{n+t-r}\bigr)^{-1}$ and $\omega_n \le W(s) \le W(\tau_0) \le 2.0051$.
\end{proof}

\paragraph{Hypothesis (F) (distractor floors).}\label{app:hypF}
\emph{Set $c_{\mathrm F} := \frac{2}{7}$. On $[\tau_0, \tau_1]$: (R) and (U) hold; every visible distractor satisfies $\omega_\delta \ge \frac{c_{\mathrm F}}{2n} = \frac{1}{7n}$; every late distractor $\delta \in \{n+4,\ldots,2n-1\}$ ($n \ge 5$) satisfies $\omega_\delta \ge \frac12$; and on $[\tau_1, \infty)$ every non-target lag $\delta \in \{2,\ldots,2n+1\}\setminus\{n\}$ satisfies $\omega_\delta \ge \frac{c_{\mathrm F}}{4n}$.}

A naive bound $\omega_1 \ge \frac{1}{n}$ cannot hold up to $\tau_1$, since $\omega_n(\tau_1) + \omega_1(\tau_1) = \frac{c_{\mathrm F}}{512n^2}$. Instead, \cref{lem:U-floor-vis}(ii) proves $\omega_\delta > \frac{1}{7n}$ unconditionally on $J$ for every $\delta \in \{2,\ldots,\min(n+1, 2n-1)\}\setminus\{n\}$ (including the bottleneck lag $\delta = 2$) via the tilted functional $\Phi_\delta$. Hypothesis~(F) supplies the remaining routing-window invariance of (R), (U), and $\int_{\tau_0}^{\tau_1}[-\dot\omega_\delta]_+\,\rmd s \le \frac12 - O(n^{-2})$ for late distractors, together with the post-routing floor on $[\tau_1, \infty)$. \Cref{tab:U-floor-audit} verifies (F) both along full \stepode{} ($38$ RK4 runs, $n \in \{4,6,8,10\}$) and in the small-$\initscale$ limiting reciprocal-weight ODE up to $n = 2{,}000$: $n\,\omega_2$ converges monotonically toward $2.22$ ($> 15\times$ above $\frac{1}{7n}$), every late distractor stays $\ge 0.9995$, $\min_t n\pi_t^\star = 1.00$, and $u_t \ge -\log 2 > -\log 15$.

Using the mobility envelope and distractor floors above, \cref{lem:U-preroute} bounds prior warm-up on $[0, \tau_0]$ (Phase~0) and full-prefix lag routing on $[\tau_0, \tau_1]$ while $a \in (0, 1/4]$ (Phase~1).

\begin{lemma}[Phases 0 and 1: Prior warm-up and fast full-prefix lag routing]
\label{lem:U-preroute}
Under \cref{def:gate} with $a(0) = a_0 > 0$, set $a_\star := \initscale^4/\sqrt{d_k}$ and $\tau^\star := \inf\{\tau \ge 0 : a(\tau) \ge 1/4 \text{ or } \max_t\gamma_t(\tau) \ge \frac{p(1-p)}{1024n}\}$ for $\initscale \le \initscale^\star(p,n,d_k)$ in \cref{eq:eps-star}. Then:
\begin{enumerate}[label=(\roman*),leftmargin=*]
\item At $\tau_0 := \frac{8(|c^\star|+1)\log(2048n/[p(1-p)])}{p(1-p)\initscale^2}$, we have $|\sigma(-\cnt(\tau_0)) - p| \le \frac{p(1-p)}{2048n}$, $a(\tau_0) \in [a_0, \frac{11}{10}a_0]$, $g(\tau_0) \in [0, \frac{1}{8}\initscale^3]$, and $|\brel_\delta(\tau_0)| \le \frac{1}{100n^2}$, and $|\sigma(-\cnt(\tau)) - p| \le \frac{p(1-p)}{1024n}$ holds for all $\tau \in [\tau_0, \tau^\star)$.
\item Whenever $\mu_t > 0$ on $[\tau_0, \infty)$, $(g, a)$ satisfy the two-sided mobility and coupled envelope bounds of \cref{lem:U-mobility}, with $g(\tau) \le 4\initscale d_k^{1/4}\sqrt{\max\{a(\tau), a_\star\}}$ for $a(\tau) \le \frac{1}{4}$.
\item Under Hypothesis~(F), there exists $\tau_1 \le \tau_0 + \frac{C n\sqrt{d_k}\log(8/r_{\mathrm{align}})}{p(1-p)\initscale^5} < \tau^\star$ at which $\min_{2\le t\le n}\pi_t^\star(\tau_1) \ge \frac{7}{8}$, $\omega_n(\tau_1) + \omega_1(\tau_1) = \frac{c_{\mathrm F}}{512n^2}$, $\min_{\delta\in\{2,\ldots,2n-1\}\setminus\{n\}} \omega_\delta(\tau_1) \ge \frac{c_{\mathrm F}}{2n}$, and $a(\tau_1) \le a_1 := \frac{C\initscale^2 n\log(8n)}{p(1-p)\sqrt{d_k}} \le \frac{1}{4}$.
\end{enumerate}
\end{lemma}
\begin{proof}[Proof of \cref{lem:U-preroute}]
\proofstep{Part (i): Prior warm-up.}
While $a \le 2a_0$ on $[0, \tau_0]$, $D_t \le 2a \le 4a_0$ and $|\dot g| \le 2\initscale^2 v^2 a \le 16\initscale^4 a_0$ keep $0 \le g(s) \le 16\initscale^4 a_0 \tau_0 \le \frac{1}{8}\initscale^3$ and $\gamma_t(s) \le 4a_0 g(s) = O(\frac{\initscale^{10}\log n}{d_k}) \ll \frac{p(1-p)}{4096n}$ for $\initscale \le \initscale^\star$. Since $\ddtau{\sigma(-\cnt)} = -\sigma'(\cnt)\dot\cnt$ and $|\frac{1}{n-1}\sum_t\E[\sigma(-M_t)\chi_t] - (\sigma(-\cnt) - p)| \le \frac{1}{4}\max_t\gamma_t$, exponential contraction drives $|\sigma(-\cnt(\tau_0)) - p| \le \frac{p(1-p)}{2048n}$ while $a(\tau_0) \le \frac{11}{10}a_0$ and $|\brel_\delta(\tau_0)| \le \frac{1}{100n^2}$.

On $[\tau_0, \tau^\star)$, whenever $|\sigma(-\cnt) - p| = \frac{p(1-p)}{1024n} \ge \max_t\gamma_t(s)$, the lower bound $\sigma'(\cnt) \ge \frac{1}{2}p(1-p)$ gives
\[
  \ddtau{\bigl(\sigma(-\cnt) - p\bigr)^2}
  \le -\initscale^2 p(1-p)\lvert\sigma(-\cnt) - p\rvert\Bigl(\lvert\sigma(-\cnt) - p\rvert - \tfrac{1}{4}\max_t\gamma_t\Bigr) < 0,
\]
trapping $|\sigma(-\cnt(s)) - p| \le \frac{p(1-p)}{1024n}$. Part~(ii) follows directly from \cref{lem:U-mobility} since $g(\tau_0)^2 \le \frac{1}{64}\initscale^6 \ll \initscale^2\sqrt{d_k}a_\star$.

\proofstep{Part (iii): Collapse of $\omega_n + \omega_1$ and routing before gate escape.}
Let $\tau_1$ be the first time $s \ge \tau_0$ at which $\omega_n(s) + \omega_1(s) = \frac{c_{\mathrm F}}{512n^2}$.
\begin{enumerate}[label=(\alph*),leftmargin=*]
\item \emph{Distractor floor across all lags (under (F)).} By \cref{eq:U-wn-w1}, $W = \omega_n + \omega_1$ decreases monotonically from $W(\tau_0) \le 2e^{1/(100n^2)} < 3$. On $[\tau_0, \tau_1]$, \cref{lem:U-floor-sign} shows that letter steps never lower a distractor weight, \cref{lem:U-floor-vis}(ii) proves $\omega_\delta > \frac{1}{7n} = \frac{c_{\mathrm F}}{2n}$ for all $\delta \in \{2,\ldots,\min(n+1,2n-1)\}\setminus\{n\}$, and (F) supplies the remaining floors together with (R) and (U). At $\tau_1$,
\[
  \sum_{r\ne t}\frac{\omega_n(\tau_1)}{\omega_{n+t-r}(\tau_1)}
  \le (n-1)\,\frac{c_{\mathrm F}/(512n^2)}{c_{\mathrm F}/(2n)}
  = \frac{n-1}{256n} < \frac{1}{16}
  \quad\implies\quad
  \min_{2\le t\le n}\pi_t^\star(\tau_1) > \frac{16}{17} > \frac{7}{8}.
\]

\item \emph{Score ceiling $a(\tau_1) \le a_1$.} Integrating \cref{eq:U-wn-w1} (using $\Omega_{\max} \le 14$ from \cref{lem:U-floor-max} while $\pi_t^\star \le \frac{1}{2}$ and $Z_t \le \frac{5}{W} + 44n^2$ when $Z_t < 2N_t$ under the $\frac{1}{7n}$ floor) yields
\[
  \int_{\tau_0}^{\tau_1}\frac{1}{n-1}\sum_{t=2}^n \gamma_t\mu_t\,\rmd s \lesssim \frac{n\log(8n)}{p(1-p)}.
\]
Since $a(\sigma'(u_t+a)+\sigma'(u_t-a)) < 4 D_t$ on $a \in (0, 1/4]$, integrating $\frac{a\,\dot a}{a + \sqrt 2\,a_\star} \le \frac{6\initscale^2}{\sqrt{d_k}}\frac{1}{n-1}\sum_{t=2}^n \gamma_t\mu_t$ from $\tau_0$ to $\tau_1$ gives $a(\tau_1) \le a_1 \le \frac{1}{4}$ and $\max_t\gamma_t(\tau_1) \le 4\initscale d_k^{1/4}a_1^{3/2} \le \frac{p(1-p)}{1024n}$ under \cref{eq:eps-star}, confirming $\tau_1 < \tau^\star$.

\item \emph{Physical routing time.} Substituting \cref{eq:U-ga-envelope} and $\frac{1}{n-1}\sum_t(\sigma'(u_t+a)+\sigma'(u_t-a))\mu_t \gtrsim \frac{p(1-p)}{n}$ into $\dot a$ gives
\[
  \dot a \gtrsim \begin{cases}
    \dfrac{p(1-p)\initscale^5}{n\sqrt{d_k}}\sqrt{a^2 - a(\tau_0)^2}, & a \le 2a_\star,\\[6pt]
    \dfrac{p(1-p)\initscale^3}{n d_k^{1/4}}\,a\sqrt{a - \one[a(\tau_0)\ge 2a_\star]a(\tau_0)}, & a \ge \max\{a(\tau_0), 2a_\star\},
  \end{cases}
\]
which integrates to $\tau_1 - \tau_0 \le \frac{C n\sqrt{d_k}\log(8/r_{\mathrm{align}})}{p(1-p)\initscale^5}$.\qedhere
\end{enumerate}
\end{proof}

Once $\min_t\pi_t^\star(\tau_1) \ge \frac{7}{8}$, the routed letter alignment $\chi_t\rho_t \ge \frac{3}{4}$ is positive on \emph{every} prompt regardless of $\cnt$. That pointwise alignment drives Phase~2 (gate escape) and Phase~3 (worst-case margin convergence).

\begin{lemma}[Phases 2 and 3: Post-routing centering and worst-case margin mastery]
\label{lem:U-postroute}
From the routed state at $\tau_1$ (\cref{lem:U-preroute}(iii)), under Hypothesis~(F) on $[\tau_1, \infty)$, for all $\tau \ge \tau_1$:
\begin{enumerate}[label=(\roman*),leftmargin=*]
\item $\min_{2\le t\le n}\pi_t^\star(\tau) > \frac{7}{8} > \frac{3}{4}$ holds permanently (with $\pi_t^\star(\tau) \to 1$ as $\tau \to \infty$), and on \emph{every prompt} $w \in \Sigma^n$,
\begin{equation}
 \frac{3}{4} \le 2\pi_t^\star(\tau) - 1 \le \chi_t\rho_t(\tau) \le 1, \qquad
 \frac{3}{4}\,\E\bigl[\sigma(-M_t)\bigr] \le \mu_t(\tau) \le \E\bigl[\sigma(-M_t)\bigr].
 \label{eq:U-pointwise-rho}
\end{equation}
\item Both $g(\tau)$ and $a(\tau)$ strictly increase on $[\tau_1, \infty)$, with $-3 \le u_t(\tau) \le a(\tau) + 3$ and $\frac{\max_t\gamma_t(\tau)}{\min_t\gamma_t(\tau)} \le \frac{4}{3}$ for all $\tau \ge \tau_1$ and all $t \in \{2,\ldots,n\}$.
\item By physical time $T_{\mathrm{step}} \le \tau_1 + \frac{C n\sqrt{d_k}\log(8/r_{\mathrm{align}})}{p(1-p)\initscale^5}$, every step $t \in \{2,\ldots,n\}$ satisfies $\min_{w\in\Sigma^n} M_t(\tau) \ge \frac{1}{2}$ for all $\tau \ge T_{\mathrm{step}}$, and $\lim_{\tau\to\infty}\Lstep(\tau) = 0$ with $\|\theta_{\mathrm s}(\tau)\|_2 \to \infty$.
\end{enumerate}
\end{lemma}
\begin{proof}[Proof of \cref{lem:U-postroute}]
Once $\pi_t^\star \ge \frac{7}{8}$, $|\chi_t\chi_r| = 1$ gives $\chi_t\rho_t = \pi_t^\star + \sum_{r\ne t}\pi_{t,r}\chi_t\chi_r \ge 2\pi_t^\star - 1 \ge \frac{3}{4}$ on every prompt, proving \cref{eq:U-pointwise-rho} and ensuring $\mu_t > 0$, $\dot g > 0$, and $\dot a > 0$ without requiring $|\cnt - c^\star| \ll 1$.

\proofstep{Parts (i)--(ii): Permanent routing and self-centering of $u_t$.}
Since $\mu_t > 0$ and $a > 0$, \cref{eq:U-post-wnw1} gives $-\ddtau{(\omega_n + \omega_1)} \ge 0$, trapping $\omega_n(\tau) + \omega_1(\tau) \le \frac{c_{\mathrm F}}{512n^2}$. Under Hypothesis~(F) on $[\tau_1, \infty)$ ($\omega_\delta(\tau) \ge \frac{c_{\mathrm F}}{4n}$ for all $\delta \in \{2,\ldots,2n-1\}\setminus\{n\}$),
\[
  \sum_{r\ne t}\frac{\omega_n(\tau)}{\omega_{n+t-r}(\tau)}
  \le (n-1)\,\frac{c_{\mathrm F}/(512n^2)}{c_{\mathrm F}/(4n)}
  < \frac{1}{128},
\]
which keeps $\min_t\pi_t^\star(\tau) > \frac{7}{8}$ permanently.

Since $\omega_n(\tau) + \omega_1(\tau) \le \frac{c_{\mathrm F}}{512n^2}$, both $e^{\brel_n(\tau)}$ and $e^{\brel_1(\tau)}$ exceed $\frac{512n^2}{c_{\mathrm F}}$ on $[\tau_1, \infty)$, while every $\delta \in (\overline{\mathcal W}_t \setminus \{1\}) \subseteq \{0\} \cup (\{2,\ldots,2n+1\}\setminus\{n\})$ satisfies $e^{\brel_\delta(\tau)} \le \frac{4n}{c_{\mathrm F}}$. Hence $(n+1)\frac{4n}{c_{\mathrm F}} \le \frac{8n^2}{c_{\mathrm F}} < \frac{1}{64}e^{\brel_1(\tau)}$, sandwiching $Z_t \in [e^{\brel_n}, \frac{129}{128}e^{\brel_n}]$ and $N_t \in (e^{\brel_1}, \frac{65}{64}e^{\brel_1})$ for every $t \in \{2,\ldots,n\}$; consequently $|u_t - u_{t'}| < \log(\frac{129}{128}\cdot\frac{65}{64}) < \log\frac{4}{3}$ and $|\pd{}{u}\log D_t| = |\xi_t| < 1$ give $\frac{\max_t\gamma_t}{\min_t\gamma_t} \le \frac{4}{3}$. At $\tau_1$, Condition~(U) gives $e^{\brel_1(\tau_1)} < N_2(\tau_1) \le 15 Z_2(\tau_1) < 16 e^{\brel_n(\tau_1)}$, and whenever $e^{\brel_1} \ge 16 e^{\brel_n}$ on $[\tau_1, \infty)$, every $t \in \{2,\ldots,n\}$ has $N_t > e^{\brel_1} > Z_t$ ($u_t < 0, \xi_t < 0$), forcing $\dot\brel_1 < 0 < \dot\brel_n$; thus $e^{\brel_1(\tau)} < 16 e^{\brel_n(\tau)}$ and $u_t(\tau) > -\log\frac{65}{4} > -3$ for all $\tau \ge \tau_1$. Conversely, whenever $\omega_n \le \frac{c_{\mathrm F}}{512n^2}$ and $\brel_n - \brel_1 \ge a + \frac{5}{2}$, $u_t > a + 2$ gives $1 + \xi_t > 2\sigma(2) > \frac{7}{4}$ and $\E[\sigma(-M_t)] - (1+\xi_t)\mu_t \le -\frac{1}{4}\E[\sigma(-M_t)] < 0$ ($\dot\brel_n < 0 < \dot\brel_1$), keeping $\brel_n(\tau) - \brel_1(\tau) \le a(\tau) + \frac{5}{2}$ and $u_t(\tau) \le a(\tau) + 3$ across $[\tau_1, \infty)$.

\proofstep{Part (iii): Gate escape, worst-case margin, and $\tau \to \infty$ limits.}
With $\gamma_{\max}(\tau) := \max_t\gamma_t(\tau)$ and $\tau_2 := \inf\{\tau \ge \tau_1 : a(\tau) \ge 4 \text{ or } \gamma_{\max}(\tau) \ge 1\}$, $|u_t - a| \le 7$ on $[\tau_1, \tau_2]$ keeps $\sigma'(u_t+a)+\sigma'(u_t-a) \ge \sigma'(7)$ and $\mu_t \gtrsim p(1-p)$, so \cref{eq:U-ga-envelope} gives $\tau_2 - \tau_1 \le \frac{C n\sqrt{d_k}\log(8/r_{\mathrm{align}})}{p(1-p)\initscale^5}$. For $\tau \ge \tau_2$, \cref{eq:U-pointwise-rho} gives
\begin{equation}
 M_t(w) \ge -|\cnt| + \frac{3}{4}\gamma_t \ge -|\cnt| + \frac{9}{16}\gamma_{\max}, \qquad
 \max_{w\,:\,\chi_t = -\operatorname{sign}(\cnt)} M_t(w) \le -|\cnt| + \gamma_{\max},
 \label{eq:U-margin-sandwich}
\end{equation}
which traps $|\cnt(\tau)| \le \frac{1}{2}\gamma_{\max}(\tau) + |c^\star| + 2$ ($\operatorname{sign}(\cnt)\dot\cnt < 0$ on the boundary) and enforces $\min_{t,w} M_t(w) \ge \frac{1}{16}\gamma_{\max}(\tau) - |c^\star| - 2$.

While $\min_{t,w} M_t(w) < \frac{1}{2}$, $\gamma_{\max} < 16(|c^\star|+3)$ keeps minority-branch derivatives $\sigma(-M_{t_\star}(w)) \gtrsim [p(1-p)]^3$ and $\dot g \gtrsim [p(1-p)]^4\initscale^4/n$, lifting $\min_{t,w} M_t \ge \frac{1}{2}$ within additional time $\lesssim \frac{n(|c^\star|+1)}{[p(1-p)]^4\initscale^4} \le \frac{Cn\sqrt{d_k}}{p(1-p)\initscale^5}$ under \cref{eq:eps-star}. Thereafter $\dot g \gtrsim \initscale^4\Lstep(\tau)$ and Barbalat's lemma on $\int_0^\infty \|\nabla_{\theta_{\mathrm s}}\Lstep\|_2^2\,\rmd\tau \le \Lstep(0)$ yield $\Lstep(\tau) \to 0$, $g(\tau) \to \infty$, and $\pi_t^\star(\tau) \to 1$.
\end{proof}

Combining Phases~0--3 yields the explicit-constant polynomial convergence guarantee for the structured step circuit.

\begin{lemma}[Explicit-constant polynomial convergence of the structured step circuit]
\label{lem:joint-explicit}
In the structured Block-$2$ circuit of \cref{def:gate}, there exist universal constants $c_0, C_0 > 0$ such that for all $n \ge 2$, $p \in (0,1)$, $r_{\mathrm{align}} \in (0,1]$, and
\begin{equation}
  0 < \initscale \le \initscale^\star(p, n, d_k) := c_0[p(1-p)]^{7/2}\frac{\min\{d_k^{1/8}, n^{1/8}\}}{n^{5/8}(\log(8n))^{3/8}},
  \label{eq:eps-star}
\end{equation}
under Hypothesis~(F), population gradient flow $\theta_{\mathrm s}(\tau)$ on $\Lstep$ satisfies $\lim_{\tau\to\infty}\Lstep(\tau) = 0$ and $\lim_{\tau\to\infty}\pi_t^\star(\brel(\tau)) = 1$ for every $t \in \{2,\ldots,n\}$, and by physical time
\begin{equation}
  T_{\mathrm{step}} \le \frac{C_0\,n\sqrt{d_k}\,\log(8/r_{\mathrm{align}})}{p(1-p)\,\initscale^5}
  \label{eq:joint-time}
\end{equation}
achieves $\pi_t^\star(\brel(\tau)) > 3/4$ and $\min_{w\in\Sigma^n} M_t(\tau) \ge 1/2$ for all $\tau \ge T_{\mathrm{step}}$, so greedy rollout from $s_1$ is exact on every prompt. Under independent isotropic draws $q(0), k(0)$, $\lvert\langle q(0),k(0)\rangle\rvert \ge \tfrac{r_{\mathrm{align}}}{2}\|q(0)\|_2\|k(0)\|_2$ holds with probability $1 - O(r_{\mathrm{align}}\sqrt{d_k})$.
\end{lemma}
\begin{proof}[Proof of \cref{lem:joint-explicit}]
Combine \cref{prop:reduced-transition} with \cref{lem:U-preroute,lem:U-postroute} under Hypothesis~(F) (using $(g,a) \mapsto (-g,-a)$ symmetry when $a(0) < 0$); $\min_{w\in\Sigma^n} M_r \ge 1/2 > 0$ for $r \in \{2,\ldots,n\}$ inductively certifies exact greedy rollout from $s_1$, and the marginal density $O(\sqrt{d_k})$ of $\langle\hat q_0, \hat k_0\rangle$ near $0$ yields the alignment probability.
\end{proof}

\begin{remark}[Limiting usage of each scaling factor in $\initscale^\star(p, n, d_k)$]
\label{rem:eps-star-roles}
Each factor in the threshold $\initscale^\star(p, n, d_k) = c_0[p(1-p)]^{7/2}\frac{\min\{d_k^{1/8}, n^{1/8}\}}{n^{5/8}(\log(8n))^{3/8}}$ of \cref{eq:eps-star} enters at a specific phase boundary of the trajectory:
\begin{itemize}[leftmargin=*,itemsep=2pt,topsep=2pt]
\item \emph{Weak-gate correlation cancellation up to $\tau_1$ (\cref{lem:U-corr}, \cref{lem:U-preroute}(iii)(b)):} Keeping $\max_t\gamma_t(\tau_1) \lesssim \initscale d_k^{1/4}a_1^{3/2} \le \frac{p(1-p)}{1024n}$ at $a_1 \asymp \frac{\initscale^2 n\log(8n)}{p(1-p)\sqrt{d_k}}$ requires $\initscale^4\frac{n^{3/2}(\log(8n))^{3/2}}{[p(1-p)]^{3/2}\sqrt{d_k}} \lesssim \frac{p(1-p)}{n}$, producing the $d_k^{1/8} n^{-5/8}(\log(8n))^{-3/8}$ scaling (which simultaneously ensures $a_1 \le \frac{1}{4}$).
\item \emph{Prior warm-up bootstrap on $[0, T_0]$ (\cref{lem:silent-routing}, Step~1):} Bounding $n^2\initscale^6 K^2 \lesssim p(1-p)d_k$ across $K = \initscale^2 T_0 = \widetilde O([p(1-p)]^{-2})$ so that $a$ and $\brel$ move negligibly before $\cnt \to c^\star$ contributes the $[p(1-p)]^{7/2}$ prefactor.
\item \emph{Post-escape minority-branch margin phase (\cref{lem:U-postroute}(iii)):} Absorbing the post-escape time $\lesssim \frac{n(|c^\star|+1)}{[p(1-p)]^4\initscale^4}$ into the saddle-escape bound $\frac{Cn\sqrt{d_k}}{p(1-p)\initscale^5}$ when $d_k \gg n$ contributes the $\min\{d_k^{1/8}, n^{1/8}\}$ cap; at linear head width $d_k = \Theta(n)$ (\cref{cor:matched}(a)), \cref{eq:eps-star} simplifies to the standard fan-in scale $\initscale = \widetilde\Theta(n^{-1/2})$.
\end{itemize}
\end{remark}

%-----------------------------------------------------------------------------
\subsection{Main Proof and Clamped Baseline}
\label{app:J-silent-obstruction}

\Cref{lem:silent-routing} bounds attention movement and loss variation across the $H(p)$ plateau, and \cref{lem:clamped-shortcut} computes early lag drift when the state-only prior is held at $\cnt \equiv 0$. Together with \cref{lem:joint-explicit}, these results prove \cref{thm:joint} and \cref{cor:joint-outcome,cor:matched}.

\begin{lemma}[Silent attention routing on the flat-loss plateau]
\label{lem:silent-routing}
Assume the hypotheses of \cref{lem:joint-explicit} with $a(0) = a_0 > 0$, let $H(p) := -p\log p - (1-p)\log(1-p)$, and set
\[
  T_0 := \frac{16(|c^\star| + 1)}{p(1-p)\initscale^2}\log\frac{2048n}{p(1-p)\initscale^2} = \widetilde O(\initscale^{-2}).
\]
\begin{enumerate}[label=(\alph*),leftmargin=*]
\item \emph{(Calibration; no Hypothesis~(F).)} $|\cnt(T_0) - c^\star| \le \initscale^2\mixrate = \initscale^2\min\{1, \mixrate\}$ (so $\cnt(T_0) = c^\star = 0$ at $p = 1/2$), $|\Lstep(T_0) - H(p)| \le \initscale^4$, $\max_{2\le t\le n}\pi_t^\star(T_0) < \frac{5}{4n}$, and $\tau_0 \le T_0 < \tau^\star$, with $\tau_0, \tau^\star$ from \cref{lem:U-preroute}.
\item \emph{(Silent routing; under Hypothesis~(F).)} Let $\tau_1$ be the time of \cref{lem:U-preroute}(iii). Then $T_0 < \tau_1$, $\min_{2\le t\le n}\mathrm{d}_{\mathrm{TV}}(\pi_{t,\cdot}(\tau_1), \pi_{t,\cdot}(T_0)) \ge \frac{7}{8} - \frac{5}{4n} \ge \frac{1}{4}$, and
\[
  0 \le \Lstep(T_0) - \Lstep(\tau_1) \le \Biggl(1 + \frac{C n^{3/2}(\log(8n))^{3/2}}{\sqrt{p(1-p)d_k}}\Biggr)\initscale^4.
\]
\item \emph{(Unchanged predictions; under Hypothesis~(F).)} Let $\Gamma_1 := 3\initscale d_k^{1/4}\max\{a_1, a_\star\}^{3/2}$ and $\mixrate_{\min} := \frac{2\Gamma_1}{p(1-p)} \le \frac{C n^{3/2}(\log(8n))^{3/2}}{(p(1-p))^{5/2}\sqrt{d_k}}\initscale^4$. If $\mixrate > \mixrate_{\min}$, then prompt-wise argmax predictions agree with the state-only transition decision throughout $[T_0, \tau_1]$.
\end{enumerate}
\end{lemma}
\begin{proof}[Proof of \cref{lem:silent-routing}]
Write $e := \cnt - c^\star$, $\gamma_{\max}(s) := \max_t|\gamma_t(s)|$, $L_0 := \log\frac{2048n}{p(1-p)\initscale^2}$, and $K := \initscale^2 T_0 = \frac{16(|c^\star|+1)L_0}{p(1-p)}$, taking $c_0 \le \frac{1}{50}$ in \cref{eq:eps-star}.

\proofstep{Step 1: A priori bounds on $[0, T_0]$.}
Since $|c^\star| + 1 \le \frac{1}{p(1-p)}$ and $\log y \le 1.5 y^{1/4}$ ($y \ge 1$), we have $L_0 \le 10.1\,n^{1/4}(p(1-p))^{-1/4}\initscale^{-1/2}$ and $2400 n^2\initscale^6 K^2 \le 6.3\times 10^7 n^{5/2}(p(1-p))^{-9/2}\initscale^5 \le p(1-p)d_k$ under $c_0 \le \frac{1}{50}$. Bootstrapping $a \in (0, 2a_0]$ on $[0, T_0]$ yields
\[
  \lvert a(s) - a_0\rvert \le \frac{6\initscale^6 K^2}{d_k}a_0 \le \frac{1}{10}a_0, \qquad
  \gamma_{\max}(s) \le \Gamma_0 := \frac{32\initscale^{10}K}{d_k} \le \frac{p(1-p)\initscale^4}{75n^2} < \frac{p(1-p)}{1024n}, \qquad
  \lvert\brel_\delta(s)\rvert \le \frac{1}{100n^2},
\]
confirming $\tau_0 \le T_0 < \tau^\star$ and $\max_t\pi_t^\star(T_0) \le \frac{e^{1/200}}{n} < \frac{5}{4n}$.

\proofstep{Step 2: Prior contraction and proof of part (a).}
Coupling $\chi_r$ ($r \in [n]$) with i.i.d.\ Rademacher signs $\zeta_r$ that disagree with probability $\frac{1-\mixrate}{2}$ (or equal $\operatorname{sign}(1-2p)$ with probability $\mixrate$) and using oddness of $h(\cnt) := \E[(\sigma(-\zeta_t(\cnt+\gamma_t\rho_t)) - \sigma(-\zeta_t\cnt))\zeta_t]$ ($\sup|h'| \le \frac{1}{10}|\gamma_t|$ since $\sup|\sigma''| = \frac{\sqrt 3}{18} < \frac{1}{10}$) gives
\[
  \dot\cnt = \initscale^2\bigl(\sigma(-\cnt) - p + \delta_c\bigr), \qquad
  \lvert\delta_c\rvert \le \gamma_{\max}\min\Bigl\{\frac{1}{4},\; \frac{\mixrate}{2} + \frac{|\cnt|}{10}\Bigr\}.
\]
This traps $\cnt(s) \in [\min\{0,c^\star\}-\initscale^2, \max\{0,c^\star\}+\initscale^2]$ (where $\sigma' \ge \frac{1}{2}p(1-p)$). Applying Gr\"onwall's inequality to $|e|$ and expanding $\ell(m) := \log(1+e^{-m})$ around $\chi_t c^\star$ gives
\begin{align*}
  |e(T_0)| &\le |c^\star|\,e^{-\frac{2}{5}\initscale^2 p(1-p)T_0} + \frac{5\Gamma_0\bigl(\frac{\mixrate}{2} + \frac{|c^\star|}{10}\bigr)}{2p(1-p)} \le \initscale^2\mixrate,\\
  -2p(1-p)\gamma_{\max} &\le \Lstep - H(p) \le 2p(1-p)\gamma_{\max} + \frac{1}{8}\bigl(|e| + \gamma_{\max}\bigr)^2,
\end{align*}
with $e \equiv 0$ at $p = 1/2$, yielding $|\Lstep(T_0) - H(p)| \le \initscale^4$ and proving~(a).

\proofstep{Step 3: Silent routing and unchanged predictions in parts (b)--(c).}
Under~(F), $\omega_n(T_0)+\omega_1(T_0) > 1 > \frac{c_{\mathrm F}}{512n^2}$ forces $T_0 < \tau_1 < \tau^\star$, and \cref{lem:U-mobility,lem:U-preroute} bound $0 \le \gamma_t(s) \le \frac{33}{16}\initscale d_k^{1/4}\max\{a_1, a_\star\}^{3/2} \le \Gamma_1$ across $[T_0, \tau_1]$, giving
\[
  \mathrm{d}_{\mathrm{TV}}\bigl(\pi_{t,\cdot}(\tau_1), \pi_{t,\cdot}(T_0)\bigr) \ge \pi_t^\star(\tau_1) - \pi_t^\star(T_0) \ge \frac{7}{8} - \frac{5}{4n} \ge \frac{1}{4},
  \qquad
  0 \le \Lstep(T_0) - \Lstep(\tau_1) \le \initscale^4 + 2p(1-p)\Gamma_1,
\]
proving~(b). When $\mixrate > \mixrate_{\min}$, $\vartheta := \max\{\initscale^2\mixrate, \frac{\Gamma_1}{p(1-p)}\} \le \frac{\mixrate}{2}$ traps $|e(s)| \le \vartheta$ on $[T_0, \tau_1]$, so $|\cnt(s) + \gamma_t(s)\rho_t(s) - c^\star| \le \vartheta + \Gamma_1 < 2\mixrate \le |c^\star|$ preserves $\operatorname{sign}(\cnt + \gamma_t\rho_t) = \operatorname{sign}(c^\star)$ on every prompt, proving~(c).
\end{proof}

\Cref{lem:clamped-shortcut} analyzes the system when the state-only prior is fixed at $\cnt(\tau) \equiv 0$ ($p \ne 1/2$): uncalibrated residuals leave $\Theta(\mixrate^2)$ wrong-letter correlation and raise non-letter lag-$1$ drift by a factor of $\Theta(n)$.

\begin{lemma}[Clamped-prior lag drift at $\cnt \equiv 0$]
\label{lem:clamped-shortcut}
In \stepode{} with the state-only prior clamped at $\cnt(\tau) \equiv 0$ ($U \equiv 0$, so $M_t = g D_t\chi_t\rho_t$) for any $n \ge 2$ and $p \in (0,1) \setminus \{1/2\}$ ($\mixrate = \lvert 1-2p\rvert > 0$):
\begin{enumerate}[label=(\alph*),leftmargin=*]
\item For every $(\brel, g, a)$ and every wrong prefix letter $r \ne t$, $\lvert\E[\sigma(-M_t)\chi_t\chi_r] - \frac{1}{2}\mixrate^2\rvert \le \frac{1}{4}\lvert g D_t\rvert$; in particular, $\E[\sigma(-M_t)\chi_t\chi_r] \ge \frac{1}{4}\mixrate^2 > 0$ whenever $\max_t\lvert g D_t\rvert \le \mixrate^2$.
\item At $\brel = 0$, define $A_n := \frac{1}{n-1}\sum_{t=2}^n \frac{n(t+2)}{(n+t+2)^2}$, $B_n := \frac{1}{n-1}\sum_{t=2}^n \frac{2n^2(t+2)}{(n+t+2)^3}$, and $C_n := \frac{1}{n-1}\sum_{t=2}^n \frac{n^2(n-t-2)}{(n+t+2)^3}$. As $a \to 0^+$ after $g \to 0^+$, the initial velocities of true letter lag $\delta = n$ and non-letter state lag $\delta = 1$ satisfy
\[
  \lim_{a\to 0^+}\lim_{g\to 0^+}\frac{2n}{ga}\,\dot\brel_n\Big|_{\brel=0,\,\cnt=0} = A_n - B_n\Bigl(\mixrate^2 + \frac{1-\mixrate^2}{n}\Bigr), \qquad
  \lim_{a\to 0^+}\lim_{g\to 0^+}\frac{2n}{ga}\,\dot\brel_1\Big|_{\brel=0,\,\cnt=0} = C_n\Bigl(\mixrate^2 + \frac{1-\mixrate^2}{n}\Bigr).
\]
Since $A_n \to \log 2 - \frac{1}{2}$, $B_n \to \frac{1}{4}$, and $C_n \to \frac{1}{4}$ as $n \to \infty$, and $C_n > 0$ iff $n \ge 6$, clamping $\cnt \equiv 0$ elevates non-letter lag-$1$ drift to $\Theta(\mixrate^2 ga/n)$ for $n \ge 6$ (versus $O(p(1-p)ga/n^2)$ at $\cnt = c^\star$), reverses initial routing ($\dot\brel_1 > \dot\brel_n$) whenever $\mixrate^2 > 2\log 2 - 1$ ($n \ge n_0(\mixrate)$), and makes $\dot\brel_n < 0$ whenever $\mixrate^2 > 4\log 2 - 2$ ($n \ge n_1(\mixrate)$).
\end{enumerate}
\end{lemma}
\begin{proof}[Proof of \cref{lem:clamped-shortcut}]
For~(a), $\sigma(0) = \frac{1}{2}$ and independence of $\chi_t, \chi_r$ ($r \ne t$) give $\E[\frac{1}{2}\chi_t\chi_r] = \frac{1}{2}\mixrate^2$, while $\frac{1}{4}$-Lipschitzness of $\sigma$ and $|M_t| \le |g D_t|$ yield $|\sigma(-M_t)\chi_t\chi_r - \frac{1}{2}\chi_t\chi_r| \le \frac{1}{4}|g D_t|$.

For~(b), at $\brel = 0$ we have $\pi_{t,r} = \frac{1}{n}$, $Z_t = n$, $N_t = t+2$, and $u_t = \log\frac{n}{t+2}$. Sending $g \to 0^+$ at fixed $a > 0$ in \cref{eq:U-beta-explicit} gives $\frac{1}{g}\gamma_t = D_t(u_t, a)$, $\E[\sigma(-M_t)] \to \frac{1}{2}$, and $\mu_t \to \frac{1}{2}m_n$ with $m_n := \mixrate^2 + \frac{1-\mixrate^2}{n}$; sending $a \to 0^+$ then gives
\[
  \frac{1}{2a}D_t \to \frac{n(t+2)}{(n+t+2)^2}, \qquad
  1 + \xi_t \to \frac{2n}{n+t+2}, \qquad
  \xi_t \to \frac{n-t-2}{n+t+2},
\]
yielding the exact finite-$n$ limits, while Riemann-sum convergence on $[0,1]$ gives $(A_n, B_n, C_n) \to (\log 2 - \frac{1}{2}, \frac{1}{4}, \frac{1}{4})$.

For $n \ge 6$, the two negative summands $t \in \{n-1, n\}$ of $C_n$ sum to $\ge -\frac{3n^2}{(2n+1)^3}$ whereas the $t = 2$ summand $\frac{n^2(n-4)}{(n+4)^3}$ exceeds $\frac{3n^2}{(2n+1)^3}$ (since $2n+1 \ge 1.3(n+4) \implies \frac{(n-4)(2n+1)^3}{(n+4)^3} \ge 2\cdot 1.3^3 > 3$), whereas $(C_2,\ldots,C_5) = (-0.037, -0.031, -0.018, -0.003) < 0$, so $C_n > 0$ iff $n \ge 6$ with $\inf_{n\ge 6} C_n > 0$ (versus $\dot\brel_1 = O(p(1-p)ga/n^2)$ at $\cnt = c^\star$). Reversal $\dot\brel_1 > \dot\brel_n$ holds iff $m_n > \frac{A_n}{B_n + C_n} \to 2\log 2 - 1$ (at $n = 8$, $(A_8, B_8, C_8) = (0.243, 0.262, 0.038)$, needing $p \lesssim 0.05$), and $\dot\brel_n < 0$ holds iff $m_n > \frac{A_n}{B_n} \to 4\log 2 - 2$.
\end{proof}

Combining \cref{lem:U-preroute,lem:U-postroute,lem:joint-explicit,lem:silent-routing,lem:clamped-shortcut} gives \cref{thm:joint}, and \cref{cor:joint-outcome} checks the $\Omega(\mixrate^{-n})$ outcome bound from the same unmasked initialization.

\par\bigskip\noindent\textbf{Theorem~\ref{thm:joint} (Silent routing under step supervision; formal statement).}\enspace{\itshape
In the Block-$2$ circuit of \cref{def:gate}, fix any $n \ge 2$, $p \in (0,1)$, initial query--key alignment $r_{\mathrm{align}} \in (0,1]$, and $\initscale \le \initscale^\star(p,n,d_k)$ \cref{eq:eps-star}, and write $\delta_{\mathrm{plat}} := \widetilde O((1 + n^{3/2}d_k^{-1/2})\initscale^4)$. Under population gradient flow on $\Lstep$, parts~(i), (ii), and~(v) hold unconditionally, and parts~(iii)--(iv) hold under Hypothesis~(F) (\cref{app:hypF}):
\begin{enumerate}[label=(\roman*),leftmargin=*,itemsep=2pt,topsep=3pt]
\item Prior calibration: by $T_0 = \widetilde O(\initscale^{-2})$, $\lvert\cnt(T_0) - c^\star\rvert \le \initscale^2\min\{1,\mixrate\}$ and $\Lstep(T_0) = H(p) + O(\delta_{\mathrm{plat}})$.
\item Distractor cancellation: throughout $[T_0, \tau^\star)$, every distractor satisfies $\lvert\E[\sigma(-M_t)\chi_t\chi_r]\rvert \le \frac{p(1-p)}{512n}$ ($r \ne t$), whereas $\E[\sigma(-M_t)\chi_t^2] \ge \frac{3}{2}p(1-p)$.
\item Silent plateau routing: by $\tau_1 \in (T_0, \tau^\star)$, $\min_{2\le t\le n}\mathrm{d}_{\mathrm{TV}}(\pi_{t,\cdot}(\tau_1), \pi_{t,\cdot}(T_0)) \ge \frac{1}{4}$ while $0 \le \Lstep(T_0) - \Lstep(\tau_1) \le \delta_{\mathrm{plat}}$; when $\mixrate > \mixrate_{\min} = O(\delta_{\mathrm{plat}})$, argmax predictions match the state-only prior throughout $[T_0, \tau_1]$.
\item Polynomial-time rollout: by $T_{\mathrm{step}} = \widetilde O(n\sqrt{d_k}\,\initscale^{-5})$, greedy rollout from $s_1$ is exact on every prompt ($\min_{t,w} M_t \ge 1/2$, $\pi_t^\star > 3/4$), with $\Lstep(\tau) \to 0$ and $\pi_t^\star(\tau) \to 1$.
\item Clamped prior ($p \ne 1/2$): holding $\cnt \equiv 0$ leaves $\E[\sigma(-M_t)\chi_t\chi_r] \ge \frac{1}{4}\mixrate^2 > 0$ ($r \ne t$) while $\max_t\lvert g D_t\rvert \le \mixrate^2$; for $n \ge 6$ at $\brel = 0$ this amplifies non-letter lag-$1$ drift by $\Theta(n)$, reversing $\dot\brel_1 > \dot\brel_n$ for $n \ge n_0(\mixrate)$ when $\mixrate^2 > 2\log 2 - 1$ ($\dot\brel_n < 0$ when $\mixrate^2 > 4\log 2 - 2$).
\end{enumerate}}\par\smallskip

\begin{proof}[Proof of \cref{thm:joint}]
Each clause follows directly from the preceding lemmas:
\begin{itemize}[leftmargin=*]
\item Part~(i) is \cref{lem:silent-routing}(a) (without Hypothesis~(F)).
\item Part~(ii) holds unconditionally across $[T_0, \tau^\star) \subseteq [\tau_0, \tau^\star)$ by \cref{lem:silent-routing}(a), \cref{lem:U-preroute}(i), and \cref{lem:U-corr}(i)--(ii), with $[T_0, \tau_1] \subset [T_0, \tau^\star)$ under~(F) by \cref{lem:U-preroute}(iii).
\item Part~(iii) is \cref{lem:silent-routing}(b)--(c) (under~(F)), and part~(iv) is \cref{lem:joint-explicit} (under~(F)).
\item Part~(v) is \cref{lem:clamped-shortcut} (without Hypothesis~(F)).\qedhere
\end{itemize}
\end{proof}

To verify an apples-to-apples dynamical separation from a shared unmasked seed, we check that \cref{def:gate} satisfies every hypothesis of \cref{thm:outcome-end} and specialize the step and outcome bounds to linear widths.

\begin{corollary}[Outcome plateau from the shared unmasked seed]
\label{cor:joint-outcome}
From the initialization of \cref{def:gate} under the full causal mask at every position and any $p \in (0,1)\setminus\{1/2\}$ ($\mixrate = |1-2p| \in (0,1)$), gradient flow on $\Lout$ training $\theta_{\mathrm s}$ or any subset of parameter tensors up to the full model (including $\brel$) obeys the $\Omega(\mixrate^{-n})$ plateau of \cref{thm:outcome-end}, while the ambient parameter space at the widths of \cref{prop:represent} contains both $\thstep$ and $\thout$.
\end{corollary}
\begin{proof}[Proof of \cref{cor:joint-outcome}]
Grouping the hidden units of each zeroed MLP in \cref{def:gate} into pairs sharing a bias ($\pm\initscale$) and up-projection row ($0$ at $\theta_0$) with opposite down-projection columns $\pm w$ ($\|w\|_2 = O(\initscale)$) cancels $w\relu(\cdot) - w\relu(\cdot) \equiv 0$ identically while satisfying $|\tilde b_m| = 1 \ge 1$, $V(0) = 0$, $\brel(0) = 0$, and $O(\initscale)$ operator norms in \cref{ass:init}, so \cref{thm:outcome-end} applies directly.
\end{proof}

\begin{corollary}[Initialization-matched guarantees]
\label{cor:matched}
From the shared unmasked initialization (\cref{def:gate}) at linear widths ($d_k \ge 2n+6$, $\dmodel, \dmlp, d_k = \Theta(n)$) with isotropic $(q_0, k_0)$: (a)~at fan-in scale $\initscale = \widetilde\Theta(n^{-1/2})$ (probability $1 - O(n^{-3/2})$), Euclidean step flow reaches exact rollout in $T_{\mathrm{step}} = \widetilde O(n^4)$ conditional on~(F); (b)~at $\initscale = \widetilde\Theta(n^{-3})$ (probability~$1$), stateless Block-RMS step flow (\cref{thm:block-rms}) reaches exact rollout in $T_{\mathrm{RMS}} = \widetilde O(n^6)$ without~(F); and from either scale, Euclidean outcome flow on any subset of weights stalls for $0 \le \tau \le c\mixrate^{-n}$ when $p \ne 1/2$ ($\theta(\tau) \equiv \theta_0$ at $p = 1/2$; \cref{thm:outcome-end} does not bound Block-RMS outcome flow).
\end{corollary}
\begin{proof}[Proof of \cref{cor:matched}]
Let $c_1 n \le d_k \le c_2 n$ ($d_k \ge 2n+6$) and $\dmlp \le c_3 n$.
\begin{itemize}[leftmargin=*]
\item \emph{Part (a):} With $c_p := \min\{c_0\min\{c_1^{1/8}, 1\}[p(1-p)]^{7/2}, c\mixrate, c c_3^{-1/2}\}$ (the last two for $p \ne 1/2$), $\initscale = c_p n^{-1/2}(\log(8n))^{-3/8}$ satisfies \cref{eq:eps-star}, $r_{\mathrm{align}} = n^{-2}$ holds with probability $1 - O(n^{-2}\sqrt{c_2 n}) = 1 - O(n^{-3/2})$, and \cref{eq:joint-time} gives $T_{\mathrm{step}} = O(n^4(\log(8n))^{23/8}) = \widetilde O(n^4)$ under~(F).
\item \emph{Part (b):} With $c_p' := \min\{2^{-30}[p(1-p)]^{3/2}, c\mixrate, c c_3^{-1/2}\}$, $\initscale = c_p' n^{-3}(\log(8n))^{-1/2}$ satisfies \cref{thm:block-rms}, $\Pr(\langle q_0, k_0\rangle \ne 0) = 1$ under isotropic $(q_0, k_0)$, and \cref{eq:rms-time} gives exact rollout by $T_{\mathrm{RMS}} = O(n^6\log(8n)) = \widetilde O(n^6)$ without Hypothesis~(F).
\item \emph{Outcome stalling:} From either initialization, \cref{cor:joint-outcome} gives $T_{\mathrm{out}} = \Omega(\mixrate^{-n})$ under \emph{Euclidean} outcome gradient flow for $p \ne 1/2$ (note that \cref{thm:outcome-end} relies on $\ddtau{\Lout} = -\|\dot\theta\|_2^2$ and does not bound Block-RMS outcome flow, which rescales $O(\mixrate^n)$ gradients to unit block RMS). At $p = 1/2$, $V(0) = 0$ vanishes all non-$V$ outcome gradients at $\theta_0$, while $h^{(2)}(\theta_0)$ at the readout position has the single-letter linear form $f_0(s_0) + \sum_{r=1}^n f_r(s_0)\chi_r$, whose product with $(-1)^{s_n} = (-1)^{s_0}\prod_{r=1}^n \chi_r$ ($n \ge 2$, $\E[\chi_r] = 0$) has expectation $0$, giving $\nabla_\theta\Lout(\theta_0) = 0$ and $\theta(\tau) \equiv \theta_0$.\qedhere
\end{itemize}
\end{proof}

\begin{remark}[Inclusion of $t = 1$]
\label{rem:first-state}
\label{rem:joint-scope}
\label{rem:J-balance}
\Cref{thm:joint} averages $\Lstep$ over $t \in \{2,\ldots,n\}$; we do not prove the extension to $t = 1$. At $t = 1$ ($i_1 = n+2$, holding $\mathrm{SEP}$), context letters $j = 1+r$ sit at $\mathcal W_1 = \{1,\ldots,n\}$ with $w_1$ at the same true lag $\delta^\star = n$ and $s_0$ at lag $n+2 \in \overline{\mathcal W}_1$. Extending to $t = 1$ requires two modifications: copying $x_0$ to $i_1$ via a Block-$1$ head so Block-$2$ queries read $h_{i_1}^{(1)}$, and re-proving the step-$1$ counterpart of \cref{lem:U-lag-cancel}(ii)--(iii) since lag $1$ reads $w_n$ at $t = 1$.
\end{remark}

% ===================================================================
\section{Proof of \texorpdfstring{\Cref{thm:block-rms}}{Theorem 3}}
\label{app:adaptive}

\paragraph{Empirical motivation.}
At $\initscale = 0.02$, cross-tensor gradient imbalance ($\mathrm{RMS}(\nabla_{W_Q}\Lstep) \approx 2.1\times 10^{-6}$ vs.\ $\mathrm{RMS}(\nabla_V\Lstep) \approx 5.0\times 10^{-3}$) stalls tuned SGD on the $H(p)$ plateau. Our Block-$2$ analysis (\cref{thm:block-rms}) proves $O((1+|c^\star|)\initscale^{-2})$ exact rollout from $s_1$ \emph{without Hypothesis~(F)} under $\initscale \lesssim [p(1-p)]^{3/2}n^{-3}(\log n)^{-1/2}$. \texttt{Block-RMS SGD} divides each tensor's gradient by its scalar RMS norm with zero optimizer state, matching AdamW's $5/5$ and $18/20$ solve rates and speed (\cref{tab:clamp-main,tab:transfer-main}a).

\paragraph{Dimensionless variables and modular proof architecture for \cref{thm:block-rms}.}
Set the normalized value scale $v_0 := 2v/\initscale \in [1, 4]$, the distractor index set $\mathcal D := \{2,\ldots,2n-1\}\setminus\{n\}$, and dimensionless time and reciprocal weights
\[
  s := \initscale\tau \ge 0, \qquad
  x(s) := \omega_n(s/\initscale), \qquad
  y_\delta(s) := \omega_\delta(s/\initscale) \;\; (\delta \in \mathcal D), \qquad
  m(s) := \min_{\delta\in\mathcal D} y_\delta(s),
\]
writing primes $f'(s) := \frac{\rmd f}{\rmd s} = \initscale^{-1}\dot f(s/\initscale)$ for $s$-derivatives. We organize the proof of \cref{thm:block-rms} into three self-contained layers:
\begin{enumerate}[leftmargin=*,itemsep=1.5pt,topsep=2pt]
\item \textbf{General ODE comparison principles (\cref{app:rms-compare-sub}):} Two model-independent lemmas establish stability of $(x, m)$ under an integrable perturbation (\cref{lem:inv-pert}) and finite-time threshold crossing / scalar barrier invariance (\cref{lem:hit-compare}).
\item \textbf{Three structural lemmas and a concise proof of \cref{thm:block-rms} (\cref{app:rms-routing-sub}):} Gate saturation (\cref{lem:rms-align}), stable distractor weights (\cref{lem:rms-contrast}), and finite-time routing (\cref{lem:rms-routing-time}) reduce \cref{thm:block-rms} to a half-page synthesis.
\item \textbf{Technical verifications (\cref{app:rms-tech-sub}):} Isolates Carath\'eodory well-posedness and symmetries (\cref{lem:rms-wellposed}), the positivity bootstrap and explicit constant checks proving \cref{lem:rms-routing-time}, and the post-routing readout barrier (\cref{lem:rms-margin-trap}).
\end{enumerate}

\subsection{Comparison Lemmas}
\label{app:rms-compare-sub}

We first isolate two model-independent ODE comparison principles in normalized time $s = \initscale\tau$: stability of the reciprocal-weight envelope under an integrable perturbation (\cref{lem:inv-pert}), and finite-time hitting under linear negative drift paired with scalar barrier invariance (\cref{lem:hit-compare}).

\begin{lemma}[Stability under integrable perturbations]
\label{lem:inv-pert}
Let $\mathcal D$ be a finite nonempty index set, let $x : [s_a, s_b) \to \R$ and $y_\delta : [s_a, s_b) \to \R$ ($\delta \in \mathcal D$, $-\infty < s_a < s_b \le \infty$) be locally absolutely continuous, and set $m(s) := \min_{\delta\in\mathcal D} y_\delta(s)$. Suppose $r : [s_a, s_b) \to [0, \infty)$ satisfies $\int_{s_a}^{s_b} r(u)\,\rmd u \le \eta < \infty$, and for almost every $s \in [s_a, s_b)$,
\[
  x'(s) \le r(s) \qquad\text{and}\qquad y_\delta'(s) \ge -r(s) \;\;\text{for every } \delta \in \arg\min_{\delta'\in\mathcal D} y_{\delta'}(s).
\]
Then $m$ is locally absolutely continuous with $m'(s) \ge -r(s)$ a.e.\ on $[s_a, s_b)$, and for all $s_a \le s_1 \le s < s_b$:
\begin{enumerate}[label=(\roman*),leftmargin=*,itemsep=1pt,topsep=2pt]
\item \emph{(Two-sided stability.)}
\begin{equation}
  x(s) \le x(s_1) + \eta, \qquad m(s) \ge m(s_1) - \eta.
  \label{eq:rms-floor-shift}
\end{equation}
\item \emph{(Ratio persistence.)} If $K x(s_1) \le m(s_1)$ ($K > 0$) and $\eta < m(s_1)$, then $\frac{x(s)}{m(s)} \le \frac{1/K + \eta/m(s_1)}{1 - \eta/m(s_1)}$.
\end{enumerate}
\end{lemma}
\begin{proof}[Proof of \cref{lem:inv-pert}]
Because $\mathcal D$ is finite and $|m(u) - m(v)| \le \sum_{\delta\in\mathcal D}|y_\delta(u) - y_\delta(v)|$, the envelope $m$ is locally absolutely continuous on $[s_a, s_b)$ and hence differentiable simultaneously with every $y_\delta$ ($\delta \in \mathcal D$) at a.e.\ $s \in (s_a, s_b)$.

At any such joint differentiability point $s$ where the hypothesis holds, pick $\delta \in \arg\min_{\delta'\in\mathcal D} y_{\delta'}(s)$, so $m(s) = y_\delta(s)$. For $h > 0$, we have $m(s + h) - m(s) \le y_\delta(s + h) - y_\delta(s)$ and $m(s) - m(s - h) \ge y_\delta(s) - y_\delta(s - h)$; dividing by $h$ and letting $h \downarrow 0$ gives $m'(s) = y_\delta'(s) \ge -r(s)$. Integrating $x'(u) \le r(u)$ and $m'(u) \ge -r(u)$ over $[s_1, s]$ yields~(i), and dividing $x(s) \le m(s_1)/K + \eta$ by $m(s) \ge m(s_1) - \eta > 0$ yields~(ii).
\end{proof}

Next, we record the hitting-time bound for $x'(s) \le -\alpha s + r(s)$ and a standard mean-value barrier test used in the prior bootstrap and post-routing margin trap.

\begin{lemma}[Quadratic-drift hitting time and scalar barrier invariance]
\label{lem:hit-compare}\leavevmode
\begin{enumerate}[label=(\roman*),leftmargin=*,itemsep=1pt,topsep=2pt]
\item \emph{(Hitting time under linear drift.)} Let $0 \le s_a < s_b \le \infty$, let $x : [s_a, s_b) \to [0, \infty)$ be locally absolutely continuous, let $r : [s_a, s_b) \to [0, \infty)$ satisfy $\int_{s_a}^{s_b} r(u)\,\rmd u \le \eta < \infty$, and let $\mathcal R \subseteq [s_a, s_b)$ be closed relative to $[s_a, s_b)$ with first hitting time $s_1 := \inf\{s \ge s_a : s \in \mathcal R\}$ ($\inf\emptyset := s_b$). If $x'(s) \le -\alpha s + r(s)$ for a.e.\ $s \in [s_a, \min\{s_1, s_b\})$ ($\alpha > 0$) and
\begin{equation}
  s_{\mathrm{route}} := s_a + \sqrt{\frac{2(x(s_a) + \eta)}{\alpha}} < s_b,
  \label{eq:s-route-general}
\end{equation}
then $s_1 \le s_{\mathrm{route}}$ and $s_1 \in \mathcal R$.
\item \emph{(Scalar barrier invariance.)} Let $\psi : [a, b) \to \R$ ($-\infty < a < b \le \infty$) be continuous on $[a, b)$ and differentiable on $(a, b)$. If $\psi(a) \ge 0$ and $\dot\psi(\tau) > 0$ at every $\tau \in (a, b)$ with $\psi(\tau) \le 0$, then $\psi(\tau) \ge 0$ for all $\tau \in [a, b)$.
\end{enumerate}
\end{lemma}
\begin{proof}[Proof of \cref{lem:hit-compare}]
For~(i), suppose $s_1 > s_{\mathrm{route}}$ and fix $s \in (s_{\mathrm{route}}, \min\{s_1, s_b\})$. Since $s > s_a \ge 0$, we have $s^2 - s_a^2 \ge (s - s_a)^2 > \frac{2(x(s_a)+\eta)}{\alpha}$. Integrating $x'(u) \le -\alpha u + r(u)$ over $[s_a, s]$ gives
\[
  x(s) \le x(s_a) - \frac{\alpha}{2}\bigl(s^2 - s_a^2\bigr) + \eta < 0,
\]
contradicting $x(s) \ge 0$; thus $s_1 \le s_{\mathrm{route}} < s_b$, and relative closedness of $\mathcal R$ gives $s_1 \in \mathcal R$.

For~(ii), if $\psi(\tau_2) < 0$ for some $\tau_2 \in (a, b)$, let $\tau_\star := \sup\{\tau \in [a, \tau_2] : \psi(\tau) \ge 0\} \in [a, \tau_2)$. Continuity gives $\psi(\tau_\star) = 0$ and $\psi < 0$ on $(\tau_\star, \tau_2]$, so by the mean value theorem there exists $\xi \in (\tau_\star, \tau_2) \subseteq (a, b)$ with $\psi(\xi) < 0$ and $\dot\psi(\xi) = \psi(\tau_2)/(\tau_2 - \tau_\star) < 0$, a contradiction.
\end{proof}

\subsection{Routing and Rollout Bounds}
\label{app:rms-struct-sub}
\label{app:rms-routing-sub}

In normalized time $s = \initscale\tau$, \cref{lem:rms-align} proves gate saturation with an integrable non-letter remainder, \cref{lem:rms-contrast} proves an unconditional $2^{-13}$ floor on distractor reciprocal weights, and \cref{lem:rms-routing-time} bounds the routing hitting time. \Cref{thm:block-rms} then follows in five short steps.

\begin{lemma}[Gate saturation and integrable non-letter perturbation]
\label{lem:rms-align}
Under \cref{def:gate} with $d_k \ge 2n + 6$, $a_0 > 0$, and $\initscale \le \frac{1}{8}$, evolve $(Q_h, k_h, w_h)$ by Block-RMS flow and $(U, \brel^{(h)})$ by Euclidean flow along the unique symmetric Carath\'eodory trajectory of \cref{lem:rms-wellposed}, and set $s_0 := \initscale d_k^{-1/2}$, $s_a := s_0 + 4 d_k^{-1/4}\sqrt{\log(8n)}$, and $\eta := (8n)^{-7}$. On any interval $(0, T)$ ($s_b := \initscale T$) where $-\pd{\Lstep}{g} > 0$ and $-\pd{\Lstep}{a} > 0$:
\begin{enumerate}[label=(\roman*),leftmargin=*,itemsep=1pt,topsep=2pt]
\item $g(s/\initscale) = v_0\initscale s \in [\initscale s,\, 4\initscale s]$ and $|\brel_\delta(s/\initscale)| \le \frac{1}{2}s^2$ on $(0, s_b)$, $m(s) > 2^{-12}$ on $[0, \min\{s_a, s_b\}]$, and $a(s/\initscale) \ge \frac{15}{16}\sqrt{d_k}\,s^2$ on $[s_a, s_b)$;
\item for all $s \in [s_a, s_b)$, $1 - D_t(s/\initscale), |\xi_t(s/\initscale)| \le 2n\,e^{-\frac{1}{2}\sqrt{d_k}\,s^2} \le 2(8n)^{-4}$; and if $Z_t > \frac{1}{2}$ on $(0, T)$, then $r(s) := \frac{1}{\initscale}g(s/\initscale)\max_{2\le t\le n}|\xi_t(s/\initscale)|$ bounds every non-letter remainder $\frac{1}{\initscale}|R_\delta(s/\initscale)| \le r(s)$ in \cref{eq:U-wd-ode} with
\[
  \int_{s_a}^{s_b} r(s)\,\rmd s \le \frac{8n}{\sqrt{d_k}}\,e^{-\frac{1}{2}\sqrt{d_k}\,s_a^2} < \eta.
\]
\end{enumerate}
\end{lemma}
\begin{proof}[Proof of \cref{lem:rms-align}]
Dividing the symmetric Euclidean gradients (\cref{lem:J-reduce,lem:rms-wellposed}) by their RMS norms on $(0, T)$ gives $\dot g = v_0\initscale^2$, $\dot q = \sqrt{2d_k}\,k/\|k\|_2$, and $\dot k = \sqrt{d_k}\,q/\|q\|_2$. With $u := \|q\|_2 \le 2\initscale + \sqrt{2d_k}\tau$, $w := \|k\|_2 \le 2\initscale + \sqrt{d_k}\tau$, and $\cos\psi := \frac{\langle q,k\rangle}{uw} \in (0,1]$,
\[
  \ddtau{\cos\psi} = \Bigl(\frac{\sqrt{2d_k}}{u} + \frac{\sqrt{d_k}}{w}\Bigr)(1 - \cos^2\psi) \ge \frac{2}{\tau + 2\initscale/\sqrt{d_k}}(1 - \cos^2\psi)
  \implies \cos\psi(\tau) \ge \frac{X^4 - 1}{X^4 + 1},
\]
where $X := 1 + \frac{\sqrt{d_k}\,s}{2\initscale^2} \ge 5$ for $s \ge s_0$ ($\cos\psi > 0.99$ independently of $r_{\mathrm{align}}$).

Integrating $\dot u \ge 0.99\sqrt{2d_k}$ and $\dot w \ge 0.99\sqrt{d_k}$ from $s_0/\initscale$ gives $a(s/\initscale) \ge 0.99^3\sqrt{d_k}(s - s_0)^2 \ge \frac{15}{16}\sqrt{d_k}\,s^2$ on $[s_a, s_b)$ ($s_0/s_a \le \frac{1}{94}$ for $d_k \ge 10$). Meanwhile \cref{lem:gate-calculus}(ii) gives $|\pd{M_t}{\brel_\delta}| \le \frac{1}{2}g$, so $|\brel_\delta'(s)| \le \frac{1}{4\initscale}g \le s$ integrates to $|\brel_\delta(s/\initscale)| \le \frac{1}{2}s^2$; since $\log(8n) \le \sqrt{2n+6} \le \sqrt{d_k}$ ($n \ge 2$) implies $\frac{1}{2}s_a^2 < 8.17 < 12\log 2$, we obtain $m(s) > 2^{-12}$ up to $s_a$, proving~(i).

For~(ii), $|u_t| \le \log n + s^2$ on $[s_a, s_b)$ gives $a - |u_t| \ge \frac{1}{2}\sqrt{d_k}\,s^2 - \log n$, so \cref{lem:gate-calculus}(i) bounds $1 - D_t$ and $|\xi_t|$ by $2n e^{-\frac{1}{2}\sqrt{d_k}\,s^2} \le 2n(8n)^{-8} \le 2(8n)^{-4}$. When $Z_t > \frac{1}{2}$, $|R_\delta| \le g\max_t|\xi_t|$ gives
\[
  \frac{1}{\initscale}\bigl|R_\delta(s/\initscale)\bigr| \le r(s) \le 8n s\,e^{-\frac{1}{2}\sqrt{d_k}\,s^2},
\]
whose integral over $[s_a, s_b)$ is at most $\frac{8n}{\sqrt{d_k}}e^{-\frac{1}{2}\sqrt{d_k}\,s_a^2} \le \frac{8n}{\sqrt{10}}(8n)^{-8} < \eta$.
\end{proof}

Because the non-letter perturbation $r(s)$ is integrable after $s_a$, combining \cref{lem:inv-pert} with a prior-independent symmetry of within-letter mismatch events yields an unconditional $2^{-13}$ floor on every distractor reciprocal weight.

\begin{lemma}[Stable distractor weights via prior-independent contrasts]
\label{lem:rms-contrast}
Write $E_{t,r} := \E[\sigma(-M_t)\chi_t\chi_r]$, $A_r := \{\chi_r \ne \chi_t\}$, and $\E[X; A] := \E[X\one_A]$.
\begin{enumerate}[label=(\roman*),leftmargin=*,itemsep=1pt,topsep=2pt]
\item \emph{(Sign structure.)} Whenever $\gamma_t \ge 0$, $E_{t,t} - \mu_t = 2\sum_{r\ne t}\pi_{t,r}\E[\sigma(-M_t); A_r] \ge 0$ ($P_n \ge 0$), $\lvert\E[\sigma(-M_t); A_r] - p(1-p)\rvert \le \frac{1}{2}\gamma_t p(1-p)$ for every $r \ne t$ and $\cnt \in \R$, and $E_{t,r} - \mu_t \le -2\pi_t^\star\E[\sigma(-M_t); A_r] \le 0$ ($P_\delta \le 0$) whenever $\delta \in \arg\min_{\mathcal D} y_{\delta'}$ and $\delta \in \mathcal W_t$ ($r = n + t - \delta$).
\item \emph{(Unconditional distractor floor.)} Under \cref{lem:rms-align} on any interval $(0, T)$ where $\mu_t > 0$ for all $t$, we have $x(s) < 2$, $Z_t(s/\initscale) > \frac{1}{2}$, \cref{eq:rms-floor-shift} holds for all $s_a \le s_1 \le s < s_b$, and
\[
  \boxed{\;m(s) \ge m(s_a) - \eta > 2^{-13} =: f \qquad \forall s \in (0, s_b).\;}
\]
\end{enumerate}
\end{lemma}
\begin{proof}[Proof of \cref{lem:rms-contrast}]
For~(i), writing $\chi_t\chi_r = 1 - 2\one_{A_r}$ ($A_t = \emptyset$) gives $E_{t,r} = \E[\sigma(-M_t)] - 2\E[\sigma(-M_t); A_r]$ and hence $E_{t,t} - \mu_t = 2\sum_{r\ne t}\pi_{t,r}\E[\sigma(-M_t); A_r] \ge 0$. For any $r \ne t$ and $\cnt \in \R$, direct evaluation gives the prior-independent identity
\[
  \E\bigl[\sigma(-\chi_t\cnt); A_r\bigr] = (1-p)p\,\sigma(-\cnt) + p(1-p)\,\sigma(\cnt) = p(1-p),
\]
and integrating $|\sigma(-M_t) - \sigma(-\chi_t\cnt)| \le \frac{1}{4}\gamma_t$ over $\Pr(A_r) = 2p(1-p)$ yields $|\E[\sigma(-M_t); A_r] - p(1-p)| \le \frac{1}{2}\gamma_t p(1-p)$.

Next, for distinct $r, r' \ne t$ with $\pi_{t,r} \ge \pi_{t,r'}$, swapping $\chi_r \leftrightarrow \chi_{r'}$ preserves the prompt law, exchanges $A_r\setminus A_{r'} \leftrightarrow A_{r'}\setminus A_r$, and increases $M_t$ on $A_r\setminus A_{r'}$ by $2\gamma_t(\pi_{t,r} - \pi_{t,r'}) \ge 0$, so $\E[\sigma(-M_t); A_r] \ge \E[\sigma(-M_t); A_{r'}]$. At any active minimizer $\delta \in \arg\min_{\mathcal D} y_{\delta'}$ with $\delta \in \mathcal W_t$, the index $r = n + t - \delta$ maximizes $\pi_{t,r'}$ over $r' \ne t$, giving $E_{t,r} - \mu_t \le -2\pi_t^\star\E[\sigma(-M_t); A_r] \le 0$.

For~(ii), \cref{eq:U-post-wnw1} gives $\omega_n + \omega_1 \le 2$ on $(0, T)$, so $x(s) < 2$ and $Z_t \ge x^{-1} > \frac{1}{2}$. Combining~(i) with \cref{lem:rms-align}(ii) gives $x'(s) \le r(s)$ and $y_\delta'(s) \ge -r(s)$ on $\arg\min_{\mathcal D} y_{\delta'}(s)$, so \cref{lem:inv-pert}(i) yields \cref{eq:rms-floor-shift} and $m(s) \ge 2^{-12} - \eta > 2^{-13}$.
\end{proof}

With $m(s) > 2^{-13}$ secured unconditionally, the target-lag reciprocal weight $x(s) = \omega_n(s/\initscale)$ obeys $x'(s) \le -\alpha s + r(s)$ with $\alpha \asymp p(1-p)/n^3$ and crosses into the routed region $\{8(n-1)x(s) \le m(s)\}$ in polynomial normalized time (proved in \cref{app:rms-tech-sub}).

\begin{lemma}[Positivity bootstrap and finite-time routing]
\label{lem:rms-routing-time}
Under \cref{lem:rms-align} with $p \in (0, 1/2]$, $a_0 > 0$, $d_k \ge 2n + 6$, and $0 < \initscale \le 2^{-30}\frac{[p(1-p)]^{3/2}}{n^3\sqrt{\log(8n)}}$, let $T := \sup\{\tau > 0 : \mu_t(u) > 0 \;\forall t \in \{2,\ldots,n\},\, u \in [0, \tau)\}$ and $s_{\mathrm{rt}} := s_a + \frac{2^{13}n^{3/2}}{\sqrt{p(1-p)}}$ ($\tau_{\mathrm{rt}} := s_{\mathrm{rt}}/\initscale$). Then the trajectory enters $\{8(n-1)x(s) \le m(s)\}$ at a time $s_1 \le s_{\mathrm{rt}} < \frac{p(1-p)}{2^{15}n\initscale}$ ($\tau_1 := s_1/\initscale < \initscale^{-2}$) with $-\log 2 \le \cnt(\tau_1) \le c^\star + \log 2$; thereafter $T = \infty$, and for all $\tau \ge \tau_1$ every step satisfies $\chi_t\rho_t(\tau) \ge \frac{3}{4}$ on every prompt and
\[
  \boxed{\;\min_{2\le t\le n}\pi_t^\star(\tau) \ge \frac{7}{8}.\;}
\]
\end{lemma}

\begin{remark}[Limiting usage of the $\initscale \lesssim n^{-3}(\log n)^{-1/2}$ scale under Block-RMS]
\label{rem:rms-eps-role}
In normalized time $s = \initscale\tau$, the target-lag drift rate $\alpha \asymp p(1-p)/n^3$ in \cref{eq:rms-Pn-lower} routes attention by $s_{\mathrm{rt}} \asymp n^{3/2}/\sqrt{p(1-p)}$. Meanwhile, the readout gain grows linearly as $g(s/\initscale) = v_0\initscale s$; keeping $g(s_{\mathrm{rt}}/\initscale) \ll \frac{p(1-p)}{n}$ throughout $[0, s_{\mathrm{rt}}]$ so the uncalibrated gate cannot flip $\mu_t > 0$ before routing completes (\cref{lem:rms-routing-time}, Step~1) requires $\initscale s_{\mathrm{rt}} \lesssim \frac{p(1-p)}{n}$, i.e.\ $\initscale \lesssim [p(1-p)]^{3/2}n^{-5/2}$ (sharpened by $(n\log(8n))^{-1/2}$ to absorb $s_a$).
\end{remark}

\par\bigskip\noindent\textbf{Theorem~\ref{thm:block-rms} (Unconditional rollout via stateless Block-RMS; formal statement).}\enspace{\itshape
In the structured Block-$2$ circuit of \cref{def:gate} with precomputed match signs and frozen values (without Hypothesis~(F)), for any $n \ge 2$, $p \in (0,1)$, $d_k \ge 2n + 6$, $\langle q(0), k(0)\rangle \ne 0$, and $0 < \initscale \le c_0\frac{[p(1-p)]^{3/2}}{n^3\sqrt{\log(8n)}}$, evolving $B \in \{Q_1, Q_2, k_1, k_2, w_1, w_2\}$ by stateless Block-RMS flow $\ddtau{\theta_B} = -\nabla_{\theta_B}\Lstep / \mathrm{RMS}(\nabla_{\theta_B}\Lstep)$ and $(U, \brel^{(h)})$ by Euclidean flow (\cref{lem:rms-wellposed}) routes attention by $\tau_{\mathrm{rt}} \le C n^{3/2}[p(1-p)]^{-1/2}\initscale^{-1} < \initscale^{-2}$ and achieves $\min_{2\le t\le n}\pi_t^\star(\tau) \ge \frac{7}{8}$ and $\min_{2\le t\le n,\,w\in\Sigma^n} M_t(\tau) \ge \frac{1}{2}$ (persistent exact greedy rollout from $s_1$, with $\lim_{\tau\to\infty}\Lstep(\tau) = 0$ and $\lim_{\tau\to\infty}\pi_t^\star(\tau) \in [\frac{7}{8}, 1)$) for all $\tau \ge T_{\mathrm{RMS}}$, with $T_{\mathrm{RMS}} \le \frac{7(1 + |c^\star|)}{\initscale^2}$ as in \cref{eq:rms-time}.}\par\smallskip

\begin{proof}[Proof of \cref{thm:block-rms}]
Take $c_0 := 2^{-30}$ and $C := 7$, and assume without loss of generality that $a_0 > 0$ and $p \le \frac{1}{2}$ ($c^\star \ge 0$) by the head-sign and label-flip symmetries of \cref{lem:rms-wellposed}.

By \cref{lem:rms-align,lem:rms-contrast}, gate saturation makes the non-letter perturbation integrable ($\int_{s_a}^{s_b} r(s)\,\rmd s < \eta$) and keeps all distractor reciprocal weights bounded below by $m(s) > 2^{-13}$.

\Cref{lem:rms-routing-time} (proved in \cref{app:rms-tech-sub}) then establishes that $\min_{2\le t\le n}\pi_t^\star(\tau)$ reaches $\frac{7}{8}$ by $\tau_1 \le \tau_{\mathrm{rt}} < \initscale^{-2}$ and remains $\ge \frac{7}{8}$ for all $\tau \ge \tau_1$, with $T = \infty$, $\chi_t\rho_t(\tau) \ge \frac{3}{4}$ on every prompt, and $-\log 2 \le \cnt(\tau_1) \le c^\star + \log 2$.

On $[\tau_1, \infty)$, \cref{lem:rms-align}(ii) gives $D_t \ge 1 - 2(8n)^{-4} \ge \frac{15}{16}$, and applying \cref{lem:hit-compare}(ii) via \cref{lem:rms-margin-trap} traps $-\bigl(\frac{1}{4}g(\tau) + 2\bigr) \le \cnt(\tau) \le \frac{1}{4}g(\tau) + c^\star + 2$. Since $\dot g \ge \initscale^2$, every prompt satisfies
\[
  M_t(\tau) \ge \frac{3}{4}\gamma_t(\tau) - |\cnt(\tau)| \ge \frac{3}{8}g(\tau) - c^\star - 2 \ge \frac{1}{2}
  \qquad\text{for all }\;
  \tau \ge T_{\mathrm{RMS}} := \max\Bigl\{\tau_1,\; \frac{3c^\star + 7}{\initscale^2}\Bigr\} \le \frac{7(1 + |c^\star|)}{\initscale^2}.
\]
Therefore, greedy rollout from $s_1$ is exact on every prompt after $T_{\mathrm{RMS}}$.

Finally, linear growth $g(\tau) = v_0\initscale^2\tau \to \infty$ sends $M_t(\tau) \to \infty$ and $\Lstep(\tau) \to 0$, while $|\dot\brel_\delta(\tau)| \le \frac{1}{4}g(\tau)\,e^{-\frac{3}{8}g(\tau) + |c^\star| + 2} \in L^1([T_{\mathrm{RMS}}, \infty))$ implies convergence $\brel(\tau) \to \brel(\infty) \in \R^{2n+2}$, so $\lim_{\tau\to\infty}\pi_t^\star(\tau) = \pi_t^\star(\infty) \in [\frac{7}{8}, 1)$.
\end{proof}

\subsection{Well-Posedness and Constants}
\label{app:rms-tech-sub}

This subsection collects the technical verifications cited in \cref{app:rms-routing-sub}: Carath\'eodory well-posedness and sign symmetries of the nine-block Block-RMS flow (\cref{lem:rms-wellposed}), the three-step proof of \cref{lem:rms-routing-time}, and the post-routing readout barrier check (\cref{lem:rms-margin-trap}).

\begin{lemma}[Carath\'eodory well-posedness, ridge selection, and sign symmetries]
\label{lem:rms-wellposed}
Under \cref{def:gate} with $\langle q(0), k(0)\rangle \ne 0$:
\begin{enumerate}[label=(\roman*),leftmargin=*,itemsep=1pt,topsep=2pt]
\item Replacing $(k_h, w_h) \mapsto (-k_h, -w_h)$ maps $(g, a) \mapsto (-g, -a)$ and $(p, \chi, U, w_h) \mapsto (1-p, -\chi, RUR, Rw_h)$ maps $(p, c^\star, \cnt) \mapsto (1-p, -c^\star, -\cnt)$ while preserving the Block-RMS equations and rollout guarantees, reducing to $a_0 > 0$ and $p \in (0, 1/2]$.
\item Because $w_h(0) = 0$ makes $\nabla_{Q_h}\Lstep(0) = 0$ and $\nabla_{k_h}\Lstep(0) = 0$ at the single instant $\tau = 0$, we take locally absolutely continuous Carath\'eodory solutions ($0/\mathrm{RMS}(0) := 0$) selected as the $\iota \downarrow 0$ limit of the smooth $\mathcal T$-equivariant ridge-regularized flows $\mathrm{RMS}_\iota(G) := (\|G\|_{\frob}^2/d_B + \iota^2)^{1/2}$ (each uniquely trapped on \cref{eq:U-manifold} by Picard--Lindel\"of). Directly at $\iota = 0$ on \cref{eq:U-manifold}, $a_0 > 0$ gives $-\pd{\Lstep}{g}(0) > 0$, so every Carath\'eodory solution satisfies $g(\tau) = v_0\initscale^2\tau > 0$ on $(0, \delta_0)$; for $\tau > 0$, $g(\tau)$ cancels out of $\nabla_{Q_h}\Lstep/\mathrm{RMS}(\nabla_{Q_h}\Lstep)$ and $\nabla_{k_h}\Lstep/\mathrm{RMS}(\nabla_{k_h}\Lstep)$, leaving a $C^\infty$ locally Lipschitz ODE on $\{q \ne 0, k \ne 0\}$ with a unique trajectory satisfying the $9$-block Block-RMS ODE at every $\tau \in (0, T)$.
\end{enumerate}
\end{lemma}
\begin{proof}[Proof of \cref{lem:rms-wellposed}]
For~(i), orthogonal reflections preserve Frobenius/RMS norms, and $M_t = \chi_t(\cnt + g D_t\rho_t)$ is invariant under $(g, a) \mapsto (-g, -a)$ ($D_t(-a) = -D_t(a)$) and under $(\chi, \cnt) \mapsto (-\chi, -\cnt)$.

For~(ii), $\mathcal T$ (\cref{lem:J-symmetry}) swaps paired blocks of equal dimension ($d_{Q_h} = 2d_k, d_{k_h} = d_k, d_{w_h} = 2$), keeping $\mathrm{RMS}_\iota(\nabla_{\theta_B}\Lstep)$ identical across paired blocks for every $\iota \ge 0$. At $\tau = 0$, $D_t(0) > 0$ and $\mu_t(0) = \frac{1}{2}(\frac{1}{n} + \frac{n-1}{n}\mixrate^2) \ge \frac{1}{2n} > 0$ give $-\pd{\Lstep}{g}(0) > 0$.
\end{proof}

We now complete the proof of \cref{lem:rms-routing-time} by combining the positivity bootstrap on $[0, s_g)$ with the hitting-time and ratio-persistence comparison lemmas of \cref{app:rms-compare-sub}.

\begin{proof}[Proof of \cref{lem:rms-routing-time}]
Set $\alpha := \frac{p(1-p)}{2^{22}n^3}$ and $s_g := \frac{p(1-p)}{2^{15}n\initscale}$ ($\tau_g := s_g/\initscale$, $s_b := \initscale T$).

\proofstep{Step 1: Prior residual and positivity bootstrap to $s_g$.}
By \cref{lem:rms-wellposed}(ii), $s_b > 0$. Combining $|\sigma(-M_t) - \sigma(-\chi_t\cnt)| \le \frac{1}{4}\gamma_t$ and $\E[\sigma(-\chi_t\cnt)] = 2p(1-p) + (1-2p)(\sigma(-\cnt) - p)$ with \cref{lem:rms-contrast}(i) gives
\[
  \bigl|\E[\sigma(-M_t)\chi_t] - (\sigma(-\cnt) - p)\bigr| \le \frac{1}{4}\gamma_t,
  \qquad
  \mu_t \ge 2p(1-p)\pi_t^\star + (1-2p)\bigl(\sigma(-\cnt) - p\bigr) - \frac{1}{2}\gamma_t.
\]
On $[0, \min\{T, \tau_g\})$ ($g < \frac{p(1-p)}{2^{13}n} < 2^{-15}$), applying \cref{lem:hit-compare}(ii) to $\sigma(-\cnt) - p + \frac{1}{4}g$ and $\cnt + \log 2$ preserves $\sigma(-\cnt) - p \ge -\frac{1}{4}g$ and $-\log 2 \le \cnt \le c^\star + \log 2$. Since \cref{lem:rms-contrast}(ii) gives $x(s) < 2$ and $m(s) > 2^{-13}$, we obtain
\[
  \pi_t^\star > \frac{1}{2^{14}n}
  \qquad\text{and}\qquad
  \mu_t \ge 2p(1-p)\pi_t^\star - \frac{3}{4}g > \frac{p(1-p)}{2^{15}n} > 0
\]
up to $\min\{T, \tau_g\}$, forcing $s_b \ge s_g$. Moreover, $s_a < 4\sqrt{\log(8n)} < \frac{2^{12}n^{3/2}}{\sqrt{p(1-p)}}$ and $\initscale \le 2^{-30}\frac{[p(1-p)]^{3/2}}{n^3\sqrt{\log(8n)}}$ imply $s_a < s_{\mathrm{rt}} < \frac{2^{14}n^{3/2}}{\sqrt{p(1-p)}} < s_g < \initscale^{-1}$.

\proofstep{Step 2: Target-lag drift and hitting time $s_1 \le s_{\mathrm{rt}}$.}
Let $\mathcal R := \{s \in [s_a, s_b) : 8(n-1)x(s) \le m(s)\}$ and $s_1 := \inf\mathcal R$. Fix $s \in [s_a, \min\{s_1, s_g\})$ ($8(n-1)x(s) > m(s)$), pick $\delta_\star \in \arg\min_{\mathcal D} y_\delta(s)$, and set $t := 2$ if $\delta_\star \le n - 1$ or $t := n$ if $\delta_\star \ge n + 1$, so $\delta_\star \in \mathcal W_t$ ($r_\star := n + t - \delta_\star \ne t$). Then
\[
  \sum_{r\ne t}\pi_{t,r} \ge \pi_{t,r_\star} = \bigl(m(s)Z_t\bigr)^{-1}, \qquad
  Z_t = x(s)^{-1} + \sum_{r\ne t} y_{n+t-r}(s)^{-1} < \frac{9(n-1)}{m(s)},
\]
with $\gamma_t/\initscale \ge \frac{1}{2}s$ ($D_t > \frac{1}{2}$) and $\E[\sigma(-M_t); A_r] \ge \frac{1}{2}p(1-p)$ (\cref{lem:rms-contrast}(i), $\gamma_t \le g < 1$). Substituting into the step-$t$ summand of $P_n$ with $m(s) > 2^{-13}$ (and $324\cdot 2^{13} = 2{,}654{,}208 < 2^{22}$) yields
\begin{equation}
  \frac{1}{\initscale}P_n(s/\initscale)
  \ge \frac{(\gamma_t/\initscale)\,p(1-p)}{2(n-1)\,m(s)\,Z_t^2}
  > \frac{s\,p(1-p)\,m(s)}{324(n-1)^3}
  > \alpha s,
  \label{eq:rms-Pn-lower}
\end{equation}
so $x'(s) \le -\alpha s + r(s)$. Since $2(x(s_a) + \eta) < 2(2 + \eta) < 8$ gives $s_{\mathrm{route}} = s_a + \sqrt{2(x(s_a)+\eta)/\alpha} < s_{\mathrm{rt}} < s_g \le s_b$, \cref{lem:hit-compare}(i) yields $s_1 \le s_{\mathrm{rt}}$ with $8(n-1)x(s_1) \le m(s_1)$.

\proofstep{Step 3: Permanent routing and $T = \infty$.}
Applying \cref{lem:inv-pert}(ii) on $[s_1, s_b)$ with $K = 8(n-1)$ and $m(s_1) > f = 2^{-13}$ gives
\[
  (n-1)\,\frac{x(s)}{m(s)} \le \frac{1/8 + (n-1)\eta/f}{1 - \eta/f} < \frac{1}{7},
\]
keeping $\pi_t^\star(\tau) \ge \frac{7}{8}$ on $[\tau_1, T)$; hence $\chi_t\rho_t(\tau) \ge 2\pi_t^\star(\tau) - 1 \ge \frac{3}{4}$ on every prompt and $\mu_t(\tau) \ge \frac{3}{4}\E[\sigma(-M_t(\tau))] > 0$, forcing $T = \infty$.
\end{proof}

Finally, we verify the boundary sign conditions used in the proof of \cref{thm:block-rms} to trap $|\cnt(\tau)| \le \frac{1}{4}g(\tau) + |c^\star| + 2$ after routing completes.

\begin{lemma}[Post-routing readout boundary signs]
\label{lem:rms-margin-trap}
When $p \in (0, 1/2]$, $g \ge 0$, $D_t \ge \frac{15}{16}$, and $\chi_t\rho_t \ge \frac{3}{4}$ on every prompt, we have $\dot\cnt < 0$ whenever $\cnt \ge \frac{1}{4}g + c^\star + 2$ and $\dot\cnt > 0$ whenever $\cnt \le -\bigl(\frac{1}{4}g + 2\bigr)$.
\end{lemma}
\begin{proof}[Proof of \cref{lem:rms-margin-trap}]
Here $\frac{3}{4}\gamma_t \ge \frac{45}{64}g > \frac{5}{8}g$. When $\cnt \ge \frac{1}{4}g + c^\star + 2$, we have $-M_t = \cnt - \gamma_t\chi_t\rho_t \ge c^\star + 2 - \frac{3}{4}g$ on $\{\chi_t = -1\}$ and $-M_t \le -(\frac{7}{8}g + c^\star + 2)$ on $\{\chi_t = +1\}$. Since
\[
  p\,\sigma\bigl(c^\star + 2 - \tfrac{3}{4}g\bigr)
  \ge \frac{1}{2}p\,e^{-\frac{3}{4}g}
  > p\,e^{-2 - \frac{7}{8}g}
  = (1-p)\,e^{-(\frac{7}{8}g + c^\star + 2)}
  > (1-p)\,\sigma\bigl(-(\tfrac{7}{8}g + c^\star + 2)\bigr),
\]
we obtain $\E[\sigma(-M_t)\chi_t] < 0$ and $\dot\cnt < 0$. Symmetrically, when $\cnt \le -(\frac{1}{4}g + 2)$, $(1-p)\sigma(2 - \frac{3}{4}g) > p\,e^{-(\frac{7}{8}g + 2)}$ gives $\dot\cnt > 0$.
\end{proof}

\paragraph{Absence of Hypothesis~(F) and freezing of $\brel$ under Block-RMS.}
Under Euclidean flow, routing finishes while $a \le \frac{1}{4}$ (\cref{lem:U-preroute}); in that regime $\xi_t = \Theta(1)$ lets non-letter drift act at first order, and after escape $\dot\brel_n / \dot g = \Theta(1)$ sends $\pi_t^\star(\tau) \to 1$. Under Block-RMS in normalized time $s = \initscale\tau$, $a$ grows as $\sqrt{d_k}\,s^2$ while $|\brel_\delta| \le \frac{1}{2}s^2$ (\cref{lem:rms-align}). Gate saturation makes $|\xi_t|$ Gaussian-integrable ($\int_{s_a}^\infty r(s)\,\rmd s < (8n)^{-7}$), so \cref{lem:inv-pert} keeps $m(s) > 2^{-13}$ without Hypothesis~(F). Once $\min_t\pi_t^\star \ge \frac{7}{8}$, linear growth of $g$ makes $\dot\brel$ integrable on $[\tau_1, \infty)$, so $\pi_t^\star(\infty) \in [\frac{7}{8}, 1)$ while $\Lstep(\tau) \to 0$.

% ===================================================================
\section{Experimental Setup and Interventions}
\label{app:experiments}
\label{app:experimental-details}
\label{app:letter-protocol}
\label{app:reproducibility}
\setcounter{totalnumber}{8}
\setcounter{topnumber}{6}
\setcounter{bottomnumber}{6}
\renewcommand{\textfraction}{0.05}
\renewcommand{\topfraction}{0.95}
\renewcommand{\bottomfraction}{0.95}

\subsection{Training and Evaluation Protocol}

All source code, configurations, seeds, and logs are available at \url{https://anonymous.4open.science/r/trace} (executed on an NVIDIA A100 80GB GPU in PyTorch and in standalone multithreaded C++20/BLAS). Paired step and outcome runs share identical initialization, data order, optimizer, and update budget (\cref{tab:protocol}; letter tasks use $\dmodel = 64$, $d_k = 24$, $\dmlp = 128$, $|\Sigma| = 4$ with $p = 1/4$ unless varied, AdamW $(\beta_1, \beta_2) = (0.9, 0.999)$ with weight decay $0.01$ and constant learning rate, supervising all continuation tokens after $\mathrm{SEP}$; \cref{fig:overview}a sets $V(0) = 0$ whereas \cref{tab:gap,tab:controls} use standard $\mathcal N(0, \initscale^2)$ on all $27$ blocks). Runs are evaluated by free-running greedy decoding on $1{,}000$--$2{,}048$ held-out prompts disjoint from the training set (\emph{Solved} $\ge 99.9\%$ exact rollout accuracy unless stated otherwise).

\begin{table}[!htbp]
\caption{\textbf{Shared experimental settings across model tiers.}}
\label{tab:protocol}
\centering\footnotesize
\fitwidth{\begin{tabular}{@{}llll@{}}
\toprule
& Letter tasks (two-block) & Boolean-8, register (GPT-2-2L) & LoRA (SmolLM2, Qwen2.5)\\
\midrule
Model & \cref{ass:arch}, $2$ blocks, $3$ heads & $2$ blocks, $4$ heads, width $128$, LayerNorm & rank $8$, $\alpha=16$, all linear layers\\
Data & $20{,}000$ train, $1{,}000$ test & $20{,}000$ train, $500$--$1{,}000$ test & $12{,}000$ train, $300$ test (Boolean-8)\\
Init / Opt & $\mathcal N(0,0.02^2)$; AdamW $\eta=10^{-3}$, bs $256$ & default; AdamW $\eta=3\times10^{-4}$, bs $128$ & default LoRA; AdamW $\eta=3\times10^{-4}$, bs $16$\\
Budget / Seeds & $8$k updates ($16$k/$32$k); seeds $42$--$51$ & $8$k updates; $3$--$5$ seeds & $800$ updates; seeds $2001$--$2003$\\
\bottomrule
\end{tabular}}
\end{table}

\paragraph{Plateau diagnostic $\hat c$ and stateless \texttt{Block-RMS SGD} (\cref{fig:overview}b,c, \cref{tab:transfer-main}a).}
On \cref{ass:arch} ($n=8, |\Sigma|=4, p=1/4$, Block-$2$ relative lag biases $\brel^{(2,h)}(0)=0$ and $V(0)=0$), \cref{fig:overview}b,c track seed~$42$ at $(\dmodel, d_k, \dmlp) = (64, 24, 128)$, $\initscale = 0.02$, AdamW $\eta = 10^{-3}$, batch size $256$, and $1{,}000$ held-out rollouts via the state-aligned logit proxy $\hat c := \frac{1}{n-1}\sum_{t=2}^n \E[x_{t-1}((z_t)_0 - (z_t)_1)]$. By \cref{lem:U-features}, the reduced circuit satisfies $\hat c = \cnt + (1-2p)\frac{1}{n-1}\sum_{t=2}^n g D_t = \cnt + O(\max_t\gamma_t)$, so $\hat c/c^\star \approx 1$ specifically diagnoses transition-prior calibration throughout the early weak-gate plateau before rollout accuracy rises.

In \cref{tab:transfer-main}a ($5$ paired seeds $42$--$46$, batch size $128$, $1{,}000$ held-out rollouts, \emph{Solved} $\ge 99.9\%$ exact rollout), tuned SGD uses $\eta=0.15$ with global gradient-norm clip $1.0$; \texttt{QKV-RMS + SGD} uses stateless block-RMS step size $2.5\times 10^{-3}$ on $(W_Q, W_K, W_V)$ and clipped SGD ($\eta=0.12$) elsewhere; \texttt{Block-RMS SGD} uses stateless block-RMS step size $2.0\times 10^{-3}$ across all $27$ weight/bias parameter blocks plus the $3$ Block-$2$ lag-bias vectors $\brel^{(2,h)}$ ($30$ active tensor slices in total, $0$ momentum, $0$ per-parameter state, updating each tensor via \texttt{rms = param.grad.pow(2).mean().sqrt() + 1e-8; param.add\_(param.grad, alpha = -eta / rms)}); and AdamW uses $\eta = 1.5\times 10^{-3}$ ($\dmodel=32, d_k=16, \dmlp=64, \initscale=0.05$) and $10^{-3}$ ($\dmodel=64, d_k=24, \dmlp=128, \initscale=0.02$). At $\initscale = 0.02$ ($B = 256$, $10{,}000$ updates, paired seeds $42$--$46$), the matched outcome \texttt{Block-RMS} run (\texttt{--mode 9}) uses the identical initial weights $\theta_0$, batch stream, test set, and $\eta = 2.0\times 10^{-3}$ over active outcome prefix positions $0,\ldots,n+2$ supervising $s_n$ at $\mathrm{SEP}$, reaching $1.0005$ bits loss ($1.0004$ at $2.4$k), $\E[p_\theta(s_n)] = 0.4999$, $49.12\%$ greedy accuracy, and $75.44\%$ Block-$1$ probe ($0/5$ solved) versus $0.0000$ bits loss, $100.00\%$ rollout, and $100.00\%$ probe ($5/5$ solved, median $1.20$k) under matched step \texttt{Block-RMS} (\texttt{--mode 8}). To build and reproduce \cref{fig:feature-learning,fig:causal-intervention,tab:clamp-main,tab:transfer-main}a from a single entry point: \texttt{make -C cpp all check \&\& ./cpp/run\_full\_plan --n 8 --p 0.25 --mode \{0..9\} --seed \{42..61\} --steps 2400 --outdir results/causal\_intervention/sweep/main \&\& python3 cpp/analysis/aggregate.py} (with \texttt{uv run python handcoded/experiments/sweeps.py zero\_readout} reproducing \cref{fig:overview}a).

\subsection{Baseline Intervention Details}
\label{app:causal-protocol}

Using a standalone multithreaded C++20 / BLAS training engine (finite-difference verified to maximum relative gradient error $2.43\times 10^{-7}$ across all parameter blocks, relative lag biases $\brel^{(2,h)}$, and projection hooks), we train the unmodified \cref{ass:arch} model at $n=8, |\Sigma|=4, (\dmodel, d_k, \dmlp)=(64, 24, 128), \initscale=0.02$ with AdamW ($\eta=10^{-3}, \beta_1=0.9, \beta_2=0.999$, weight decay $0.01$, batch size $B=256$, $2{,}400$ updates, $20$ seeds $42$--$61$ per condition, and $1{,}000$ held-out free-running rollouts per evaluation at updates $0, 25, 50, 100, 200, 300, 400, 600, 800, 1{,}000, 1{,}200, 1{,}600, 2{,}000, 2{,}400$; \cref{fig:causal-intervention} plots fine-cadence $100$-update trajectories on seeds $42$--$46$).

Using the compact step prefix $(s_0, w^\star, w_1, \ldots, w_n, \mathrm{SEP}, s_1, \ldots, s_{n-1})$ of length $2n+2$ (omitting the redundant terminal answer copy and supervising all $t \in \{1,\ldots,n\}$, so positional embeddings and attention also learn the fixed $t=1$ lookback to $s_0$ while summary diagnostics $\hat c, E_{\mathrm{true}}, E_{\mathrm{dist}}$ average over $t \in \{2,\ldots,n\}$), at each supervised position $i_t = n+1+t$ ($t \in \{1,\ldots,n\}$) let $\Delta z_t(m) := (z_{i_t}(m))_0 - (z_{i_t}(m))_1$ denote the binary state logit difference on minibatch prompt $m \in [B]$. Under \emph{Suppressed} ($\Pi_0$), we apply the exact minibatch orthogonal projection
\[
  \tilde z_{i_t,s}(m) \;=\; z_{i_t,s}(m) \,-\, \frac{1}{2}(1-2s)\,x_{t-1}(m)\,\widehat c_{t,B}, \qquad
  \widehat c_{t,B} := \frac{1}{B}\sum_{m'=1}^B x_{t-1}(m')\,\Delta z_t(m') \quad (s \in \{0,1\}),
\]
differentiating through the linear projection map during backpropagation so gradients live strictly in the subspace $\widehat c_{t,B} = 0$ (plus light nullspace damping $0.1\,\widehat c_{t,B}$ on the orthogonal component to prevent drift). Because $\E_{\mathcal D_p}[x_{t-1}\chi_t] = 1-2p$, the centered parity predictor $\Delta z_t = M x_{t-1}(\chi_t - (1-2p))$ satisfies $\E_{\mathcal D_p}[x_{t-1}\Delta z_t] = 0$ identically while achieving arbitrarily large positive margins $2Mp$ (mismatch) and $2M(1-p)$ (match) as $M\to\infty$, so $\Pi_0$ removes \emph{only} the state-only transition baseline without restricting parity expressivity.

We evaluate the following companion interventions:
\begin{itemize}[leftmargin=*]
\item \emph{Weight-Calibrated} ($V(0)=c^\star$): sets rows $0,1$ of $V(0)$ to $\pm \frac{c^\star}{2\|v_{\mathrm{state}}\|_2^2} v_{\mathrm{state}}^\top$ along $v_{\mathrm{state}} := \frac{1}{2}(E_{\mathrm{tok}}[0] - E_{\mathrm{tok}}[1])$ so $(V(0)h_{i_t}^{(0)})_0 - (V(0)h_{i_t}^{(0)})_1 = c^\star x_{t-1}$ at update~$0$, with unconstrained updates thereafter.
\item \emph{External Offset} ($+c^\star$): leaves $V(0)=0$ and adds $\pm \frac{1}{2}c^\star x_{t-1}$ externally to $z_{i_t,0}, z_{i_t,1}$ (along with three diagnostic ablations: removing the offset at update $300$; drawing rows of $V(0)$ as $\pm v$ with $v \sim \mathcal N(0, 0.15^2 I)$ projected orthogonal to $v_{\mathrm{state}}$ to match natural $\|V\|$; and multiplying the learning rate of $V$ by $10$).
\item \emph{Restored} and \emph{\texttt{Block-RMS SGD}}: applies $\Pi_0$ for $t < T_{\mathrm{rel}} \in \{600, 1{,}000\}$ before training freely thereafter, while \texttt{Block-RMS SGD} uses step size $2\times 10^{-3}$.
\end{itemize}

A ridge-regularized ($10^{-2}$) linear probe, fit on half of the held-out prompts and scored on the other half, predicts $b_r$ from the Block-$1$ output at $1+r$. Writing $\delta y_t := (\hat y_{i_t} - y_t)_0 - (\hat y_{i_t} - y_t)_1$ for the signed probability residual at step $t$, we log held-out rollout accuracy, Block-$2$ true-lag prefix share, and the residual-match correlations
\begin{align*}
  E_{\mathrm{true}} &:= \frac{1}{n-1}\sum_{t=2}^n \Bigl|\E\bigl[\delta y_t\,x_{t-1}\chi_t\bigr]\Bigr|, &
  E_{\mathrm{dist}} &:= \frac{1}{(n-1)^2}\sum_{t=2}^n\sum_{r\ne t}\Bigl|\E\bigl[\delta y_t\,x_{t-1}\chi_r\bigr]\Bigr|,
\end{align*}
which satisfy $E_{\mathrm{dist}}(0) = \frac{1}{2}(1-2p)^2$ and $E_{\mathrm{dist}}(c^\star) = 0$ in population, with finite-sample floor $\approx 0.01$ at $1{,}000$ prompts. \Cref{tab:causal-shortcut} reports every condition across $p \in \{0.05, 0.10, 0.15, 0.25, 0.50, 0.75\}$ (sharing paired seeds $42$--$61$ across conditions, with marginal two-sided Fisher $p$-values summarizing solve counts); at extreme skew $p \in \{0.05, 0.10\}$ ($\mixrate^2 \in \{0.64, 0.81\} > 2\log 2 - 1$), calibrating $\cnt \to c^\star$ shrinks the true-lag residual variance $E_{\mathrm{true}}(c^\star) = 2p(1-p)$ relative to the uncalibrated $1/2$ scale, lengthening post-routing gate opening as $[p(1-p)]^{-1}$ (\cref{thm:joint,rem:shortcut-paradox}) beyond the fixed $2.4$k-step cap so Natural remains at $0/20$ with Block-$1$ probe accuracy near the $\max(p, 1-p) = 90$--$95\%$ majority baseline while $\Pi_0$ solves $3/20$ and $1/20$ seeds.

\begin{table}[!htbp]
\caption{\textbf{Full-model prior interventions, all conditions} ($n=8, |\Sigma|=4, \dmodel=64, \initscale=0.02$, $2.4$k updates, $20$ seeds each; solved means $\ge 99.9\%$ exact free-running rollout on $1{,}000$ held-out prompts; Median is over solved seeds of the first evaluation solved; Fisher $p$-values cited in the text are two-sided and unadjusted across tested arms; $E_{\mathrm{dist}}$ and $\hat c$ are seed means at $100$ updates; Probe is Block-$1$ match-bit probe accuracy at $2.4$k, majority baseline $\max(p, 1-p)$).}
\label{tab:causal-shortcut}
\centering\footnotesize
\fitwidth{\setlength{\tabcolsep}{4.0pt}\begin{tabular}{@{}llccccccr@{}}
\toprule
$p$ ($c^\star$) & Condition & Solved & Wilson $95\%$ & Median & $E_{\mathrm{dist}}$ ($0.1$k) & $\hat c$ ($0.1$k) & $\|V\|$ ($0.4$k) & Probe ($2.4$k)\\
\midrule
$0.25$ ($+1.10$) & Natural & $18/20$ & $[0.70, 0.97]$ & $1.2$k & $0.012$ & $1.10$ & $0.50$ & $94.0\%$\\
 & Prior suppressed ($\Pi_0$) & $11/20$ & $[0.34, 0.74]$ & $2.0$k & $0.125$ & $0.00$ & $0.06$ & $90.3\%$\\
 & Suppressed, released at $0.6$k & $19/20$ & $[0.76, 0.99]$ & $1.6$k & $0.125$ & $0.00$ & $0.06$ & $95.6\%$\\
 & Suppressed, released at $1.0$k & $15/20$ & $[0.53, 0.89]$ & $1.6$k & $0.125$ & $0.00$ & $0.06$ & $90.2\%$\\
 & Weight-calibrated $V(0)=c^\star$ & $16/20$ & $[0.58, 0.92]$ & $0.8$k & $0.011$ & $1.09$ & $7.23$ & $89.7\%$\\
 & External offset $+c^\star x_{t-1}$ & $13/20$ & $[0.43, 0.82]$ & $1.6$k & $0.011$ & $1.10$ & $0.07$ & $86.9\%$\\
 & \quad offset removed at $0.3$k & $19/20$ & $[0.76, 0.99]$ & $1.2$k & $0.011$ & $1.10$ & $0.54$ & $94.6\%$\\
 & \quad with $V(0) \ne 0$, $\perp$ state axis & $13/20$ & $[0.43, 0.82]$ & $1.6$k & $0.011$ & $1.10$ & $1.37$ & $86.8\%$\\
 & \quad with $10\times$ learning rate on $V$ & $12/20$ & $[0.39, 0.78]$ & $2.0$k & $0.011$ & $1.10$ & $0.24$ & $85.0\%$\\
 & Natural, \texttt{Block-RMS SGD} & $18/20$ & $[0.70, 0.97]$ & $1.6$k & $0.022$ & $1.06$ & $0.22$ & $98.6\%$\\
\midrule
$0.05$ ($+2.94$) & Natural & $0/20$ & $[0.00, 0.16]$ & --- & $0.009$ & $3.13$ & $0.24$ & $95.0\%$\\
 & Prior suppressed ($\Pi_0$) & $3/20$ & $[0.05, 0.36]$ & $1.6$k & $0.404$ & $0.00$ & $0.04$ & $96.0\%$\\
\midrule
$0.10$ ($+2.20$) & Natural & $0/20$ & $[0.00, 0.16]$ & --- & $0.011$ & $2.31$ & $0.31$ & $90.3\%$\\
 & Prior suppressed ($\Pi_0$) & $1/20$ & $[0.01, 0.24]$ & $1.6$k & $0.320$ & $0.00$ & $0.04$ & $91.5\%$\\
\midrule
$0.15$ ($+1.73$) & Natural & $8/20$ & $[0.22, 0.61]$ & $2.0$k & $0.012$ & $1.79$ & $0.37$ & $91.2\%$\\
 & Prior suppressed ($\Pi_0$) & $4/20$ & $[0.08, 0.42]$ & $2.0$k & $0.245$ & $0.00$ & $0.04$ & $88.6\%$\\
\midrule
$0.50$ ($0$) & Natural & $19/20$ & $[0.76, 0.99]$ & $1.2$k & $0.013$ & $-0.00$ & $1.37$ & $95.0\%$\\
 & Prior suppressed ($\Pi_0$) & $19/20$ & $[0.76, 0.99]$ & $1.2$k & $0.013$ & $0.00$ & $1.28$ & $94.4\%$\\
\midrule
$0.75$ ($-1.10$) & Natural & $15/20$ & $[0.53, 0.89]$ & $2.0$k & $0.012$ & $-1.12$ & $0.39$ & $97.7\%$\\
 & Prior suppressed ($\Pi_0$) & $10/20$ & $[0.30, 0.70]$ & $1.6$k & $0.122$ & $0.00$ & $1.90$ & $96.7\%$\\
\bottomrule
\end{tabular}}
\end{table}

% ===================================================================
\section{Additional Experimental Results}
\label{app:exp-other}

\subsection{Scaling and Robustness Sweeps}
\label{app:exp-regime}
\label{app:regime}
\label{app:exp-spectra}
\label{app:spectrum}
\label{app:validation-spectra}
\label{app:exp-plateau}
\label{app:tracking}
\label{app:exp-lambda}

Using two blocks of width $\dmodel = 16$, three heads of dimension $d_k = 8$, MLP width $\dmlp = 16$ with $|\tilde b_m| \ge 1$, and $p = 1/4$ ($\mixrate = 1/2$) in float64 population arithmetic over all $2^n$ words, \cref{tab:exit-times} confirms that outcome escape time grows as $\Omega(\mixrate^{-n})$ both at the certified radius of $\initball$ and at macroscopic radius $r = 0.1$ (with zero ReLU sign flips prior to exit), while step escape time is $O(1)$. On matched initializations ($\initscale = 0.02$), step rule credit scales as $\Theta(\initscale^6/n)$ while outcome interaction credit (Hoeffding--Efron--Stein ANOVA degrees $|S| \ge 2$ \citep{hoeffding1948class,efron1981jackknife,odonnell2014analysis}) scales as $O(\initscale^6\mixrate^{n-3})$, with over $99.6\%$ of initial logit-Jacobian energy sitting at ANOVA degree~$0$.

\begin{table}[!htbp]
\caption{\textbf{Escape time from $\initball$ ($\initscale=0.002$) and from radius $r=0.1$ ($\initscale=0.02$)} under exact population GD on count mod $2$ ($p=1/4$, $\mixrate=1/2$, $3$ seeds).}
\label{tab:exit-times}
\centering\footnotesize
\fitwidth{\begin{tabular}{@{}lcccccc@{}}
\toprule
Quantity & $n=4$ & $n=6$ & $n=8$ & $n=10$ & $n=12$ & $n=14$\\
\midrule
Outcome exit time from $\initball$ & $2.8$--$6.8\mathrm{e}2$ & $1.4$--$3.4\mathrm{e}3$ & $0.7$--$1.6\mathrm{e}4$ & $3.2$--$8.2\mathrm{e}4$ & $1.5$--$3.5\mathrm{e}5$ & $0.7$--$1.8\mathrm{e}6$\\
Step exit time from $\initball$ & $0.11$--$0.16$ & $0.11$--$0.16$ & $0.10$--$0.16$ & $0.10$--$0.16$ & $0.10$--$0.16$ & $0.10$--$0.16$\\
Outcome exit time ($r=0.1$) & $0.9$--$1.5\mathrm{e}4$ & $0.5$--$0.8\mathrm{e}5$ & $2.3$--$3.8\mathrm{e}5$ & $1.0$--$1.8\mathrm{e}6$ & $4.7$--$7.8\mathrm{e}6$ & $2.2$--$3.8\mathrm{e}7$\\
Step exit time ($r=0.1$) & $4.47$--$5.55$ & $4.33$--$5.38$ & $4.26$--$5.29$ & $4.22$--$5.25$ & $4.20$--$5.23$ & $4.18$--$5.22$\\
\bottomrule
\end{tabular}}
\end{table}

Under standard $\mathcal N(0, 0.02^2)$ initialization and AdamW, the initial outcome gradient decays by $\approx 0.508$ per letter (step-to-outcome norm ratio rising from $21\times$ at $n=2$ to $11{,}708\times$ at $n=10$), and the Jacobian energy share at degrees $\ge 2$ stays $\le 2.4\times 10^{-6}$ throughout the $6$k-update plateau before jumping five orders of magnitude to $0.155$ as the trajectory begins escaping at $8$k updates. Across a $20$-cell grid varying $|\Sigma| \in \{4,6,10,20\}$ ($\mixrate = 1 - 2/|\Sigma|$) and $n \in \{6,8,10,12,16\}$, outcome accuracy correlates with $\mixrate^n$ at Spearman $r_{\mathrm S} = 0.91$ (mean $+0.6$ percentage points above $(1+\mixrate^n)/2$, vs.\ $r_{\mathrm S} = -0.36$ with $n$). \Cref{tab:gap,tab:controls} report robustness across horizons, $20$ seeds, initialization scales, deeper models ($4$--$6$ blocks), and the first-occurrence control.

\begin{table}[!htbp]
\caption{\textbf{Extended horizon and $20$-seed sweeps on count mod $2$} ($\mathcal N(0, 0.02^2)$ init, held-out exact-continuation accuracy \%; solved $\ge 99.9\%$).}
\label{tab:gap}
\centering\footnotesize
\fitwidth{\begin{tabular}{@{}llccc@{}}
\toprule
Sweep & Runs & Step accuracy & Outcome accuracy & Chance\\
\midrule
Extended budget ($32$k updates), $n \in \{8,12,16,24\}$ & $12$ & $12/12$ solved within $14$k; $11/12$ at $100\%$ at $32$k & $46$--$74\%$ at $32$k ($0/12$ solved) & $50\%$\\
$20$ seeds per horizon ($8$k updates), $n \in \{6,8,10,12,14,16\}$ & $120$ & $40/120$ solved ($60/120$ at $\ge 90\%$) & $47$--$62\%$ ($0/120$ solved) & $50\%$\\
\bottomrule
\end{tabular}}
\end{table}

\begin{table}[!htbp]
\caption{\textbf{Initialization scale, depth, and first-occurrence controls} (held-out exact-continuation accuracy \%).}
\label{tab:controls}
\centering\footnotesize
\fitwidth{\begin{tabular}{@{}llccc@{}}
\toprule
Setting & Runs & Step accuracy & Outcome accuracy & Chance\\
\midrule
Init scale $\initscale=0.1$, $n=8$ / $n=16$ & $3$ / $3$ & $100\%$ / $100\%$ & $63$--$65\%$ / $48$--$50\%$ & $50\%$\\
Init scale $\initscale=0.5$, $n=8$ / $n=16$ & $3$ / $3$ & $8$--$69\%$ / $0$--$1\%$ & $50$--$53\%$ / $50$--$52\%$ & $50\%$\\
Deeper models ($4$ or $6$ blocks), $n \in \{8,12\}$ & $12$ & $31$--$100\%$ ($10/12$ reach $\ge 80\%$) & $47.8$--$50.9\%$ ($0/12$) & $50\%$\\
First-occurrence control ($\ge 31\%$ deg-$1$ power), $n \in \{8,12,16,24\}$ & $12$ & $80$--$99\%$ ($12/12$ reach $\ge 90\%$ by $16$k) & $71$--$100\%$ ($11/12$ reach $\ge 90\%$ by $16$k) & $25$--$28\%$\\
\bottomrule
\end{tabular}}
\end{table}

\subsection{Automata and Pretrained Models}
\label{app:circuits}
\label{app:bbh}
\label{app:p-half}
\label{app:reach}

In Boolean-8, each prompt initializes $s_0 \in \{0,1\}^4$ ($16$ state tokens) and samples $n = 8$ invertible gates uniformly over gate types (\texttt{NOT}/\texttt{XOR-1}, \texttt{XNOR}, \texttt{SWAP}, \texttt{CNOT}, \texttt{TOFFOLI}) and distinct wire tuples (all even permutations on $\{0,1\}^4$, generating the alternating group $A_{16}$ with $\sigma_2(\bar P) \approx 0.84 < 1$). In the modulo-$17$ register machine, each prompt initializes two registers uniformly in $\mathbb Z_{17}$ and samples $n = 16$ instructions uniformly over (\texttt{INC}, \texttt{DEC}, \texttt{ADD}, \texttt{SUB}, \texttt{SWAP}, \texttt{CMOV}) and operands before querying a target register (including non-bijective \texttt{CMOV}, testing high-degree state tracking beyond \cref{cor:reversible}). LoRA runs serialize each state/instruction trace in plain-text \texttt{s0 | op1 -> s1 | ... -> sn} format (and \texttt{s0 | op1 ... opn -> sn} for outcome). Exact ANOVA decompositions place $98.9\%$ (Boolean-8) and $99.9\%$ (register mod $17$) of target variance at interaction degrees $|S| \ge 2$. Both GPT-2-2L trained from scratch and SmolLM2-135M / Qwen2.5-0.5B fine-tuned with rank-$8$ LoRA exhibit a sharp step-versus-outcome separation (\cref{tab:main}), which persists under both an $8\times$ outcome budget ($6{,}400$ updates versus $800$ step updates at $n=8$, matching total supervised tokens) and answer-first trace ordering.

\begin{table}[!htbp]
\caption{\textbf{Held-out accuracy (\%) across tasks and model tiers} (mean and min--max range over seeds). Counting rows report seeds solved ($\ge 99.9\%$ exact) at the stated update budget; at $32$k updates the $p=1/2$ outcome runs reach $48$--$99\%$ with $0/10$ solved. First-occurrence counts use the $\ge 90\%$ criterion of \cref{tab:controls}.}
\label{tab:main}
\centering\footnotesize
\fitwidth{\begin{tabular}{@{}llccccc@{}}
\toprule
Task & Model & $n$ & Seeds & Outcome & Step & Chance\\
\midrule
count mod $2$ ($p=1/4$, $8$k) & two-block & $16$ & $10$ & $49.8$ ($46.9$--$51.6$, $0/10$) & $93.1$ ($47.3$--$100$, $6/10$) & $50.0$\\
count mod $2$ ($p=1/2$, $16$k) & two-block & $16$ & $10$ & $50.7$ ($49.0$--$52.6$, $0/10$) & $95.2$ ($51.7$--$100$, $9/10$) & $50.0$\\
first occurrence & two-block & $8$--$24$ & $12$ & $71$--$100$ ($11/12$ $\ge 90\%$) & $80$--$99$ ($12/12$ $\ge 90\%$) & $25$--$28$\\
Boolean-8 & GPT-2-2L & $8$ & $5$ & $7.7$ ($6.3$--$8.9$) & $82.0$ ($77.6$--$88.2$) & $6.25$\\
Register mod $17$ & GPT-2-2L & $16$ & $3$ & $6.3$ ($6.2$--$6.4$) & $99.9$ ($99.8$--$100$) & $5.88$\\
Boolean-8 (LoRA) & SmolLM2-135M / Qwen2.5-0.5B & $8$ & $3$ & $9.9$ ($9.7$--$10.1$) / $11.1$ ($10.0$--$12.3$) & $89.7$ ($84.8$--$94.6$) / $98.3$ ($97.5$--$99.0$) & $6.25$\\
\bottomrule
\end{tabular}}
\end{table}

\subsection{ODE Verification}
\label{app:proof-checks}
\label{app:verify}
\label{app:J-check}
\label{app:J-generalK}
\label{tab:verify}

\begin{table}[!htbp]
\caption{\textbf{Step-ODE audit of the distractor floors (F).} Top ($n \le 10$): full \stepode{} RK4 ($\initscale = 0.08$, $d_k = 2n+6$, step capped at $2/\initscale^2$; window $[\tau_0, T_{\mathrm{solve}}]$), reporting extremes over $p \in \{0.1, 0.25, 0.5, 0.75\}$ and $r_{\mathrm{align}} \in \{0.05, 0.5\}$ ($\initscale = 0.05$ also at $n \le 8$; at $n = 10$, $p \in \{0.25, 0.5, 0.75\}$, $r_{\mathrm{align}} = 0.5$; $38$ runs). Bottom ($n \ge 20$): small-$\initscale$ limiting reciprocal-weight ODE on $[\tau_0, \tau_1]$ ($\mathrm{d}\tilde s = \frac{2p(1-p)ga}{n-1}\,\rmd s$, parameter-free in $(p, d_k, r_{\mathrm{align}}, \initscale)$; extending to $n = 2{,}000$ gives $\min n\,\omega_{\mathrm{vis}} = 2.22$, $\min\omega_{\mathrm{late}} = 0.9999$, and $u_t \in [-0.01, 7.18]$). In every run the visible minimum is attained at $\delta = 2$, staying $> 15\times$ above $\frac{1}{7n} \approx \frac{0.14}{n}$.}
\label{tab:U-floor-audit}
\centering\footnotesize
\begin{tabular}{@{}rccccc@{}}
\toprule
$n$ & $\min n\,\omega_{\mathrm{vis}}$ & $\min\omega_{\mathrm{late}}$ & $\sup\Omega_{\max}$ & $\min n\,\pi^\star$ & $u_t$ range\\
\midrule
4   & $4.00$ & ---      & $46.4$ & $1.00$ & $[-0.41, 0.74]$\\
6   & $5.40$ & $1.0000$ & $23.2$ & $1.00$ & $[-0.29, 1.14]$\\
8   & $4.77$ & $1.0000$ & $12.8$ & $1.00$ & $[-0.22, 1.44]$\\
10  & $4.33$ & $1.0000$ & $7.2$  & $1.00$ & $[-0.18, 1.68]$\\
\midrule
20  & $3.38$ & $1.0000$ & $3.9$  & $1.00$ & $[-0.10, 2.45]$\\
50  & $2.75$ & $1.0000$ & $2.4$  & $1.00$ & $[-0.04, 3.44]$\\
100 & $2.51$ & $0.9995$ & $2.1$  & $1.00$ & $[-0.02, 4.16]$\\
200 & $2.37$ & $0.9996$ & $1.9$  & $1.00$ & $[-0.01, 4.87]$\\
\bottomrule
\end{tabular}
\end{table}

\begin{itemize}[leftmargin=*]
\item \textbf{Single-head untied routing.} Relaxing the two-head antisymmetric pairing of \cref{def:gate} to single-head untied initialization ($d_k = 2n+6$, active ReLU MLPs, population expectation over $2^n$ binary match patterns) and training query--key blocks, relative lag bias $\brel$, and readout jointly by full-batch GD (step size $2$, $\initscale = 0.3$, $n \in \{6,8,12,16\}$), all $17/17$ step runs reach $\ge 99.9\%$ exact rollout accuracy ($16$ at $100\%$), whereas $0/17$ outcome runs solve within $300$k updates ($50$--$55\%$ except one late-escaping $n=6$ seed at $94\%$).
\item \textbf{Unmasked \stepode{} solve-time scaling (\cref{thm:joint,tab:J-ode}).} Integrating both unconstrained matrix gradient flow on $\theta_{\mathrm s}$ and \stepode{} via adaptive Radau~IIA ($\mathrm{rtol}=10^{-10}, \mathrm{atol}=10^{-13}$, matching to $<10^{-12}$ max-norm discrepancy) across all $60$ unmasked trajectories ($p \in \{0.10, 0.25, 0.50, 0.75\}$, three seeds spanning $\operatorname{sign}\langle q(0),k(0)\rangle = \pm 1$, $\initscale \in \{0.20, 0.14, 0.10, 0.07, 0.05\}$ at $n=4, d_k=8$), all $60/60$ runs reach worst-case margin $\min_{w\in\Sigma^n} M_t \ge 0.5$ with $\min_t\pi_t^\star \in [0.977, 0.997]$ and empirical $\initscale$-exponents in $[5.035, 5.048]$, consistent with $\Theta(\initscale^{-5})$. Across $n \in \{2,4,6,8,10\}$ at $\initscale = 0.10$, the empirical $n$-exponent is $1.020$ ($p=0.25$) and $1.021$ ($p=0.50$), consistent with $\Theta(n)$.

Even at fan-in scale $\initscale = 0.25\,n^{-1/2}$ ($d_k = 2n+6$, $r_{\mathrm{align}} = n^{-2}$) outside the small-$\initscale$ threshold of \cref{thm:block-rms}, all $18/18$ population trajectories of the block-RMS attention flow ($n \in \{2,4,6,8,10,12\}$, $p \in \{0.10, 0.25, 0.50\}$) reach $\min_t\pi_t^\star \ge 3/4$ and $\min_{w\in\Sigma^n} M_t \ge 0.5$ with horizon ratio $T_{\mathrm{RMS}}/n \in [9.69, 13.01]$ ($p=0.50$), $[11.28, 15.38]$ ($p=0.25$), and $[16.95, 24.38]$ ($p=0.10$).
\item \textbf{Score-degree and rule-credit exponents (\cref{prop:represent,prop:score-degree}).} The measured $\initscale$-exponents of $\sum_{|S|=r}\|(\Jz)_S\|$ at degrees $r = 1, \ldots, 5$ equal $2.0, 6.1, 7.0, 11.1, 12.0$ (consistent with the lower bounds $2, 6, 6, 10, 10$ of \cref{prop:score-degree}); step rule credit matches $\initscale^6 P(\tilde\theta)$ to $6$ digits across $n \in [2,12]$ with \mbox{$n\|\nabla_\theta\rulecredit\|_2/\initscale^6 \in [130, 324]$} and $\initscale$-exponent $6.00$; and $\thstep, \thout$ achieve $100\%$ exact-trace accuracy at $n = 1{,}000$ ($\dmodel = 2{,}025$).
\end{itemize}

\begin{table}[!htbp]
\caption{\textbf{Numerical check of the unmasked step system \stepode{} and $\initscale^{-5}$ solve-time scaling} ($n=4$, $d_k=8$, worst-case margin $\min_{w\in\Sigma^n} M_t \ge 0.5$, $60/60$ runs solved with $\min_t\pi_t^\star \ge 0.977$).}
\label{tab:J-ode}
\centering\footnotesize
\fitwidth{\begin{tabular}{@{}lcccccc@{}}
\toprule
Match rate $p$ & $\operatorname{sign}\langle q_0, k_0\rangle$ & $\initscale=0.20$ & $\initscale=0.10$ & $\initscale=0.05$ & $\min_t\pi_t^\star$ & Slope in $\log(1/\initscale)$\\
\midrule
$p = 0.10$ & $+1$ / $-1$ & $(1.34\text{--}1.77)\times 10^{6}$ & $(4.47\text{--}5.85)\times 10^{7}$ & $(1.46\text{--}1.90)\times 10^{9}$ & $0.997$ & $5.035\text{--}5.045$\\
$p = 0.25$ & $+1$ / $-1$ & $(6.42\text{--}8.48)\times 10^{5}$ & $(2.14\text{--}2.80)\times 10^{7}$ & $(7.01\text{--}9.13)\times 10^{8}$ & $0.994$ & $5.036\text{--}5.046$\\
$p = 0.50$ & $+1$ / $-1$ & $(4.80\text{--}6.35)\times 10^{5}$ & $(1.61\text{--}2.10)\times 10^{7}$ & $(5.25\text{--}6.84)\times 10^{8}$ & $0.977$ & $5.037\text{--}5.048$\\
$p = 0.75$ & $+1$ / $-1$ & $(6.42\text{--}8.48)\times 10^{5}$ & $(2.14\text{--}2.80)\times 10^{7}$ & $(7.01\text{--}9.13)\times 10^{8}$ & $0.994$ & $5.036\text{--}5.046$\\
\bottomrule
\end{tabular}}
\end{table}
